% Taxer mieux plutôt que davantage : quatre réformes à recette constante (rapport Sorbon, partie III)
% Note unique issue de la fusion des notes ir_csg, foncier, is et pfu (novembre 2026).
% Compilation depuis fiscalite/synthese/ : pdflatex note-reforme-fiscale.tex (deux passes)
% Chiffres : sorties/perte_seche_ircsg_purs.txt (IR-CSG), ../foncier/sorties/simulation.txt,
% ../is/sorties/simulation_is.txt, sorties/perte_seche_autres.txt ; méthode : ../revue/methode_perte_seche.md
\documentclass[11pt,a4paper]{article}
% Préambule commun au brief et à la revue de littérature sur la taxe sur la valeur foncière.
% Même composition que le cours (cours/refaire-l-impot.tex) ; figures « nues » produites par
%   FIG_NU=1 python lvt/graphiques_lvt.py
\usepackage[T1]{fontenc}
\usepackage[utf8]{inputenc}
\usepackage{lmodern}
\usepackage[french]{babel}
\usepackage[a4paper,left=2.8cm,right=2.8cm,top=2.6cm,bottom=2.8cm]{geometry}
\usepackage{graphicx}
\usepackage[export]{adjustbox}
\usepackage{booktabs}
\usepackage{tabularx}
\usepackage{array}
\usepackage{xcolor}
\usepackage{titlesec}
\usepackage[labelfont=bf,labelsep=period,font=small,singlelinecheck=false,
            justification=raggedright,skip=6pt]{caption}
\usepackage{fancyhdr}
\usepackage{enumitem}
\usepackage{microtype}
\usepackage{textcomp}
\usepackage{amsmath}
\providecommand{\euro}{\texteuro}
\usepackage[hidelinks]{hyperref}

\definecolor{encre}{HTML}{1A1A19}
\definecolor{bleu}{HTML}{003399}
\definecolor{vert}{HTML}{1F7A4D}
\definecolor{ocre}{HTML}{C79100}
\definecolor{brique}{HTML}{993333}
\definecolor{filet}{HTML}{C8C6BE}

\titleformat{\section}{\normalfont\large\bfseries\color{bleu}}{\thesection.}{0.6em}{}
\titleformat{\subsection}{\normalfont\normalsize\bfseries\color{bleu}}{\thesubsection}{0.6em}{}
\titlespacing*{\section}{0pt}{2.4ex plus 1ex minus .2ex}{1.2ex plus .2ex}
\titlespacing*{\subsection}{0pt}{2ex plus 1ex minus .2ex}{0.8ex plus .2ex}

\pagestyle{fancy}
\fancyhf{}
\renewcommand{\headrulewidth}{0.4pt}
\fancyhead[R]{\small\thepage}

\setlength{\parindent}{0pt}
\setlength{\parskip}{0.62em}
\linespread{1.05}
\renewcommand{\arraystretch}{1.2}
\setlist{itemsep=2pt,topsep=2pt,parsep=0pt,leftmargin=1.4em}

\newenvironment{encadre}[1]
  {\par\medskip\noindent\begin{minipage}{\linewidth}%
   \color{encre}\setlength{\parskip}{0.5em}%
   \hrule height 0.8pt width \linewidth \vspace{6pt}%
   {\footnotesize\bfseries\color{vert}\MakeUppercase{#1}}\par\vspace{-2pt}%
   \small}
  {\par\vspace{4pt}\hrule height 0.4pt width \linewidth\end{minipage}\par\medskip}

\newcommand{\figbloc}[3]{%
  \begin{figure}[htbp]\centering
  \caption{#2}\label{#3}
  \includegraphics[max width=\linewidth,max height=0.40\textheight]{../figures/nu/#1.pdf}
  \IfFileExists{../figures/nu/#1_note.tex}{\par\vspace{4pt}%
    \parbox{\linewidth}{\raggedright\scriptsize\color{encre}\setlength{\parskip}{2pt}%
    \input{../figures/nu/#1_note.tex}}}{}%
  \end{figure}}

\newcommand{\pc}{\,\%}

\renewcommand{\figbloc}[3]{%
  \begin{figure}[!htbp]\centering
  \caption{#2}\label{#3}
  \includegraphics[max width=\linewidth,max height=0.28\textheight]{../figures/nu/#1.pdf}
  \par\vspace{3pt}\parbox{\linewidth}{\raggedright\scriptsize\color{encre}\setlength{\parskip}{2pt}%
    \input{../figures/nu/#1_note.tex}}%
  \end{figure}}
\newcommand{\figgrand}[3]{\figbloc{#1}{#2}{#3}}
\usepackage{needspace}
\usepackage{etoolbox}
\preto\section{\needspace{8\baselineskip}}
\preto\subsection{\needspace{5\baselineskip}}
\renewcommand{\topfraction}{0.9}\renewcommand{\bottomfraction}{0.8}
\renewcommand{\textfraction}{0.08}\renewcommand{\floatpagefraction}{0.85}
\setcounter{topnumber}{3}\setcounter{totalnumber}{5}
\fancyhead[L]{\small\itshape Taxer mieux plutôt que davantage}

\begin{document}
\renewcommand{\figurename}{Figure}\renewcommand{\tablename}{Tableau}

\thispagestyle{empty}
\vspace*{0.6cm}
{\footnotesize\color{ocre}\bfseries NOTE DE TRAVAIL\quad·\quad NOVEMBRE 2026\par}
\vspace{8pt}
{\fontsize{22}{26}\selectfont\bfseries\color{bleu}Taxer mieux plutôt que davantage\par}
\vspace{6pt}
{\large Quatre réformes à recette constante : impôts sur le revenu, immobilier, sociétés, revenus du capital\par}
\vspace{8pt}
{\small Colin Baget. Droit de 2026 ; données fiscales de la DGFiP (revenus 2024, fiscalité locale 2024, IS
2022), ventes immobilières DVF 2023-2024, comptes nationaux de l'Insee ; programmes \texttt{fiscalite/synthese/}
et dossiers \texttt{ir\_csg/}, \texttt{foncier/}, \texttt{is/}, \texttt{pfu/}. Vérifications et revues :
\texttt{fiscalite/revue/}.\par}
\vspace{6pt}
{\color{bleu}\hrule height 1.2pt}
\vspace{4pt}

\section*{Introduction}

Cette note examine quatre réformes avec un seul critère : à recette à peu près constante, quel système
fiscal gêne le moins la croissance, c'est-à-dire décourage le moins de travailler, d'investir, de déménager
ou de déclarer ses revenus ? La redistribution n'est pas l'objectif, mais deux contraintes s'imposent. La
somme des réformes ne doit pas alourdir nettement la facture des plus grandes fortunes : leur départ coûterait
plus qu'il ne rapporterait. Et le système doit être, à la marge, plus favorable aux familles avec enfants.
Il s'agit de taxer mieux, pas davantage.

La méthode est la même pour les quatre impôts. On part d'une référence simple : l'assiette la plus large,
le même taux pour tout le monde, sans niche. Puis on fait apparaître chacun des endroits où un choix doit
être fait (une franchise pour les revenus d'activité, un traitement des familles, la vitesse de la
déduction de l'investissement, le sort des plus-values logées dans une holding) et on chiffre chaque
choix de la même façon : ce qu'il coûte en recette, ce qu'il coûte ou rapporte en perte sèche, et qui
paie. La note ne tranche pas ces arbitrages à la place du lecteur ; elle dit ce que chacun coûte.

\begin{encadre}{Comment la perte sèche est calculée}
\textbf{Ce qu'elle mesure.} Un impôt fait renoncer à des choses qui auraient valu plus qu'elles ne coûtaient :
une heure de travail, une machine, une vente de logement. La perte sèche est la valeur de ces renoncements,
au-delà de l'impôt lui-même. Elle se mesure par la réaction des quantités taxées à l'impôt.

\textbf{La formule.} La perte sèche ordinaire, ½·ε·τ²/(1$-$τ)·base, suppose une seule marge et aucun autre impôt
dessus. On utilise sa généralisation (Harberger ; Hines, 1999 ; Chetty, 2009) : la variation de perte sèche
est la somme, sur chaque marge touchée, du coin fiscal total qui frappe cette marge multiplié par la
variation de la quantité taxée. Un euro de salaire perdu coûte à la collectivité tout ce qu'il aurait payé :
impôt sur le revenu et CSG, mais aussi cotisations non contributives et prestations perdues. Trois
conséquences comptent :
\begin{itemize}\setlength{\itemsep}{0pt}
\item \emph{le coin total} : relever l'impôt là où les cotisations et le retrait des prestations prélèvent
déjà 60~\% d'un euro de salaire en plus (autour de 1 à 1,5 Smic) coûte beaucoup plus que le calcul sur le
seul impôt ne le dit ;
\item \emph{les déplacements entre impôts} : un revenu qui quitte le barème pour l'impôt sur les sociétés ou le
prélèvement forfaitaire n'est perdu que pour l'écart de taux ; une réaction de calendrier (distribuer un
dividende un an plus tard) ne coûte presque rien ;
\item \emph{la décision de travailler} : on compte aussi la réaction de ceux qui hésitent à prendre un emploi,
qui dépend de ce que l'impôt prélève sur l'ensemble du premier salaire.
\end{itemize}

\textbf{Comparer des variantes.} Quand deux variantes ne lèvent pas la même recette, la différence doit être
levée ailleurs. On la valorise au « prix d'un milliard » levé par le barème fusionné : 0,16 à 0,60~Md€ de perte
sèche par milliard selon le coin retenu, 0,38 au centre (méthode et paramètres : \texttt{revue/methode\_perte\_seche.md}).

\textbf{Ce qui reste fragile.} Les élasticités, et surtout deux paramètres sans estimation française : la part
de la réaction des hauts revenus qui n'est qu'un déplacement vers l'impôt sur les sociétés, et la part de
l'ajustement des distributions de dividendes qui n'est qu'un report. Chaque résultat est donné avec sa
fourchette.
\end{encadre}


\section{Fusionner l'impôt sur le revenu, la CSG et la CRDS : même revenu, même impôt}\label{sec:ircsg}
% Partie IR-CSG (méthode v2 : revue/methode_v2.md ; chiffres : sorties/architectures_ircsg.txt,
% sorties/controle_hauts_purs.txt ; tableaux générés par tables_ircsg.py)

L'impôt sur le revenu, la CSG et la CRDS prélèvent 220~Md€ sur les revenus soumis au barème : salaires,
pensions, bénéfices des indépendants, loyers, allocations chômage. Cette somme se décompose en 80~Md€ d'IR,
132~Md€ de CSG et 7~Md€ de CRDS. L'IR des revenus du capital au prélèvement forfaitaire n'y est pas compté :
ces revenus gardent leur régime (section~\ref{sec:pfu}). Les trois impôts se superposent sans se coordonner,
et la charge varie selon le profil :
\begin{itemize}\setlength{\itemsep}{0pt}
\item Pour une personne seule salariée, ils prélèvent 9,5~\% du salaire brut jusqu'à 23\,000~€ environ. Le
taux marginal saute alors à 21~\%, puis à 32~\% vers 40\,000~€ et à 40~\% au-delà de 110\,000~€
(figure~\ref{fig:baremes}).
\item Les retraités paient moins à revenu égal, du fait de l'abattement de 10~\% (5,3~Md€) et des taux
réduits de CSG. Ces taux réduits créent cependant des marches : quand le revenu fiscal franchit un seuil, le
nouveau taux s'applique à toute la pension, et 1\,000~€ de pension en plus peuvent coûter 620 à 770~€.
\item Le quotient familial des enfants réduit l'IR d'environ 13~Md€.
\end{itemize}

\subsection{Ce que dit la théorie}

Un taux marginal coûte par ce qu'il fait renoncer à gagner. Ce coût est d'autant plus élevé que l'euro en
question est déjà lourdement prélevé, et que les revenus concernés sont nombreux et capables de réagir.
\begin{itemize}\setlength{\itemsep}{0pt}
\item Une assiette large permet des taux plus bas là où ils coûtent. Les pensions déjà liquidées et les revenus
aujourd'hui exonérés par des niches (au moins 24~Md€) ne réagissent presque pas ; les taxer comme les salaires
permet, à recette égale, d'abaisser les taux qui pèsent sur le travail. C'est le mécanisme principal de la
fusion, et il fait des retraités les premiers contributeurs de la réforme.
\item L'élasticité des hauts revenus dépend des règles. Leur revenu déclaré réagit fortement au taux, mais
surtout par la forme, la date ou le lieu de la déclaration ; une assiette sans niche réduit ces possibilités.
Ce gain va à la recette, non à l'activité.
\item Les marches et les hauts taux du milieu coûtent plus que leur rendement ne le laisse croire. Entre 1 et
1,5~Smic, les cotisations non contributives et la sortie de la prime d'activité prélèvent déjà près de 60~\%
d'un euro de salaire en plus ; y ajouter un saut d'IR coûte beaucoup plus que ne le dit le seul taux de
l'impôt.
\item La décision de travailler dépend de l'impôt sur tout le premier salaire. Peu de salariés ajustent leurs
heures au taux marginal ; beaucoup plus choisissent de prendre un emploi, de reprendre après une interruption
ou d'ajouter un second revenu au foyer. Une franchise sur les premiers euros d'activité baisse l'impôt sur ce
premier salaire sans toucher au taux marginal des autres : c'est la logique du crédit d'impôt sur les revenus
d'activité introduit en Suède en 2007, dont les effets sur l'emploi restent d'ailleurs discutés.
\end{itemize}

\subsection{Ce qu'on peut chiffrer}

On mesure chaque architecture avec la méthode de l'encadré. Toutes partagent la même base :
\begin{itemize}\setlength{\itemsep}{0pt}
\item l'assiette est le revenu après cotisations, le même pour un salarié, un indépendant, un retraité ou un
chômeur, sans réduction, abattement ni taux réduit ;
\item l'impôt est individuel, et le même barème vaut pour tous ;
\item le taux marginal monte en ligne droite entre quelques seuils, jusqu'à 50~\% ;
\item la recette est toujours de 220~Md€.
\end{itemize}

\textbf{Le barème neutre.} Parmi ces barèmes, on retient d'abord celui qui s'écarte le moins de ce que chaque
foyer paie aujourd'hui : 2~\% au premier euro, 14~\% à 20\,000~€, 23~\% à 30\,000~€ et 50~\% à 250\,000~€ par
adulte (figure~\ref{fig:baremes}).
\begin{itemize}\setlength{\itemsep}{0pt}
\item \emph{Qui paie plus ou moins.} Aucun décile ne voit sa charge moyenne bouger de plus de 600~€
(tableau~\ref{tab:ircsg_dec}). Il reste qu'un foyer sur quatre paie plus de 500~€ de plus, à cause des niches
supprimées, de l'imposition individuelle et de l'alignement des retraités, qui paient 7~Md€ de plus.
\item \emph{Gain d'efficacité.} Le gain est de 4,0~Md€ par an, entre 1,9 et 9,2~Md€ selon les hypothèses.
\item \emph{D'où il vient.} Pour un peu moins de la moitié (1,8~Md€), de la participation : le taux d'entrée de
2~\% remplace les 9,7~\% de CSG et de CRDS prélevés dès le premier euro. Pour le reste, du dixième le plus
aisé : son taux marginal baisse, alors que son impôt moyen monte un peu avec la fin des niches.
\item \emph{Effet sur la recette.} La moindre optimisation des hauts revenus rapporterait en plus 1,5 à
4,5~Md€. Ce supplément est trop incertain pour financer quoi que ce soit, et il n'entre pas dans le gain.
\end{itemize}

\textbf{Ce qu'on peut obtenir au mieux.} Le barème simple qui minimise la perte sèche, franchise d'activité
comprise, montre ce que dirait la seule efficacité. Il sert ici d'étalon, non de proposition.
\begin{itemize}\setlength{\itemsep}{0pt}
\item Il gagne 7,8~Md€ (3,8 à 18,2).
\item Il est presque proportionnel : 16~\% jusqu'à 60\,000~€, et le taux de 50~\% n'est atteint qu'à 2~M€.
\item Le dernier décile y gagne 25~Md€, et la moitié des foyers perdent plus de 500~€ par an.
\end{itemize}

\figbloc{ircsg_p02_frontiere}{Gain d'efficacité et nombre de perdants : rendements décroissants}{fig:frontiere}

\textbf{Les rendements sont décroissants} (figure~\ref{fig:frontiere}).
\begin{itemize}\setlength{\itemsep}{0pt}
\item La fusion à répartition inchangée capte la moitié du gain atteignable.
\item Une franchise sur les revenus d'activité porte cette part à 59~\% avec 4\,000~€ par adulte, à 64~\% avec
8\,000~€. C'est le seul dispositif qui ajoute du gain, par la marge extensive.
\item Aller plus loin demande d'aplatir le barème au-dessus de 45\,000~€. Les 30 premiers pour cent du chemin
vers l'optimum portent la part de 52 à 69~\%, mais le dernier décile passe de +2,5 à $-$6~Md€. Le dernier tiers
du gain exige de lui rendre 20 à 25~Md€, et fait perdre plus de 500~€ à la moitié des foyers.
\end{itemize}
On obtient donc environ deux tiers du gain atteignable sans bouleverser la répartition par décile.

\begin{table}[!htbp]\centering\scriptsize\setlength{\tabcolsep}{3pt}
\caption{Architectures comparées, à recette égale (220~Md€)}\label{tab:ircsg_arch}
\resizebox{\linewidth}{!}{%
\begin{tabular}{@{}lrrrrrrr@{}}
\toprule
& Points & Gain d'efficacité & Part du & Perdants & Retraités & Couple, & 10~\% \\
& pour tous & Md€ par an & gain & $>$ 500~€ & Md€ & 2 enfants & du haut \\
& & (fourchette) & atteignable & & & € par foyer & Md€ \\
\midrule
Barème neutre & 0,0 & 4,0 (1,9 à 9,2) & 52~\% & 25~\% & +7,0 & +570 & +2,5 \\
\multicolumn{8}{@{}l}{\emph{Un dispositif ajouté au barème neutre}} \\
\quad Franchise d'activité 4~000~€ & +1,8 & 4,6 (2,1 à 10,2) & 59~\% & 30~\% & +14,0 & +217 & +3,6 \\
\quad Franchise d'activité 8~000~€ & +3,6 & 5,0 (2,2 à 11,0) & 64~\% & 33~\% & +21,4 & $-$172 & +4,7 \\
\quad Franchise de 8~000~€ sur tous les revenus & +5,8 & 3,3 (1,4 à 7,4) & 43~\% & 21~\% & +5,4 & +775 & +9,6 \\
\quad Franchise de 6~000~€ par enfant & +1,5 & 3,6 (1,7 à 8,4) & 47~\% & 27~\% & +13,0 & $-$983 & +0,6 \\
\quad Chèque de 742~€ par enfant & +0,9 & 3,1 (1,5 à 7,3) & 40~\% & 26~\% & +10,7 & $-$361 & +3,1 \\
\quad Quotient conjugal et familial & +2,0 & 3,9 (1,9 à 9,0) & 50~\% & 28~\% & +12,1 & $-$965 & $-$3,0 \\
\quad Abattement de 10~\% sur les pensions & +0,6 & 3,4 (1,6 à 8,0) & 44~\% & 24~\% & +0,9 & +923 & +2,7 \\
\quad Petites pensions exonérées & +0,4 & 3,7 (1,7 à 8,5) & 47~\% & 16~\% & +3,2 & +786 & +4,0 \\
\midrule
\textbf{Proposition}, franchise d'activité 4~000~€ & --- & 4,9 (2,3 à 11,1) & 63~\% & 29~\% & +15,2 & $-$1\,195 & $-$1,7 \\
Proposition, franchise d'activité 8~000~€ & --- & 5,2 (2,4 à 11,8) & 67~\% & 26~\% & +22,8 & $-$1\,629 & $-$1,5 \\
\midrule
Optimum d'efficacité, sans franchise & --- & 7,4 (3,7 à 17,4) & 95~\% & 58~\% & +13,7 & $-$519 & $-$25,7 \\
Gain atteignable (optimum, franchise 4~000~€) & --- & 7,8 (3,8 à 18,2) & 100~\% & 49~\% & +20,3 & $-$869 & $-$25,1 \\
\bottomrule
\end{tabular}}
\par\vspace{3pt}\parbox{\linewidth}{\raggedright\scriptsize\color{encre}
Gain d'efficacité : baisse de la perte sèche sur l'activité réelle par rapport à aujourd'hui, hypothèses centrales (basses et hautes entre parenthèses) ; la réaction du revenu déclaré n'y est pas comptée. Points pour tous : hausse des taux du barème neutre qui finance le dispositif. Proposition : franchise d'activité, 6~000~€ par enfant, pensions exonérées jusqu'à 12~500~€ puis de moins en moins jusqu'à 25~000~€ ; barème le plus efficace parmi ceux qui ne changent la charge moyenne d'aucun décile de plus de 500~€ et n'alourdissent pas celle du 1~\% le plus aisé. Retraités, couples, 10~\% du haut : variation de leur impôt (positif : ils paient plus).}
\end{table}


\textbf{Chaque dispositif a un prix} (tableau~\ref{tab:ircsg_arch}). Les dispositifs qui ne touchent pas à
l'activité, aide aux enfants ou ménagement des pensions, se financent par des taux un peu plus élevés pour
tous. Chacun réduit le gain de 0,1 à 0,9~Md€. Quatre arbitrages ressortent.
\begin{itemize}\setlength{\itemsep}{0pt}
\item \emph{Activité ou tous les revenus.} Une même franchise de 8\,000~€ étendue à tous les revenus coûte
deux points de plus et fait perdre un tiers du gain. Elle donne aux retraités ce que la franchise d'activité
réserve aux actifs, sans aucun effet sur l'emploi.
\item \emph{Ce que couvre la franchise d'activité.} Le modèle ne distingue pas, dans les revenus d'activité,
les salaires des loyers et des bénéfices. Il ne peut donc pas chiffrer une franchise réservée aux salaires.
Une franchise sur les loyers ne rend pourtant aucun travail plus payant.
\item \emph{Pensions.} Garder l'abattement de 10~\% ramène la hausse des retraités de 7 à 1~Md€, pour
0,6~point de taux pour tous. Exonérer les petites pensions la ramène à 3~Md€, pour 0,4~point.
\item \emph{Enfants.} Une franchise de 6\,000~€ par enfant est l'option la plus favorable aux familles :
1\,000~€ de moins en moyenne pour un couple avec deux enfants. Le chèque de 742~€ par enfant, qui reverse à
parts égales le coût actuel du quotient, l'est moins (360~€). Le quotient conjugal et familial profite surtout
aux couples aisés : le dixième le plus aisé y gagne 3~Md€.
\end{itemize}

\subsection{Ce que nous proposons}

\figbloc{ircsg_p01_baremes}{Taux moyen et taux marginal : aujourd'hui, barème neutre et proposition}{fig:baremes}

\textbf{La proposition} retient les trois dispositifs que la note défend, et fixe la répartition avant de
choisir le barème.
\begin{itemize}\setlength{\itemsep}{0pt}
\item \emph{Les dispositifs.} Une franchise de 4\,000~€ par adulte sur les revenus d'activité, pour l'emploi.
Une franchise de 6\,000~€ par enfant, déduite du revenu du parent qui gagne le plus, pour les familles.
L'exonération des pensions jusqu'à 12\,500~€, à peu près le minimum vieillesse, puis de moins en moins
jusqu'à 25\,000~€, pour que le minimum vieillesse ne soit pas imposé.
\item \emph{La contrainte de répartition.} La charge moyenne d'aucun décile ne doit bouger de plus de 500~€, et
celle du 1~\% le plus aisé ne doit pas monter.
\item \emph{Le barème.} Parmi ceux qui respectent la contrainte, on prend le plus efficace : 14~\% jusqu'à
20\,000~€ imposés, puis une montée en ligne droite jusqu'à 30~\% à 45\,000~€, puis une montée lente jusqu'à
50~\% à 1~M€.
\end{itemize}
Le gain d'efficacité est de 4,9~Md€ par an (2,3 à 11,1), soit 63~\% du gain atteignable ; 29~\% des foyers
perdent plus de 500~€.

\textbf{Qui paie.} Les actifs paient 15~Md€ de moins, les retraités 15~Md€ de plus. Le tableau~\ref{tab:ircsg_dec}
donne la variation par décile.
\begin{itemize}\setlength{\itemsep}{0pt}
\item \emph{Familles.} Un couple avec deux enfants paie 1\,200~€ de moins en moyenne (1\,360~€ à 50\,000~€ de
revenu), un parent isolé 230~€ de moins.
\item \emph{Salariés.} Un salarié au Smic paie à peu près ce qu'il paie aujourd'hui (+60~€) ; une personne
seule active, 320~€ de moins en moyenne.
\item \emph{Retraités.} Un retraité seul paie 600~€ de plus en moyenne, un couple de retraités 2\,300~€. Un
retraité seul à 17\,000~€ de pension, aujourd'hui exonéré, paierait 1\,260~€.
\end{itemize}
Entre 12\,500 et 25\,000~€ de pension, la sortie de l'exonération fait monter le taux marginal à 28~\%. C'est
sans effet sur l'activité, mais visible ; une sortie plus lente l'adoucirait, au prix d'un peu plus de taux pour
tous.

Avec une franchise de 8\,000~€, le gain monte à 5,2~Md€ (67~\%). Les retraités paient alors 23~Md€ de plus, et un
couple de retraités 3\,200~€. C'est l'arbitrage central de la réforme : l'emploi et les familles se financent
pour l'essentiel sur les pensions.

\begin{table}[!htbp]\centering\scriptsize\setlength{\tabcolsep}{3pt}
\caption{Variation de l'impôt par décile de revenu fiscal, toutes variantes (€ par foyer et par an)}\label{tab:ircsg_dec}
\resizebox{\linewidth}{!}{%
\begin{tabular}{@{}lrrrrrrrrrrrr@{}}
\toprule
& D1 & D2 & D3 & D4 & D5 & D6 & D7 & D8 & D9 & D10 & 1~\% (Md€) & Retraités (Md€) \\
\midrule
Réforme naïve (quotient, 3,99~\%) & $-$11 & $-$37 & +162 & +105 & $-$543 & $-$551 & $-$570 & +95 & $-$524 & +1\,841 & +2,6 & +8,2 \\
Réforme naïve, variante individuelle (1,64~\%) & $-$57 & $-$217 & $-$171 & $-$256 & $-$906 & $-$938 & $-$794 & +36 & +277 & +2\,953 & +2,8 & +2,7 \\
\textbf{Proposition} & $-$84 & $-$279 & +80 & +251 & $-$215 & $-$155 & $-$320 & +354 & $-$71 & +436 & 0,0 & +21,5 \\
Variante : franchise de 4\,000~€ & $-$85 & $-$149 & +147 & +225 & $-$286 & $-$129 & $-$237 & +304 & $-$158 & +367 & 0,0 & +14,3 \\
Variante : sans franchise & +2 & +19 & +243 & +318 & $-$225 & $-$202 & $-$414 & +113 & $-$259 & +399 & 0,0 & +6,9 \\
Quotient conjugal et familial & $-$84 & $-$285 & $-$18 & +119 & $-$276 & $-$137 & $-$149 & +468 & $-$99 & +460 & 0,0 & +18,9 \\
Abattement de 10~\% sur les pensions & $-$83 & $-$277 & +42 & +273 & $-$200 & $-$138 & $-$287 & +314 & $-$94 & +448 & 0,0 & +14,8 \\
Chèque par enfant (même coût) & $-$168 & $-$354 & +16 & +197 & $-$226 & $-$99 & $-$227 & +392 & +11 & +455 & 0,0 & +21,6 \\
Le 1~\% peut payer moins & $-$84 & $-$278 & +105 & +335 & $-$86 & $-$11 & $-$197 & +493 & +162 & $-$423 & $-$3,8 & +22,5 \\
Barème neutre (méthode v2) & 0 & +83 & +265 & +165 & $-$365 & $-$221 & $-$324 & +130 & $-$27 & +290 & $-$0,5 & +5,0 \\
Ancienne proposition (1~\% en baisse) & $-$85 & $-$145 & +190 & +342 & $-$150 & $-$12 & $-$121 & +450 & 0 & $-$453 & $-$2,5 & +13,9 \\
\bottomrule
\end{tabular}}
\par\vspace{3pt}\parbox{\linewidth}{\raggedright\scriptsize\color{encre}
Déciles de foyers fiscaux classés par revenu fiscal de référence (IRCOM, revenus 2024 ; montants de 2025). Moyennes de foyers représentatifs : elles cachent les écarts à l'intérieur de chaque décile (sous-partie~\ref{sec:ircsg}.4).}
\end{table}


\textbf{Les très hauts revenus.} Pour les 3\,800 foyers qui déclarent plus de 3~M€, le résultat dépend d'une
donnée que la DGFiP ne publie pas : la part de leurs revenus imposée au prélèvement forfaitaire.
\begin{itemize}\setlength{\itemsep}{0pt}
\item \emph{Si leurs revenus du barème paient aujourd'hui ce que dit le barème légal}, 44 à 49~\% : la
proposition baisse leur impôt de 6 à 13~points de ce revenu. Seuls les 700 foyers au-delà de 9~M€ paient
1~point de plus.
\item \emph{Dans la convention centrale du modèle}, qui leur prête un taux actuel de 32 à 37~\%, difficile à
concilier avec le barème : ils paieraient 2 à 16~points de plus.
\end{itemize}
Dans les deux cas, la comparaison omet la CEHR, qui disparaîtrait sur les revenus du barème et leur retirerait
3 à 4~points. Leurs revenus du capital restent au prélèvement forfaitaire, majoration comprise
(section~\ref{sec:pfu}). Caler cette part sur les données de la DGFiP par tranche est le premier travail à faire
avant de chiffrer définitivement.

\textbf{La transition.}
\begin{itemize}\setlength{\itemsep}{0pt}
\item La collecte ne change presque pas : employeurs, caisses de retraite et France Travail prélèvent déjà à
la source l'IR et la CSG sur les mêmes revenus. Le taux individualisé est la règle par défaut du prélèvement à
la source depuis le 1\textsuperscript{er}~septembre~2025 (loi de finances pour 2024, art.~19).
\item Une convergence linéaire sur quatre ans limiterait la hausse annuelle du retraité à 17\,000~€ à environ
320~€. Comme gagnants et perdants avancent au même rythme, la recette reste à peu près constante chaque année.
\end{itemize}

\textbf{Le risque constitutionnel.} Le Conseil constitutionnel a censuré :
\begin{itemize}\setlength{\itemsep}{0pt}
\item en 2000, une réduction de CSG qui ignorait les autres revenus du foyer et les charges de famille
(décision 2000-437~DC) ;
\item en 2012, une contribution assise sur chaque personne sans tenir compte du foyer (2012-662~DC).
\end{itemize}
Un impôt individuel avec une franchise réservée aux revenus d'activité réunit ces deux traits. La proposition
s'en distingue parce qu'elle remplace l'impôt de foyer au lieu de s'y ajouter et qu'elle tient compte des
enfants. L'argument reste à construire. La variante « quotient », qui rétablit l'imposition commune des
couples, est la solution de repli : elle coûte peu en efficacité (3,9~Md€ de gain au lieu de 4,0 pour le
barème neutre).


\section{Immobilier : taxer le sol plutôt que les ventes}\label{sec:foncier}
% Partie foncier, extraite de ../foncier/note-foncier.tex puis retouchée
La France taxe beaucoup l'immobilier, et elle le taxe au mauvais endroit. Elle prélève lourdement au
moment où un logement change de mains : un ménage qui achète un logement ancien paie 5,8 à 6,3~\% du prix en
droits de mutation, en plus du prix, à l'instant précis où son apport est le plus contraint. Elle prélève
aussi beaucoup sur la détention, mais mal : la taxe foncière repose sur des valeurs locatives de 1970, qui
font payer proportionnellement plus les logements des communes pauvres que ceux des communes chères.

L'économie de ce choix est défavorable. Un impôt sur les ventes décourage les ventes : il retient dans leur
logement des ménages qui voudraient en changer, freine les déménagements vers un logement plus grand quand
la famille s'agrandit, et pèse d'abord sur ceux qui achètent, souvent jeunes. Un impôt récurrent sur la
valeur du sol, à l'inverse, ne décourage ni le travail ni la construction : le terrain existe de toute
façon. Il a un défaut, qui explique qu'on s'en serve si peu : il se voit, chaque année, sur un avis
d'imposition.

Cette note dresse l'état des lieux de tous les prélèvements assis sur l'immobilier, résume ce que la
recherche établit sur les droits de mutation et sur la taxe foncière, puis simule une bascule : diviser
par deux les droits de mutation, et reporter la recette sur une taxe foncière refondée, assise sur la
valeur de marché, avec un taux plus élevé sur le terrain que sur le bâti. L'objectif est une fiscalité
favorable à la croissance et à l'accession à la propriété, en particulier des jeunes ménages, sans
déstabiliser brutalement les budgets des propriétaires en place.

\subsection{État des lieux : 116 milliards, dont un quart sur les ventes}

Les prélèvements assis sur l'immobilier atteignent environ 116~Md€ en 2024, 4~\% du PIB
(figure~\ref{fig:inv}). La détention en fournit 69~Md€ : la taxe foncière sur les propriétés bâties
(42~Md€, 43,1 en 2025, dont environ 31,6~Md€ sur les logements), la taxe d'enlèvement des ordures ménagères
qui lui est adossée (9~Md€), la cotisation foncière des entreprises (8~Md€), l'impôt sur la fortune
immobilière (2,7~Md€) et les taxes d'habitation qui subsistent sur les résidences secondaires et les
logements vacants. Les transactions en fournissent 25~Md€ : droits de mutation (14,1~Md€ en 2024,
16,7 en 2025), TVA sur les logements neufs et les terrains à bâtir (7,1~Md€), impôt sur les plus-values
immobilières (2,8~Md€), contribution de sécurité immobilière (0,7~Md€). S'y ajoutent les droits de
succession et de donation sur l'immobilier (environ 12,5~Md€, estimation faute de ventilation
officielle), l'impôt sur les revenus locatifs (8,4~Md€) et la taxe d'aménagement (1,5~Md€).

\figbloc{fon_03_inventaire}{Les prélèvements assis sur l'immobilier en 2024}{fig:inv}

Selon l'OCDE, la France prélevait en 2023 en impôts périodiques sur l'immobilier 1,90~\% du PIB, le double
de la moyenne des pays membres (0,95~\%), et en impôts sur les transactions 0,72~\% contre 0,37~\%. Elle taxe
donc lourdement les deux. L'argument de cette note n'est pas que la détention serait peu taxée, mais qu'un
euro de droits de mutation coûte davantage à l'économie qu'un euro de taxe foncière, surtout s'il porte sur
le sol.

\textbf{Les droits de mutation.} Pour un logement ancien, ils cumulent la part départementale (4,50~\%),
la taxe communale (1,20~\%) et les frais d'assiette de l'État (2,37~\% de la part départementale), soit
5,81~\% du prix. La loi de finances pour 2025 a permis aux départements de porter leur part à 5~\% pour les
ventes conclues du 1\textsuperscript{er}~avril~2025 au 31~mars~2028 ; 88 collectivités sur une centaine
l'ont fait, ce qui porte le taux à 6,32~\%. Les primo-accédants achetant leur résidence principale sont
exclus de cette majoration ; la possibilité d'une réduction propre aux primo-accédants n'est utilisée que
par la Savoie.

\textbf{La taxe foncière.} Elle repose sur la moitié de la valeur locative cadastrale, établie selon les loyers
de 1970 et revalorisée depuis par des coefficients nationaux. La révision des valeurs des logements a été
reportée à 2031 par la loi de finances pour 2026. Rapportée à la valeur de marché, elle pèse 0,52~\% dans le
dixième des logements situés dans les communes les moins chères et 0,18~\% dans le dixième le plus cher
(figure~\ref{fig:tfv}) ; à Paris, 0,12~\%, malgré la hausse du taux de 13,5 à 20,5~\% en 2023. L'Insee aboutit au même constat par une autre voie : la taxe
foncière est régressive rapportée au patrimoine immobilier et progressive rapportée au revenu (André et
Meslin, 2023).

\subsection{Ce que dit la recherche}

\subsubsection*{Les droits de mutation réduisent les ventes et la mobilité}

Trois résultats sont solides. D'abord, les ventes réagissent fortement au calendrier : le mois qui précède
une hausse, elles bondissent de 26~\% en France en 2014 (Bérard et Trannoy, 2018), de 41~\% en Allemagne ;
ce sont des ventes avancées, pas des ventes créées. Ensuite, au-delà de ces effets de calendrier, la plupart
des études trouvent que chaque point de droits supprime environ 8~\% des ventes, dans une fourchette de 0 à
18~\% et sur des horizons de quelques mois à quatre ans (Dachis, Duranton et Turner, 2012, à Toronto ; Fritzsche et Vandrei, 2019, en
Allemagne ; Best et Kleven, 2018, au Royaume-Uni). Enfin, la mobilité des propriétaires baisse de 8 à 14~\%
par point (Pays-Bas, Finlande, Toronto). En Finlande, ce sont surtout les déménagements courts et les
déménagements vers plus grand qui disparaissent, pas les réductions de taille (Eerola et al., 2021) : la taxe
frappe le ménage qui veut une pièce de plus, ce qui est souvent le cas quand la famille s'agrandit (l'étude
n'observe pas les naissances).

Pour la France, les preuves sont plus minces et partagées. Après la hausse de 0,7~point de 2014, Bérard et
Trannoy trouvent une baisse des ventes de 6~\% sur trois mois, puis plus rien ; Bachelet et Poulhès (2017)
trouvent une baisse durable de 10 à 12~\%, de 20~\% dans les marchés détendus et nulle dans les marchés
tendus, et un prix hors droits inchangé : l'acheteur a supporté toute la hausse. En Allemagne, au contraire,
une hausse d'un point des droits réduit les prix d'environ 3~\% : les vendeurs la supportent (Dolls, Fuest,
Krolage et Neumeier, 2025). L'incidence d'une baisse est donc incertaine : à court terme, elle profite
plutôt aux acheteurs ; à plus long terme, elle passe en partie dans les prix au profit des vendeurs, et une
exonération ciblée sur les primo-accédants au Royaume-Uni a augmenté leurs achats de 11~\% tout en relevant
de 1~\% les prix qu'ils payaient (HMRC, 2023).

Le coût économique des droits de mutation croît vite avec leur taux, comme celui de tout impôt : autour de
1 à 2~\%, il représente 5 à 30~\% de la recette ; vers 4 à 5~\%, il approche 0,7 à 0,9~€ par euro
supplémentaire collecté (Buettner, en Allemagne ; Määttänen et Terviö, 2022). Au taux français, l'extrapolation
du modèle finlandais donne de l'ordre de 1 à 2~€ par euro, mais personne ne l'a estimé pour la France. Le seul
modèle français (Bontemps, Cherbonnier et Magnac, 2023) trouve un gain modeste à remplacer les droits par
un impôt foncier (+0,2~\% de consommation), qui s'inverse si l'on finance la baisse par l'impôt sur le
revenu : c'est la raison de la reporter sur la détention. Dans ce modèle, la baisse profite surtout aux
jeunes ménages aisés, dont l'accès à la propriété progresse, alors que celui des jeunes ménages modestes
recule.

\subsubsection*{La taxe foncière : un bon impôt mal assis}

Le principal inconvénient d'un impôt récurrent sur l'immobilier est la transition. Une hausse permanente se
capitalise dans les prix : le propriétaire présent au moment de l'annonce perd la valeur actuelle de toutes
les hausses futures, et les acheteurs suivants achètent moins cher (Koster et Pinchbeck, 2022, pour la
council tax anglaise). C'est un transfert des propriétaires en place vers les futurs acheteurs. Pour une
taxe assise sur le seul sol, l'ampleur de cette capitalisation reste discutée (Danemark : complète selon
Høj et al., 2018, nulle selon Nielsson et al., 2024). Un ménage âgé riche en patrimoine mais modeste en
revenu peut aussi avoir du mal à payer ; les reports de paiement facultatifs sont peu utilisés, sauf quand
ils sont proposés par défaut, comme au Danemark.

L'assiette est l'autre difficulté. Les évaluations de masse sont souvent régressives, parce qu'elles suivent
mal la valeur des quartiers chers (Berry, 2021 ; Avenancio-León et Howard, 2022). Philadelphie a corrigé
trente ans de gel par une réévaluation générale en 2014, accompagnée d'un abattement et d'un dispositif
pour les occupants de longue date. Pittsburgh, à l'inverse, a abandonné en 2001 son double taux après une
réévaluation brutale qui a fait bondir la valeur du sol de 81~\% sans lissage. Un plafond de hausse attaché
au propriétaire plutôt qu'au bien crée un verrou : en Californie, la Proposition~13 retient les ménages dans
leur logement (Ferreira, 2010).

Le double taux, plus élevé sur le terrain que sur le bâti, a un effet modeste mais positif sur la
construction : en Pennsylvanie, un rapport de quatre pour un a accru d'environ 3~points la croissance
décennale du nombre de pièces (Banzhaf et Lavery, 2010). La part de l'impôt qui porte sur le bâti décourage l'investissement là où l'offre
de logements peut réagir (Lutz, 2015). Le Bade-Wurtemberg est allé au bout de la logique : depuis 2025, sa
taxe foncière ne porte que sur la valeur du sol, établie par des valeurs de zone, avec un taux réduit de
30~\% pour le logement ; la plus haute juridiction fiscale allemande l'a jugée conforme en mai 2026, des recours
constitutionnels restant pendants. À Karlsruhe, à recette constante, les maisons individuelles paient
76~\% de plus et les copropriétés 20~\% de moins : l'assiette sol charge les grandes parcelles et allège la ville
dense.

\subsection{La proposition et ce qu'elle donne}

\begin{encadre}{Comment la bascule est chiffrée}
\textbf{Champ.} Les 30,5~millions de logements des particuliers (parc du recensement moins les HLM) de France
métropolitaine hors Alsace-Moselle, où les ventes ne sont pas publiées. Leur valeur est de 6~760~Md€ : prix
médian au m² des ventes 2023-2024 de la commune (sinon de l'intercommunalité ou du département), multiplié
par la surface médiane des logements vendus.

\textbf{Part du terrain.} Pour les maisons, prix au m² des terrains à bâtir multiplié par la surface de la
parcelle ; quand la commune compte trop peu de ventes de terrains, le prix départemental est ajusté selon
l'écart de prix des maisons (élasticité estimée de 1,54 sur 3~300 communes). Pour les appartements, valeur
du logement moins le coût de reconstruction du bâti, calé pour que l'ensemble retrouve la part du terrain
des comptes de patrimoine des ménages (45~\%). Elle va de 16~\% dans le dixième le moins cher à 67~\% dans le
plus cher.

\textbf{Charges actuelles.} Taxe foncière : rapport taxe/valeur de chaque commune selon l'Insee (2017),
vieilli selon la croissance de la taxe foncière de la commune de 2017 à 2024 (fichiers REI), appliqué aux
valeurs de 2024 et calé sur le total (23,6~Md€ pour ce champ). Droits de mutation : 5,81~\% des
ventes de logements anciens, 11,4~Md€ par an en moyenne en 2023-2024.

\textbf{Ce qui n'est pas modélisé.} Les effets sur les prix (capitalisation), l'âge et le revenu des
propriétaires, les locaux professionnels, les logements sociaux et les résidences secondaires en tant que
telles.
\end{encadre}

\textbf{Première étape : diviser par deux les droits de mutation.} La mesure porte sur les logements ; les
locaux professionnels, qui portent le total des droits à 16,7~Md€ en 2025, ne sont pas traités ici. Le taux
passe de 5,81~\% (taux de droit commun, une fois éteinte la majoration temporaire à 6,32~\% votée jusqu'en
mars 2028) à 2,9~\%. La perte de
recette est de 5,7~Md€ si les ventes ne réagissent pas, de 4,6~Md€ si chaque point de baisse augmente les ventes
de 6~\%, valeur centrale de la littérature (4,3 à 5,7~Md€ selon l'hypothèse). Une baisse de 30~\% coûterait 2,6
à 3,4~Md€. La baisse doit s'appliquer dès son annonce, sans délai : sinon les ventes s'arrêtent jusqu'à son
entrée en vigueur. La Finlande l'a fait en 2023 en la rendant rétroactive à la date du projet de loi. Par
prudence, on finance la baisse comme si les ventes ne réagissaient pas ; leur reprise, si elle a lieu, est un
gain.

\textbf{Deuxième étape : reporter la recette sur la taxe foncière.} Trois manières de le faire ont des
effets très différents (figures~\ref{fig:tfv} et~\ref{fig:var}).

La plus simple relève la taxe foncière actuelle de 24~\% partout (20~\% si les ventes réagissent) : entre 110
et 240~€ de plus par logement et par an. Elle ne bouleverse rien, mais elle aggrave l'injustice des
valeurs de 1970 : la taxe atteindrait 0,65~\% de la valeur dans les communes les moins chères contre 0,22~\%
dans les plus chères.

La plus ambitieuse remplace la taxe foncière par un impôt national à deux taux sur la valeur de marché :
0,73~\% sur le terrain et 0,18~\% sur le bâti. Elle supprime les distorsions de 1970, mais d'un coup : la taxe
augmenterait de 215~\% en moyenne dans le dixième des logements le plus cher et un logement sur quatre
verrait sa taxe augmenter de plus de moitié. Sans une transition de plusieurs décennies, c'est impossible.

\figbloc{fon_01_tf_valeur}{La taxe foncière rapportée à la valeur des logements, aujourd'hui et après la bascule}{fig:tfv}

\textbf{Notre proposition} combine les deux logiques. Chaque commune garde sa recette de taxe foncière et son
pouvoir de taux, mais l'assiette devient la valeur de marché, avec un taux sur le terrain égal à quatre fois
le taux sur le bâti, dans la fourchette pratiquée en Pennsylvanie. La perte de droits de mutation est
remplacée par un taux national unique de 0,19~\% sur la valeur des terrains, reversé aux départements selon
leurs pertes ; si les ventes repartent, il peut descendre vers 0,15~\%. Le principal effet économique est
de déplacer l'impôt du bâti vers le sol : la part de la taxe foncière qui porte sur les constructions passe
d'environ 14 à 7~Md€. Le taux national évite une difficulté
sérieuse : si chaque département remplaçait seul ses droits perdus, il lui faudrait un taux de 0,12~\% de la
valeur du sol à Paris mais de 0,85~\% dans la Creuse, où les terrains valent peu.

\figbloc{fon_02_variation}{Variation de la charge par logement et par an, selon le prix au m² de la commune}{fig:var}

La taxe foncière augmente de 3~\% dans le dixième des logements le moins cher et de 69~\% dans le plus
cher. En comptant la baisse des droits de mutation, rapportée à l'ensemble des logements de la commune, la
charge totale varie de $-$63 à +133~€ par logement et par an selon le décile. Rapportée à la valeur, la taxe
va de 0,54~\% dans les communes les moins chères à 0,30~\% dans les plus chères : l'écart actuel se resserre
sans être renversé. Les hausses se concentrent : un logement sur quatre (23~\%) voit sa taxe foncière
augmenter de plus de 300~€, et un sur vingt (5~\%, surtout à Paris et dans les communes littorales les plus
chères) la voit plus que doubler. Le partage entre maisons et appartements, lui, n'est pas robuste : il
dépend de la façon de mesurer le terrain. Avec notre méthode, les appartements paient davantage (+240~€
contre +150~€) ; si l'on mesure le terrain des maisons comme celui des appartements, par différence entre la
valeur et le coût de reconstruction, c'est l'inverse (+130~€ contre +230~€). Il faudra simuler l'assiette
légale, valeur de zone multipliée par la parcelle et répartie entre copropriétaires, avant de conclure.

Le gain d'efficacité est réel mais modeste. Si chaque euro de droits supprimé coûte en moyenne 0,3 à 0,9~€ à
l'économie, et chaque euro de taxe sur le sol presque rien, déplacer 5,7~Md€ rapporte de l'ordre de 1,5 à
4~Md€ par an, 0,05 à 0,15~\% du PIB : un ordre de grandeur, pas une estimation.

\figbloc{fon_04_primo}{Ce que gagne un ménage qui achète le logement médian de sa commune}{fig:primo}

\textbf{Pour les jeunes acheteurs, le gain dépend de qui supporte l'impôt} (figure~\ref{fig:primo}). Si
les prix ne bougent pas, comme après la hausse française de 2014, un ménage qui achète le logement médian de
sa commune paie entre 2~300~€ (un appartement à Saint-Étienne) et 22~700~€ (une maison à Arcachon) de droits
en moins : 8~100~€ pour un appartement à Lyon, 13~000~€ pour une maison à Nantes, 12~100~€ pour un
appartement à Paris. La hausse de taxe foncière qu'il paie ensuite (de $-$100 à +660~€ par an selon la
commune) ne rattrape cette économie qu'au bout de 17 à 23~ans dans les grandes villes, davantage ailleurs.
Le perdant est alors le propriétaire installé qui ne vend pas : il paie la hausse sans profiter de la baisse,
jusqu'à 660~€ par an pour l'appartement médian parisien, 22~000~€ en valeur actuelle.

Si au contraire les deux mouvements passent dans les prix, la baisse des droits les fait monter et la hausse
de taxe foncière les fait baisser ; à recette constante, le prix moyen bouge peu. La bascule transfère alors
de la valeur entre logements, des plus fonciers et des moins souvent vendus vers les autres, plutôt qu'entre
générations : l'appartement médian parisien perdrait environ 8~000~€, 2~\% de sa valeur. Dans les deux cas,
l'argument le plus solide en faveur des jeunes est la liquidité : les droits se paient comptant, sur
l'apport. Si l'apport représente 10 à 20~\% du prix, la baisse de 2,9~points en représente un septième à un
tiers.

\textbf{Pour les familles}, le gain tient au type de déménagement que les droits découragent le plus :
l'agrandissement. Avec des droits divisés par deux, passer d'un trois-pièces à un quatre-pièces à la
naissance d'un deuxième enfant coûte la moitié moins en impôt. Aucune étude
ne relie directement les droits de mutation à la natalité : l'effet attendu passe par le logement, pas par
une incitation directe.

\textbf{Pour les ménages les plus aisés}, l'effet net est limité. La taxe foncière augmente davantage dans
les communes chères (en moyenne 560~€ par an dans le dixième le plus cher), mais
les acheteurs de biens chers économisent davantage de droits. L'impôt sur la fortune immobilière n'est pas
modifié.

\subsection{Comment passer d'un système à l'autre}

\textbf{Construire l'assiette.} La France n'a pas de valeurs du sol par zone, comme les
\emph{Bodenrichtwerte} allemandes. Les ventes publiées (DVF) et leur version enrichie réservée aux acteurs
publics (DV3F) permettent de les établir par commune ou par quartier, en combinant les ventes de terrains et
les prix des logements, comme le fait cette note. L'administration fiscale établirait des valeurs de zone,
révisées tous les trois ans et lissées sur trois ans ; le propriétaire pourrait contester par expertise
quand sa valeur s'écarte nettement du marché, comme en Bade-Wurtemberg au-delà de 30~\%. Il faut écarter
la révision de la valeur à chaque vente, adoptée en commission le 9~octobre~2026 (amendement I-CF1872), puis
tombée avec le rejet de la première partie du projet de loi de finances : elle créerait le même verrou que la Proposition~13.

\textbf{Un calendrier.} Les valeurs du sol n'existent pas encore ; le Bade-Wurtemberg a mis cinq ans à les
établir (loi de 2020, premier impôt en 2025). Or la baisse des droits ne peut pas attendre : annoncée pour
plus tard, elle gèlerait les ventes jusqu'à son entrée en vigueur. Elle s'applique donc dès son vote et
elle est financée, le temps d'établir les valeurs, par une hausse uniforme et temporaire de la taxe foncière
actuelle (20 à 24~\%), remplacée ensuite par la nouvelle assiette. Celle-ci se substitue à la révision des
valeurs locatives reportée à 2031. Le taux national sur le terrain s'applique en totalité dès la première
année : il est faible (0,19~\% de la valeur du sol). Seule la nouvelle répartition de la taxe communale entre
logements est étalée, sur dix ans, hausses et baisses au même rythme, pour que la recette de chaque commune
reste constante chaque année. Les propriétaires âgés aux revenus modestes, déjà exonérés au-delà de 75~ans
sous condition de ressources, se verraient proposer par défaut un report de paiement portant intérêt,
remboursé à la vente ou à la succession. Le seuil de ressources, fondé sur le revenu fiscal de référence,
devra être redéfini si l'impôt sur le revenu et la CSG fusionnent.

\textbf{Garder l'autonomie locale.} Les communes votent leur taux sur la nouvelle assiette ; la loi fixe le
rapport entre le taux sur le terrain et le taux sur le bâti. Le Conseil constitutionnel n'exige pas de
compensation à l'euro près quand une recette locale est supprimée (décision 2019-796 DC) ; le taux national
sur le terrain, reversé aux départements, assure néanmoins cette compensation. Elle a un prix : les
départements échangent la moitié de leur principal impôt dont ils votent le taux contre une dotation, et
certains territoires paient plus qu'ils ne perdent (Paris environ 270~M€, les Hauts-de-Seine 130~M€). La taxe
d'enlèvement des ordures ménagères, assise sur la base de la taxe foncière, devra en être découplée.

\textbf{Limites.} Les valeurs sont des médianes communales, pas des valeurs de biens ; la part du terrain
des appartements est un calage sur les comptes nationaux, que l'Insee obtient lui-même par différence ; la
charge actuelle de chaque commune repose sur un rapport de 2017, vieilli par la croissance de sa taxe
foncière. L'effet de la taxe foncière française sur
les loyers et la construction n'a jamais été mesuré, et la répartition par âge de la bascule reste à
établir sur des données individuelles.


\section{Impôt sur les sociétés : déduire l'investissement plutôt que baisser le taux}\label{sec:is}
% Partie IS (source : ../is/note-is.tex, corrigée selon ../revue/verif_is_pfu.md ; chiffres : ../is/sorties/simulation_is.txt, euros de 2027)
L'impôt sur les sociétés (IS) n'appelle pas de grande réforme. Ce que lui reprochent les comparaisons
internationales, c'est son taux affiché, le plus élevé de l'OCDE pour les plus grands groupes en 2025 et 2026
avec la contribution exceptionnelle, sa complexité et son instabilité, pas le traitement de
l'investissement, qui est dans la moyenne. La référence garde donc le taux normal à 25~\% ; la question est
de savoir s'il faut déduire l'investissement plus vite, et comment le payer.

\subsection{Ce que les comparaisons reprochent à l'IS français}

La Tax Foundation classe la France dernière des 38 pays de l'OCDE en 2025 et en 2026, au total comme pour
l'IS, à cause du taux de 36,1~\% des plus grands groupes (24,2~\% en moyenne), de quatre taux différents
(2,3 en moyenne), d'une surtaxe, d'une boîte à brevets et de la troisième subvention implicite à la
recherche la plus généreuse. Cet indice pénalise par construction toute incitation ciblée. Il ne pénalise
pas la déduction de l'investissement : la France est 15\up{e} sur 38 en 2025 et 16\up{e} en 2026 pour la
récupération des coûts, et les déductions d'une machine y valent 88~\% de son prix en valeur actuelle, contre
86~\% en moyenne. Les taux effectifs de l'OCDE vont dans le même sens. Le taux moyen effectif, qui compte pour
choisir où implanter un projet rentable, est de 23,6~\% contre 22,2~\% (13\up{e} sur 38) ; le taux marginal
effectif, qui compte pour décider d'un investissement de plus, est de 15,0~\% contre 16,5~\%, et de 9,8~\%
contre 10,4~\% pour les équipements (chiffres de 2025, qui semblent ignorer la contribution exceptionnelle).
Le Conseil des prélèvements obligatoires (2023, 2025) et le Sénat (2026) font de l'instabilité le premier
grief : la contribution exceptionnelle, présentée comme temporaire, a été prorogée une fois, et son assiette
est la moyenne de l'IS de deux exercices, dont l'un était clos quand elle a été votée.

\subsection{Déduire l'investissement plutôt que baisser le taux}

Une entreprise achète une machine de 100~€ qui dure sept ans. Elle en déduit aujourd'hui le coût par
fractions ; actualisées à 5~\%, ces déductions valent 85 à 90~€ selon que l'amortissement est linéaire ou
dégressif. Si elle peut déduire les 100~€ l'année de l'achat, le total ne change pas, seulement sa date ; mais
l'État rend alors 25~€ à l'achat et prélève ensuite 25~\% de ce que la machine rapporte. Il devient un
associé, à condition que les pertes soient elles aussi déduites sans limite, ce que le plafonnement de leur
imputation empêche en partie. Un investissement qui rapporte tout juste le rendement normal du capital ne
paie plus d'impôt ; restent imposés les profits qui dépassent ce rendement et ceux du capital déjà installé.
Une baisse du taux allège au contraire tous les profits, y compris ceux des machines achetées il y a dix ans et
ceux des rentes, qui ne dépendent pas de la décision d'investir. Sur la réforme américaine de 2017, un point
de capital supplémentaire au bout de dix ans coûte environ 10~\% de la recette de l'IS par an avec une baisse
du taux et 2,5~\% avec la déduction immédiate (Chodorow-Reich, Smith, Zidar et Zwick, 2026). L'effet sur
l'investissement est établi : +10 et +17~\% d'investissement éligible pour les déductions partielles
américaines de 2001-2004 et 2008-2010, deux fois plus dans les petites entreprises et seulement chez celles
qui payaient l'impôt (Zwick et Mahon, 2017) ; +11~\% pour les entreprises moyennes britanniques (Maffini, Xing
et Devereux, 2019). Ces mesures étaient temporaires et ont en partie avancé des achats ; d'où l'importance
d'une mesure permanente. Il n'existe aucune évaluation causale du suramortissement français de 2015.

Le coût d'usage du capital, ce que doit rapporter une machine pour couvrir son financement, son usure et
l'impôt, s'écrit $c = (r+\delta)(1-\tau z)/(1-\tau)$, où $z$ est la valeur actuelle des déductions pour 1~€.
Avec $z$ entre 0,85 et 0,90, l'impôt le majore de 3,5 à 5~\% ; avec la déduction immédiate, plus du tout. Si
le stock d'équipements des sociétés imposables (environ les trois quarts des 541~Md€ des sociétés non
financières) réagit à peu près proportionnellement au coût d'usage (élasticité de 0,5 à 1,6), le PIB gagnerait
0,04 à 0,2~\% à long terme, 0,1~\% au centre. Ce gain de production n'est pas un gain de bien-être du même
montant : il rémunère d'abord le capital supplémentaire qui le produit. La perte sèche évitée, le triangle
sur ce capital, n'est que de 0,05 à 0,16~Md€ par an, parce que le coin supprimé est petit ; l'essentiel de
l'intérêt de la mesure tient à ce qu'elle coûte peu.

\subsection{Combien cela coûte, et comment le payer}

Les sociétés non financières ont investi 84~Md€ en 2025 en machines, équipements et matériels informatiques ;
on suppose que les trois quarts sont faits par des sociétés imposables, hypothèse non mesurée à laquelle le
coût est proportionnel. La première année (2027), une déduction immédiate totale coûterait 14 à 15,5~Md€ ;
le coût baisse ensuite, parce que les machines achetées après la réforme n'ont plus d'amortissement à
déduire. À partir de la huitième année ne reste qu'un coût permanent de 1,4 à 2,0~Md€ par an, qui vient
seulement de ce que l'investissement croît. En valeur actuelle, une année d'investissement coûte 1,8 à
2,6~Md€ au taux d'actualisation des entreprises et 1,2 à 1,7~Md€ à celui de l'État, soit 0,5 à 0,8 point de
taux normal (un point vaut 3,2~Md€ en 2027 une fois le taux réduit supprimé).

\figbloc{is_01_cout_deduction}{Coût annuel de la déduction immédiate des équipements et ressources pour la financer}{fig:cout}

\begin{table}[!htbp]\centering\scriptsize
\caption{Les arbitrages de l'IS (Md€ de 2027 par an)}\label{tab:is_arb}
\begin{tabularx}{\linewidth}{@{}>{\raggedright\arraybackslash}p{3.4cm}>{\raggedright\arraybackslash}X>{\raggedright\arraybackslash}p{4.2cm}@{}}
\toprule
Choix & Ce qu'il change & Qui paie \\
\midrule
Déduction immédiate : 50~\% ou 100~\% & 50~\% : 7 à 7,8 la 1\up{re} année, gain de PIB divisé par deux ; 100~\% : 14 à 15,5, puis 1,4 à 2 & --- \\
Calendrier : 100~\% dès 2027, ou 50~\% jusqu'en 2032 puis 100~\% & dette maximale (intérêts à 3,25~\%) : 17 à 45 ou 2 à 32 ; remboursée en 10 à 28 ou 4 à 25~ans & --- \\
Financement par l'assiette & Fin du taux réduit (2,2 à 3,1), CIR apprécié au niveau du groupe (0,6), fin des exonérations zonées hors outre-mer (0,55), boîte à brevets à 17,5~\% (0,2 à 0,6) : 3,5 à 4,9 & petites sociétés ; grands groupes de recherche ; 1~065 utilisateurs de la boîte à brevets \\
Financement par le taux (28~\%) & 5,1 à 8,7 ; perte sèche +0,04 sur le capital et 0 à +3,8 selon le coût réel du transfert de bénéfices & toutes les sociétés, surtout les profits mobiles \\
Taux réduit maintenu & il manque 2,2 à 3,1 ; la dette de transition n'est jamais remboursée & --- \\
Taux réduit réservé aux sociétés opérationnelles & rendement non chiffrable (aucune ventilation publiée) & holdings patrimoniales, sociétés immobilières à l'IS \\
Excédent une fois la dette remboursée & rendu par une baisse du taux normal de 0,6 à 1,2 point & --- \\
Contribution exceptionnelle & s'éteint (référence) : $-$7 par rapport à 2026, $-$5 par rapport au PLF 2027 & environ 450 grands groupes \\
\bottomrule
\end{tabularx}
\end{table}

\textbf{Le taux ou l'assiette.} Une hausse du taux normal de 25 à 28~\% rapporterait 8,7~Md€ à bénéfices
inchangés et 5,1~Md€ si les bénéfices déclarés baissent de 1,5~\% par point, l'ordre de grandeur trouvé sur les
filiales de multinationales (Beer, de Mooij et Liu, 2020). L'objection n'est pas son montant, mais ce qu'elle
frappe : la déduction immédiate cesse d'imposer le rendement normal des investissements nouveaux, alors que le
taux porte d'abord sur les profits exceptionnels, qui sont aussi les plus mobiles. Un groupe choisit le pays
où il installe un projet très rentable selon le taux moyen effectif, et celui où il déclare ses bénéfices
selon le taux affiché. Si les bénéfices déplacés sont une perte réelle pour la France, la perte sèche de
28~\% atteint plusieurs milliards ; si ce ne sont que des écritures, elle est presque nulle (tableau~\ref{tab:is_arb}).
Le financement par l'assiette (3,5 à 4,9~Md€) dépasse le coût permanent de la déduction ; l'excédent rembourse
la transition, puis peut être rendu par une baisse du taux normal de 0,6 à 1,2 point.

\textbf{La transition} est la vraie difficulté. Une déduction totale dès 2027, financée par l'assiette,
creuserait une dette de 17 à 45~Md€, intérêts compris, remboursée en 10 à 28~ans. Un calendrier voté dès 2027,
50~\% jusqu'en 2032 puis 100~\%, limite la dette à 2 à 32~Md€, remboursée en 4 à 25~ans ; il évite de laisser le
second palier à l'appréciation des budgets futurs, ce qui reproduirait l'instabilité reprochée à l'IS.

\textbf{La fin du taux réduit : qui y perd.} Les sociétés dont le chiffre d'affaires est inférieur à 10~M€ et
le capital détenu à 75~\% par des personnes physiques paient 15~\% au lieu de 25~\% sur les 42\,500 premiers
euros de bénéfice ; depuis une décision du Conseil d'État de 2025, le chiffre d'affaires d'une société
membre d'un groupe s'apprécie au niveau du groupe. Le dispositif coûte environ 2,7~Md€ (bénéfices de 2022).
Pour une société au-dessus du plafond, il revient à verser chaque année 4\,250~€, quoi qu'elle fasse : la
somme ne dépend d'aucune décision, sauf celle de rester petite ou de se découper (on observe un pic de
sociétés juste sous le plafond). Mais beaucoup de bénéficiaires sont sous le plafond (environ 24\,000~€ de
bénéfice imposé à 15~\% en moyenne en 2018) : pour eux, la suppression relève l'impôt de chaque euro
supplémentaire, de deux tiers en tout (3\,000~€ pour 30\,000~€ de bénéfice), et renchérit tout leur capital.
Bernini et Treibich (2016) trouvent que l'introduction du taux réduit en 2001-2003 a accru de 29~\% le
capital des petites sociétés et leurs exportations, mais ralenti leur emploi (résumé seulement lu). La
déduction immédiate ne compense qu'une partie : pendant la phase à 50~\%, elle avance 2\,300~€ d'impôt la
première année pour 20\,000~€ d'équipements, et ne compense durablement la perte qu'au-delà de 220\,000~€
d'investissement annuel (110\,000~€ à 100~\%). Perdent le plus les sociétés qui investissent peu en
équipements : professions libérales, conseil, petites holdings et sociétés immobilières à l'IS ; aucune
statistique ne ventile la base imposée à 15~\% par type de société. Pour l'associé qui se fait distribuer le
bénéfice, le taux cumulé de l'IS et du PFU sur les 42\,500 premiers euros passe de 41,7 à 48,6~\%, celui de
toutes les autres sociétés.

\textbf{Ce que la référence écarte.} Un barème progressif de l'IS : la taille d'une entreprise ne dit rien de
la richesse de ses propriétaires, et les seuils multiplieraient les découpages. Un taux minimal national :
le minimum mondial de 15~\% s'applique déjà aux groupes de plus de 750~M€ de chiffre d'affaires (loi de
finances pour 2024, art.~33) ; pour les autres, un minimum assis sur l'impôt de l'année frapperait les
entreprises qui investissent beaucoup cette année-là. La contribution exceptionnelle s'éteint : surtaxe en
partie rétroactive et renouvelée d'année en année, elle fait payer les investissements passés et rend
incertain le rendement des investissements futurs.


\section{Revenus du capital : garder le taux, fermer la sortie sans impôt}\label{sec:pfu}
% Partie PFU (source : ../pfu/note-pfu.tex, corrigée selon ../revue/verif_is_pfu.md)
Le prélèvement forfaitaire unique (PFU) impose depuis 2018 les dividendes, les intérêts et les plus-values
mobilières à un taux unique, 31,4~\% en 2026, au lieu du barème progressif. La référence le garde tel qu'il
est :
\begin{itemize}\setlength{\itemsep}{0pt}
\item Quand le taux a monté en 2013, les actionnaires qui contrôlent leur société ont laissé les bénéfices
dedans. Quand il a baissé en 2018, ils les ont distribués.
\item Cette réaction porte sur le moment où l'on se verse le revenu plus que sur l'activité. Elle dépend sans
doute du taux qu'on attend pour demain : le principal service à rendre est de ne plus le changer.
\item La faiblesse du système n'est pas dans le taux du PFU, mais dans le circuit des holdings.
\end{itemize}

\subsection{Le PFU en 2026}

Le PFU additionne 12,8~\% d'impôt sur le revenu et 18,6~\% de prélèvements sociaux.
\begin{itemize}\setlength{\itemsep}{0pt}
\item Depuis la LFSS pour 2026, la CSG sur les revenus du capital est de 10,6~\%. Elle reste à 9,2~\% pour une
liste fermée de revenus (assurance-vie, épargne logement, revenus fonciers, plus-values immobilières) : deux
taux coexistent, sans justification économique.
\item Le contribuable peut opter pour le barème, globalement et, depuis 2026, de façon révocable. Sur leurs
revenus de 2024, 2,5~millions de foyers l'ont fait.
\item En haut de la distribution, la contribution exceptionnelle sur les hauts revenus (CEHR) ajoute 3 puis
4~points au-delà de 250\,000~€ de revenu (500\,000~€ pour un couple). La contribution différentielle (CDHR)
garantit un impôt d'au moins 20~\% du revenu, ce qui porte à 38,6~\% le taux d'un foyer qui vit de
dividendes ; elle est prorogée jusqu'à ce que le déficit du budget de l'État repasse sous 3~\% du PIB.
\end{itemize}

Le PFU ne s'apprécie qu'avec l'IS déjà payé par la société. Sur 100~€ de bénéfice, une société paie 25~€
d'IS (25,8~€ si son impôt dépasse 763\,000~€, avec la contribution sociale).
\begin{itemize}\setlength{\itemsep}{0pt}
\item Si elle distribue le reste, l'actionnaire paie 31,4~\%. Le bénéfice distribué supporte donc 48,6 à
49,1~\% au total, et jusqu'à 54,5~\% pour un foyer à la CDHR (figure~\ref{fig:taux}). C'est proche du taux
marginal supérieur sur un salaire (55,4~\% hors cotisations sociales).
\item Ce niveau est élevé en Europe. La France est sixième sur 35 pour le taux personnel supérieur sur les
dividendes (Tax Foundation, 2026, avec le taux de 34~\% de 2025) ; à cinq fois le salaire moyen, l'OCDE
trouvait en 2023 un taux moyen intégré de 49~\% que seuls le Danemark et l'Espagne dépassent.
\item Si la société verse son dividende à une holding de l'actionnaire qui le garde, la holding ne paie que
1,25~\% (régime mère-fille) et le bénéfice n'a supporté que 26,8~\%. C'est un report, pas un taux final :
distribué un jour, il paiera le PFU.
\end{itemize}
Imposer le bénéfice quand il sort des sociétés est d'ailleurs, en principe, ce qui gêne le moins la
croissance, puisque le capital circule d'une entreprise à l'autre sans impôt. Le report ne devient une
exonération que si la sortie n'a jamais lieu, ce que permet l'effacement au décès.

\figbloc{pfu_01_taux_cumule}{Impôt total sur 100~€ de bénéfice, selon ce qu'en fait l'actionnaire (2026)}{fig:taux}

\subsection{Ce que 2013 et 2018 ont appris}

Les dividendes ont suivi les deux réformes, en sens inverse :
\begin{itemize}\setlength{\itemsep}{0pt}
\item En 2013, l'option pour un prélèvement libératoire disparaît et tous les dividendes passent au barème.
Les foyers qui avaient opté pour le prélèvement libératoire déclarent 40~\% de dividendes en moins, sans
hausse de leurs salaires ; les entreprises contrôlées par des particuliers distribuent moins, accumulent des
actifs financiers et n'investissent pas moins (Bach et al., 2019).
\item Le taux marginal supérieur intégré n'était pourtant passé que de 57,6 à 60,1~\% : les actionnaires ont
pu juger la hausse provisoire et attendre, ou garder le bénéfice dans la société pour le sortir plus tard en
plus-value ou le transmettre (Bach et al., 2024).
\item En 2018, les dividendes remontent (+15~\% dès la première année). Le coût de la réforme tombe de 0,9 à
0,4~Md€ selon l'IPP (de 1,05 à 0,74~Md€, voire un gain dans l'hypothèse haute, selon l'Insee).
\item Mais 2013 et 2018 forment un seul épisode : les bénéfices retenus de 2013 à 2017 ont été distribués à
partir de 2018. L'« autofinancement » que France Stratégie (2023) déduit d'une part causale supposée, faute
de pouvoir la mesurer, ne vaut que pour les premières années, rattrapage compris.
\end{itemize}

Pour une société qui investit ses bénéfices, l'impôt sur le dividende ne change pas le choix d'investir :
100~€ distribués aujourd'hui rapportent 68,6~€ nets ; réinvestis à 5~\%, ils donneront 105~€ l'an prochain,
soit 72,0~€ nets. Le PFU prend la même fraction dans les deux cas : il change seulement la date et la forme
du revenu. Le raisonnement ne vaut pas pour une jeune entreprise qui lève des fonds propres à l'extérieur.
Les effets réels mesurés sont minces :
\begin{itemize}\setlength{\itemsep}{0pt}
\item après 2018, l'effet le plus net est une création d'entreprises un peu plus rapide dans les secteurs
exposés, qu'on ne sait pas séparer de celui de la baisse de l'IS (IPP, 2023) ;
\item la baisse américaine de 2003 n'a pas davantage modifié l'investissement (Yagan, 2015).
\end{itemize}

Sur le niveau du taux :
\begin{itemize}\setlength{\itemsep}{0pt}
\item Le taux qui maximise les recettes sur les revenus du capital, prélèvements sociaux compris, est estimé
à 43~\% en comptant la réaction des revenus du travail, à 57~\% sans elle (Lefebvre, Lehmann et Sicsic,
2025). Le PFU est en dessous, mais la marge est mince : avec l'élasticité estimée (0,77), un passage de 31,4
à 33~\% ne garde qu'un peu moins des deux tiers de son rendement mécanique.
\item Cette élasticité ne mesure pas une perte sèche : la réaction des dividendes consiste surtout à laisser le
bénéfice dans la société, ce qui reporte l'impôt sans détruire d'activité.
\item Les dividendes sont concentrés : 1~\% des foyers en déclarent 96~\% (2021). L'exil est faible (un point
d'imposition effective en plus fait partir à long terme 0,02 à 0,23~\% du centile aux plus hauts revenus du
capital ; Bach et al., 2025) et non mesuré au-delà.
\item La commission d'enquête de l'Assemblée sur les plus hauts patrimoines recommande de « maintenir la
stabilité » des taux (rapport n°~3056, 2026).
\end{itemize}

\subsection{Les arbitrages}

\textbf{Le taux.}
\begin{itemize}\setlength{\itemsep}{0pt}
\item Un PFU à 25~\% coûterait au plus 5,1~Md€ à comportements inchangés (un majorant, à cause des foyers qui
optent pour le barème) et 3,2 à 3,9~Md€ après réaction. Le gain irait d'abord au centile le plus aisé
(44~\% du gain du PFU selon l'estimation du Gouvernement en 2017).
\item Cette baisse ne rapporterait rien en activité : après 2018, aucun effet n'a été détecté sur
l'investissement ni sur les salaires des entreprises les plus exposées (France Stratégie, 2023). Elle
coûterait le prix des 3,7~Md€ à lever ailleurs, environ +0,3~Md€ de perte sèche (0,15 à 0,6).
\item La réaction des dividendes est une affaire de recette, pas d'efficacité : si une partie était une
évasion définitive que la baisse ferait cesser, la recette perdue serait moindre, sans gain d'activité.
\item Une hausse rapporterait peu : 1,2~Md€ nets pour 3~points, de 30 à 33~\%, selon le CAE (2025).
\end{itemize}

\textbf{Le report de l'apport-cession.} Un entrepreneur qui apporte ses titres à une holding avant de les
vendre place sa plus-value en report : elle n'est imposée que s'il vend les titres de la holding, ou si la
holding vend dans les trois ans sans réinvestir au moins 70~\% du produit. Au décès, le report devient une
exonération définitive ; en cas de donation, il s'efface si le donataire ne vend pas pendant six ans (onze
dans certains cas). D'après le rapport n°~3056 :
\begin{itemize}\setlength{\itemsep}{0pt}
\item 33,6~Md€ de plus-values ont été placées en report en 2023 et 22,2~Md€ en 2024, plus que les plus-values
imposées au PFU en 2024 (18,8~Md€ d'après les déclarations) ;
\item plus des deux tiers des entrées de 2023 (23,5~Md€) se concentrent sur 84 foyers ;
\item le stock atteint 153~Md€ en 2023 et 173~Md€ en 2024, et seuls 0,4 à 1,5~Md€ de reports ont expiré
chaque année de 2022 à 2024.
\end{itemize}
Trois traitements sont possibles :
\begin{itemize}\setlength{\itemsep}{0pt}
\item \emph{L'expiration à toute transmission} (PLF 2027, art.~5) : 0,5~Md€ annoncés pour 2027 et, selon
notre calcul, au moins 2~Md€ par an en régime de croisière. Mais elle taxe une structure plutôt qu'un gain
réalisé : une holding familiale transmise sous pacte Dutreil paierait sans avoir rien vendu, et l'étalement
sur cinq ans ne vise que l'impôt sur le revenu. Elle alourdit aussi de 90 à 100~M€ chacun la facture des plus gros apporteurs.
\item \emph{L'expiration seulement quand la holding vend.}
\item \emph{La transmission du report aux héritiers sans effacement} : la plus-value est imposée quand le
capital sort réellement des sociétés, sans impôt sans liquidité ni hausse de facture à la transmission. Elle
rapporte moins et plus tard, et sa solidité juridique reste à expertiser.
\end{itemize}

\textbf{Les autres choix.}
\begin{itemize}\setlength{\itemsep}{0pt}
\item Dans l'impôt fusionné, la CEHR disparaît pour les revenus du barème. La garder sur les revenus du
capital, sous la forme d'une majoration de 3 puis 4~points du PFU, maintient le taux actuel (tant que la CDHR
existe, un foyer qui vit de dividendes, déjà au plancher de 20~\%, n'y gagnerait de toute façon rien). Durcir
la CDHR, comme le propose la commission d'enquête, alourdirait la facture des mêmes foyers.
\item Aligner les deux taux de prélèvements sociaux à 18,6~\% rapporterait de l'ordre de 1,3~Md€, payés
surtout par la classe moyenne épargnante.
\item Un prélèvement modéré sur la trésorerie passive des holdings (l'Irlande taxe à 20~\% les revenus passifs
que les sociétés fermées ne distribuent pas dans les dix-huit mois) viserait le bon problème, mais
alourdirait lui aussi la facture des grandes fortunes. Un impôt de 25~\% sur les revenus non réinvestis des
holdings risque de heurter la directive mère-fille, selon l'administration fiscale, et, s'il est établi chez
l'associé, la jurisprudence constitutionnelle (imposer des bénéfices qu'il n'a pas perçus) ; il taxerait le
capital avant qu'il sorte des sociétés.
\item Pour une détention longue, la correction de principe est d'imposer le gain réel, en indexant le prix
d'acquisition sur l'inflation, plutôt que de rétablir des abattements pour durée de détention. Ces
abattements avaient compliqué l'impôt de 2013 à 2017.
\end{itemize}


\section{L'ensemble}\label{sec:ensemble}
% Bilan d'ensemble (chiffres : sorties/architectures_ircsg.txt, sorties/perte_seche_autres.txt,
% ../foncier/sorties/simulation.txt, ../is/sorties/simulation_is.txt)

\begin{table}[!htbp]\centering\scriptsize
\caption{Les quatre réformes, à recette constante (Md€ par an, régime permanent)}\label{tab:ensemble}
\begin{tabularx}{\linewidth}{@{}>{\raggedright\arraybackslash}p{3.2cm}>{\raggedright\arraybackslash}p{2.6cm}>{\raggedright\arraybackslash}p{2.4cm}>{\raggedright\arraybackslash}X@{}}
\toprule
Réforme & Recette & Gain d'efficacité (fourchette) & Qui paie plus \\
\midrule
Fusion IR-CSG-CRDS : franchise d'activité 4~000~€, 6~000~€ par enfant, petites pensions exonérées & 0 par
construction ; moindre optimisation : +1,9 à +5,7 & 4,9 (2,3 à 11,1) & retraités (+15~Md€) ; 29~\% des foyers
plus de 500~€ \\
Immobilier : droits de mutation $-$50~\%, taxe foncière à deux taux, taux national sur le sol & 0 (transition :
hausse temporaire de la taxe actuelle) & 1,9 (1,0 à 4,2) & propriétaires de terrains chers (Paris,
Hauts-de-Seine) \\
IS : déduction immédiate de l'investissement, financée par l'assiette & 0 en régime permanent ; dette de
transition de 2 à 32 & 0,1 (0,05 à 0,16), moins 0 à 0,8 si fin du taux réduit & sociétés au taux réduit,
crédit d'impôt recherche \\
PFU maintenu, report d'apport-cession & 0 ; recette différée selon le traitement du report & 0 & selon le
traitement du report \\
\midrule
Ensemble & 0 & environ 7 (2,5 à 15,5) & \\
\bottomrule
\end{tabularx}
\end{table}

\textbf{Le gain total est modeste et robuste} (tableau~\ref{tab:ensemble}). Il est d'environ 7~Md€ par an au
centre, 0,2~\% du PIB. Il pourrait être trois fois moindre ou deux fois plus grand selon les comportements.
Dans tous les jeux d'hypothèses, la fusion et la réforme de l'immobilier réduisent la perte sèche, et aucun
gain ne repose sur une réaction déclarative. Plus des deux tiers du gain viennent de la fusion de l'impôt sur le revenu et de la CSG, un peu plus
d'un quart du remplacement des droits de mutation. L'IS et le PFU servent surtout à ne pas aggraver les choses :
ne pas monter des taux que les bénéfices et les dividendes savent fuir, et rendre l'impôt prévisible. Le gain
de la fusion représente environ deux tiers de ce qu'un barème optimisé sans contrainte obtiendrait. Le dernier
tiers demanderait de rendre plus de 20~Md€ au dernier décile, sur la foi d'élasticités que personne n'a
mesurées pour un tel changement.

\textbf{Les grandes fortunes.} Trois profils sont à distinguer.
\begin{itemize}\setlength{\itemsep}{0pt}
\item \emph{Le dirigeant salarié.} Son taux marginal passe de 55~\% environ (avec la CEHR) à 50~\%. Le 1~\% le
plus aisé paie 2,7~Md€ de moins sur ses revenus du barème.
\item \emph{Le rentier.} Il garde le PFU, majoration comprise. Il paie plus de taxe foncière s'il possède des
terrains chers, et moins de droits quand il achète.
\item \emph{L'entrepreneur qui a apporté ses titres à une holding.} Son sort dépend du traitement du report
d'apport-cession (section~\ref{sec:pfu}). L'expiration à toute transmission alourdirait de 90 à 100~M€ la
facture de chacun des plus gros apporteurs ; la transmission du report sans effacement n'augmente rien à la
transmission.
\end{itemize}
Une incertitude demeure pour les 3\,800 foyers au-delà de 3~M€ de revenu fiscal : selon la part (non publiée)
de leurs revenus au PFU, la fusion baisse leur impôt sur le barème ou l'augmente de plusieurs points. Il faut
la lever avec les données de la DGFiP avant d'arrêter le barème.

\textbf{Les familles} gagnent : 1\,200~€ par an en moyenne pour un couple avec deux enfants, 230~€ pour un
parent isolé. La franchise de 6\,000~€ par enfant vaut 6\,000~€ fois le taux marginal du parent : 840~€ par
enfant sous 20\,000~€ de revenu imposé, 1\,800~€ vers 45\,000~€, davantage au-delà, contre 742~€ en moyenne
pour le quotient familial actuel. Les droits de mutation, payés au moment où une famille achète plus grand,
sont divisés par deux.

\textbf{Les retraités} sont les principaux contributeurs. La fusion leur fait payer 15~Md€ de plus,
2\,300~€ par an pour un couple, car elle met fin à leurs taux réduits et à leur abattement. Propriétaires plus
souvent que la moyenne, ils paient aussi la hausse de taxe foncière, sans profiter de la baisse des droits
s'ils ne déménagent pas. C'est le prix du gain d'efficacité : les pensions ne réagissent pas à l'impôt, et les
taxer comme les salaires permet de baisser ce qui pèse sur le travail. Trois leviers atténuent cette charge :
\begin{itemize}\setlength{\itemsep}{0pt}
\item l'exonération des petites pensions, retenue ici ;
\item le maintien de l'abattement de 10~\%, pour 0,6~point de taux pour tous ;
\item une transition sur quatre ans, et, pour la taxe foncière, le report de paiement des propriétaires âgés.
\end{itemize}

\textbf{Les dates.} Les chiffres décrivent le droit de 2026. L'IR et la CSG sont mesurés sur les revenus de
2024, l'IS en euros de 2027, l'immobilier sur les ventes de 2023-2024 et la fiscalité locale de 2024.


\clearpage
\appendix
\renewcommand{\thesection}{Annexe \Alph{section}}
\section{Fusion de l'IR, de la CSG et de la CRDS : calibrage, hypothèses et variantes}\label{ann:ircsg}
\IfFileExists{parties/annexe_ircsg.tex}{% Annexe de la partie IR-CSG (réécrite le 16/11/2026 ; base 2025, taux marginaux recalculés et cotisations retirées le
% 17/11/2026 ; coupée le 17/11/2026, CONSIGNES_COUPE.md). Chiffres : sorties/*_ircsg.txt, controles_purs.txt,
% controle_hauts_purs.txt, architectures_ircsg.txt ; tableaux : tables_ircsg.py ; figures : figures_ircsg_n.py.

\subsection{Champ, données et modèle}

\paragraph{Le champ.} En 2025, l'IR, la CSG et la CRDS prélèvent 229,9~Md€ sur les revenus imposés au barème : 86,6~Md€
d'IR, 135,8~Md€ de CSG (108,5 sur les revenus d'activité, 26,7 sur les pensions, 0,6 sur le chômage) et 7,6~Md€ de CRDS.
Les revenus du capital au prélèvement forfaitaire unique gardent leur régime (section~\ref{sec:pfu}) ; on leur impute
11,5~Md€ de l'IR de l'IRCOM des revenus 2024 (12,8~\% de 90~Md€ de revenus du capital hors barème, convention), 12,4~Md€
en 2025. Ces 90~Md€ comprenant 7,6~Md€ d'assurance-vie et 1,7~Md€ de plus-values immobilières, hors PFU, l'IR imputé est
surestimé d'environ 1,4~Md€ : corrigé, la recette à remplacer passerait à 231,3~Md€ et le 1~\% paierait environ 1,1~Md€
de moins (modèle non recalé ; section~\ref{sec:ens_haut}). La CASA, le prélèvement de solidarité sur les loyers et la
CEHR restent hors champ.

\paragraph{La base 2025.} La répartition est celle des revenus de 2024 (IRCOM), chaque montant étant porté en 2025 avec
la croissance de son agrégat, la même dans toutes les tranches : IR des comptes nationaux (96,0~Md€ en 2024, 103,6~Md€
en 2025, $+$7,9~\%), CSG sur les revenus d'activité ($+$2,3~\%) et de remplacement ($+$3,3~\%) (rapport « Financement »
annexé au projet de loi d'approbation des comptes de la Sécurité sociale de 2025), CRDS ($+$2,2~\%) ; les revenus
grandissent comme l'assiette de la CSG qui les frappe. La hausse de l'IR encaissé tient en partie au jeu des acomptes et
des soldes (Insee) ; le budget de l'État la confirme (IR net de 94,6~Md€, en hausse de 7,0~Md€, selon la commission des
finances du Sénat). Si elle est passagère, la recette à remplacer est surestimée d'environ 4,5~Md€ (0,4 point de taux).
Restent en 2024 : la répartition entre foyers, la composition des revenus, la CSG des pensions de départ, les niches et
le taux marginal actuel par tranche. Le droit de l'IR est celui des revenus 2025.

\paragraph{Les données et les foyers.} Les 25 tranches de revenu fiscal de l'IRCOM sont découpées en 60 foyers
représentatifs chacune ; la composition des revenus vient des totaux par case des déclarations 2042, la CSG d'une
reconstitution calée sur les comptes de la Sécurité sociale (25,4~Md€ sur les pensions en 2024, chiffre de la CCSS non
relu à la source). L'assiette est de 1\,520~Md€ après cotisations en 2025, dont 420~Md€ de pensions. Chaque foyer est
scindé en un foyer d'actifs et un foyer de retraités de même revenu, dans la proportion des pensions de sa tranche :
27,1~millions de foyers d'actifs, 14,5~millions de foyers de retraités. Les couples partagent leurs revenus selon
l'Insee (Insee Première 1492) ; les enfants ne sont que chez les actifs (un par parent isolé, deux par couple « avec
enfants »). La charge actuelle de chaque tranche est la charge observée, répartie entre types de foyers selon des
cas-types OpenFisca sans niche ; les retraités paient ainsi 50,2~Md€, 12,0~\% de leurs pensions, et les actifs 16,3~\%
de leurs revenus (ni la DGFiP ni la DREES ne publient l'impôt des foyers de retraités). Ce sont des ordres de grandeur,
pas une microsimulation : la dispersion réelle des gagnants et des perdants est plus forte.

\paragraph{Le nouvel impôt.} Les options du modèle (\texttt{modele\_purs.impot\_adultes}) sont la franchise d'activité
(salaires, indépendants et loyers, que le modèle ne sait pas séparer), l'exonération des petites pensions (totale
jusqu'à 12\,500~€, à peu près le minimum vieillesse de 12\,523~€ en 2026, puis réduite d'un euro par euro de pension,
d'un demi-euro dans la variante « étalée »), l'abattement sur les pensions, la franchise ou le chèque par enfant, les
quotients aux plafonds actuels et la déduction de l'épargne retraite (10,8~Md€ de versements en 2025), supposée
proportionnelle aux revenus d'activité de la moitié la plus aisée, faute de répartition publiée ; elle coûte 3,0~Md€,
soit 0,29 point de taux pour tous.

\paragraph{L'efficacité.} Le gain est la baisse de la perte sèche statique sur l'activité, à court terme (méthode v2,
\texttt{revue/methode\_v2.md}). Seul le revenu d'activité des actifs réagit ; chaque euro perdu coûte tout ce que la
collectivité prélevait dessus (annexe de méthode). La participation réagit au taux de prélèvement du premier salaire,
sous deux Smic ; la réaction du revenu déclaré du dixième le plus aisé est comptée comme un effet sur la recette
(tableau~\ref{tab:ircsg_jeux}).

\begin{table}[!htbp]\centering\scriptsize\setlength{\tabcolsep}{4pt}
\caption{Les jeux d'élasticités}\label{tab:ircsg_jeux}
\begin{tabular}{@{}lrrrr@{}}
\toprule
& Bas & Central & Haut & Long terme \\
\midrule
Revenu réel, moitié du bas & 0,02 & 0,05 & 0,10 & 0,10 \\
Revenu réel, 50 à 90~\% & 0,05 & 0,10 & 0,20 & 0,20 \\
Revenu réel, 90 à 99~\% et 1~\% du haut & 0,05 & 0,10 & 0,25 & 0,25 \\
Revenu déclaré, 90 à 99~\% / 1~\% (effet sur la recette) & 0,10 / 0,15 & 0,20 / 0,30 & 0,30 / 0,45 & 0,20 / 0,30 \\
Participation, sous deux Smic & 0,10 & 0,25 & 0,50 & 0,10 \\
\bottomrule
\end{tabular}
\par\vspace{3pt}\parbox{\linewidth}{\raggedright\scriptsize\color{encre}
Appuis : Lehmann, Marical et Rioux (2013), Cabannes, Houdré et Landais (2014), Sicsic (2022), Chetty, Guren, Manoli et
Weber (2013) ; \texttt{cours/elasticites-france.tex}. Ce sont des choix, pas des estimations. Jeu de long terme :
\texttt{revue/avis\_cochrane.md}. Déplacement du revenu déclaré vers l'IS et le PFU : la moitié, à 48,6~\%
(convention).}
\end{table}

\paragraph{Les cas-types et les taux marginaux.} OpenFisca-France 176.1, législation des revenus 2025
(\texttt{openfisca\_cas\_ircsg.py}) : salariés du privé, retraités nés en 1950, aucune niche, couples mariés ou pacsés,
Smic de 21\,622~€ brut par an (sources secondaires). Les seuils de CSG des retraités y sont de 12\,817, 16\,755 et
26\,004~€ de revenu fiscal pour une part, non relus ; ceux de 2026 (13\,048, 17\,057 et 26\,472~€) ont été vérifiés sur
service-public.fr. Le taux de CSG dépendant du revenu fiscal de l'avant-dernière année, qui dépend de la CSG déductible,
on suppose la même pension chaque année et l'on retient le taux le plus élevé dont le seuil est dépassé par le revenu
fiscal calculé à ce taux. Les taux marginaux sont calculés tous les 100~€, sans lissage ; leurs pics sont de vrais sauts
du droit : seuil de CSG appliqué à toute la pension, abattement des plus de 65 ans qui baisse d'un coup de 1\,411~€ à
chacun de ses deux seuils (vers 33\,500~€ de pension, 100~€ de plus coûtent 247~€) et, dans la proposition, sortie de
l'exonération des petites pensions.

\subsection{Situation actuelle et réforme naïve : compléments}

\figbloc{ircsg_n02_ocde}{La France prélève sur le revenu un peu plus que la moyenne de l'OCDE, pour près de 60~\% par
la CSG et la CRDS}{fig:ircsg_ocde}

\figbloc{ircsg_n05_quotient}{Le quotient familial ne vaut rien sous 45\,000~€ de revenu, et jusqu'à 7\,200~€ par an
au-delà de 114\,000~€}{fig:ircsg_quotient}

\figbloc{ircsg_n07_naive_marges}{La fusion naïve garde des marches, à d'autres endroits}{fig:ircsg_naive_marges}

Appliqué à chaque adulte sur l'assiette large, le barème de l'IR rapporte 205~Md€ et le taux proportionnel tombe à
1,64~\%. Le dixième le plus aisé paie 2\,950~€ de plus par foyer (12,5~Md€), le 1~\% le plus aisé 2,8~Md€ ; un couple
avec deux enfants à un seul revenu de 50\,000~€, 2\,360~€ de plus. Le gain est de 2,6~Md€ (1,1 à 5,6), presque
entièrement sur la participation. Le passage de l'IR actuel (86,6~Md€) à ce barème se décompose ainsi : règles actuelles
sans niche, 115,1~Md€ ; assiette large, $+$47,5 ; sans quotient familial, $+$17,4 ; sans décote, $+$7,0 ; sans quotient
conjugal, $+$18,0. L'écart de 28,5~Md€ entre 115,1 et 86,6 recouvre les niches (au moins 24~Md€) et des revenus exonérés
présents dans le revenu fiscal ; il n'est pas bouclé.

\subsection{Le choix de la proposition}

La recherche porte sur la franchise d'activité (0 à 8\,000~€ par adulte), le taux d'entrée (0 à 16~\%), le taux à
20\,000~€ (8 à 26~\%) et le seuil intermédiaire (25\,000 à 90\,000~€). Le taux au seuil intermédiaire est calé sur la
recette, le seuil du taux de 50~\% sur une charge inchangée du 1~\% le plus aisé (à 0,05~Md€ près). Sur 10\,332 barèmes
essayés, 3\,743 tiennent la recette et la charge du 1~\%, et 1\,119 laissent aussi chaque décile à moins de 500~€ de sa
charge actuelle (figure~\ref{fig:ircsg_deciles}). La règle de choix, le gain sur la participation le plus élevé (à
0,02~Md€ près) puis le gain total, retient une franchise de 8\,000~€, à la borne de l'intervalle, et un taux d'entrée de
15~\%. Le barème le plus efficace a la même franchise et un taux d'entrée de 16~\% : il gagne autant (4,69~Md€ contre
4,67), mais 2,1~Md€ seulement sur la participation (tableau~\ref{tab:ircsg_design}, figure~\ref{fig:ircsg_designs}). Le
barème retenu passe par 15~\% à 20\,000~€, 28,8~\% à 30\,000~€, 30~\% à 45\,000~€, 36~\% à 100\,000~€, 41,5~\% à
150\,000~€ et 50~\% à 231\,000~€.

\figbloc{ircsg_n10_deciles}{Avec la proposition, aucun décile ne bouge de plus de 500~€}{fig:ircsg_deciles}

\begin{table}[!htbp]\centering\scriptsize\setlength{\tabcolsep}{3pt}
\caption{Franchise d'activité contre taux d'entrée : le meilleur barème pour chaque combinaison}\label{tab:ircsg_design}
\resizebox{\linewidth}{!}{%
\begin{tabular}{@{}rrrrrrrl@{}}
\toprule
Franchise & Taux & Gain & dont & Retraités & \multicolumn{2}{c}{Premier salaire (\%)} & Barème \\ \cmidrule(lr){6-7}
(€) & d'entrée & Md€ & participation & Md€ & mi-temps & Smic & \\
\midrule
0 & 0~\% & 3,26 & 0,88 & +6,9 & 4,6 & 9,2 & 22~\% à 20~k€, 22,5~\% à 30~k€, 50~\% à 212~k€ \\
0 & 2~\% & 3,49 & 0,78 & +6,7 & 5,7 & 9,7 & 21~\% à 20~k€, 21,2~\% à 30~k€, 50~\% à 194~k€ \\
0 & 4~\% & 3,69 & 0,82 & +6,6 & 6,7 & 9,8 & 19~\% à 20~k€, 24,4~\% à 45~k€, 50~\% à 201~k€ \\
0 & 6~\% & 3,89 & 0,69 & +6,4 & 7,8 & 10,3 & 18~\% à 20~k€, 21,7~\% à 37,5~k€, 50~\% à 184~k€ \\
0 & 8~\% & 4,07 & 0,76 & +6,2 & 8,8 & 10,4 & 16~\% à 20~k€, 24,5~\% à 45~k€, 50~\% à 194~k€ \\
0 & 10~\% & 4,20 & 0,96 & +6,2 & 9,5 & 10,1 & 13~\% à 20~k€, 25,2~\% à 37,5~k€, 50~\% à 219~k€ \\
2\,000 & 0~\% & 3,26 & 1,22 & +10,6 & 3,1 & 8,0 & 24~\% à 20~k€, 28,8~\% à 90~k€, 50~\% à 176~k€ \\
2\,000 & 2~\% & 3,51 & 1,28 & +10,3 & 4,0 & 8,3 & 22~\% à 20~k€, 26,3~\% à 45~k€, 50~\% à 231~k€ \\
2\,000 & 4~\% & 3,74 & 1,16 & +10,1 & 5,0 & 8,9 & 21~\% à 20~k€, 32,3~\% à 90~k€, 50~\% à 212~k€ \\
2\,000 & 6~\% & 3,96 & 1,20 & +9,9 & 5,9 & 9,1 & 19~\% à 20~k€, 24,6~\% à 37,5~k€, 50~\% à 212~k€ \\
2\,000 & 8~\% & 4,17 & 1,23 & +9,8 & 6,8 & 9,4 & 17~\% à 20~k€, 25,3~\% à 37,5~k€, 50~\% à 219~k€ \\
2\,000 & 10~\% & 4,36 & 1,26 & +9,6 & 7,6 & 9,6 & 15~\% à 20~k€, 26,0~\% à 37,5~k€, 50~\% à 227~k€ \\
2\,000 & 12~\% & 4,56 & 1,13 & +9,3 & 8,6 & 10,2 & 14~\% à 20~k€, 24,7~\% à 37,5~k€, 50~\% à 205~k€ \\
4\,000 & 0~\% & 3,30 & 1,62 & +14,4 & 1,8 & 6,8 & 26~\% à 20~k€, 28,8~\% à 90~k€, 50~\% à 182~k€ \\
4\,000 & 2~\% & 3,52 & 1,50 & +14,2 & 2,6 & 7,4 & 25~\% à 20~k€, 27,4~\% à 90~k€, 50~\% à 156~k€ \\
4\,000 & 4~\% & 3,78 & 1,54 & +13,9 & 3,4 & 7,7 & 23~\% à 20~k€, 32,0~\% à 90~k€, 50~\% à 212~k€ \\
4\,000 & 6~\% & 4,01 & 1,67 & +13,8 & 4,0 & 7,8 & 20~\% à 20~k€, 25,9~\% à 30~k€, 50~\% à 247~k€ \\
4\,000 & 8~\% & 4,23 & 1,57 & +13,4 & 4,8 & 8,4 & 19~\% à 20~k€, 26,4~\% à 37,5~k€, 50~\% à 231~k€ \\
4\,000 & 10~\% & 4,44 & 1,58 & +13,2 & 5,6 & 8,8 & 17~\% à 20~k€, 27,0~\% à 37,5~k€, 50~\% à 239~k€ \\
4\,000 & 12~\% & 4,64 & 1,45 & +12,9 & 6,4 & 9,4 & 16~\% à 20~k€, 25,7~\% à 37,5~k€, 50~\% à 215~k€ \\
4\,000 & 14~\% & 4,83 & 1,45 & +12,7 & 7,2 & 9,8 & 14~\% à 20~k€, 26,4~\% à 37,5~k€, 50~\% à 223~k€ \\
6\,000 & 6~\% & 4,05 & 1,93 & +17,6 & 2,4 & 6,9 & 23~\% à 20~k€, 25,2~\% à 25~k€, 50~\% à 239~k€ \\
6\,000 & 8~\% & 4,28 & 1,92 & +17,3 & 2,9 & 7,4 & 21~\% à 20~k€, 26,1~\% à 30~k€, 50~\% à 239~k€ \\
6\,000 & 10~\% & 4,50 & 1,91 & +17,1 & 3,5 & 7,8 & 19~\% à 20~k€, 26,2~\% à 30~k€, 50~\% à 239~k€ \\
6\,000 & 12~\% & 4,71 & 2,00 & +17,0 & 4,0 & 8,0 & 16~\% à 20~k€, 27,8~\% à 30~k€, 50~\% à 274~k€ \\
6\,000 & 14~\% & 4,89 & 1,87 & +16,6 & 4,5 & 8,7 & 15~\% à 20~k€, 26,4~\% à 30~k€, 50~\% à 239~k€ \\
7\,000 & 6~\% & 4,08 & 2,03 & +19,6 & 1,6 & 6,5 & 25~\% à 20~k€, 31,3~\% à 90~k€, 50~\% à 200~k€ \\
7\,000 & 8~\% & 4,28 & 2,01 & +19,3 & 2,0 & 6,9 & 23~\% à 20~k€, 25,5~\% à 30~k€, 50~\% à 227~k€ \\
7\,000 & 10~\% & 4,53 & 2,09 & +19,1 & 2,4 & 7,2 & 20~\% à 20~k€, 27,1~\% à 30~k€, 50~\% à 256~k€ \\
7\,000 & 12~\% & 4,73 & 2,07 & +18,9 & 2,8 & 7,7 & 18~\% à 20~k€, 27,2~\% à 30~k€, 50~\% à 256~k€ \\
7\,000 & 14~\% & 4,89 & 2,04 & +18,6 & 3,2 & 8,1 & 16~\% à 20~k€, 27,2~\% à 30~k€, 50~\% à 247~k€ \\
7\,000 & 16~\% & 5,02 & 1,80 & +18,2 & 3,7 & 8,9 & 16~\% à 20~k€, 24,3~\% à 30~k€, 50~\% à 201~k€ \\
8\,000 & 8~\% & 4,31 & 2,21 & +21,4 & 1,2 & 6,3 & 24~\% à 20~k€, 25,8~\% à 25~k€, 50~\% à 239~k€ \\
8\,000 & 10~\% & 4,51 & 2,19 & +21,1 & 1,4 & 6,8 & 22~\% à 20~k€, 26,5~\% à 30~k€, 50~\% à 239~k€ \\
8\,000 & 12~\% & 4,72 & 2,23 & +21,0 & 1,6 & 7,1 & 19~\% à 20~k€, 26,4~\% à 25~k€, 50~\% à 247~k€ \\
8\,000 & 14~\% & 4,91 & 2,27 & +20,8 & 1,9 & 7,4 & 16~\% à 20~k€, 27,4~\% à 25~k€, 50~\% à 265~k€ \\
8\,000 & 16~\% & 5,04 & 2,06 & +20,3 & 2,2 & 8,2 & 16~\% à 20~k€, 24,7~\% à 25~k€, 50~\% à 215~k€ \\
\bottomrule
\end{tabular}}
\par\vspace{3pt}\parbox{\linewidth}{\raggedright\scriptsize\color{encre}
Pour chaque franchise d'activité (par adulte) et chaque taux d'entrée (taux sur le premier euro imposé après la franchise), le barème le plus efficace parmi ceux qui respectent la recette, les déciles à $\pm$500~€ et la neutralité du 1~\%. Sont aussi fixées : l'exonération des petites pensions, la franchise de 6\,000~€ par enfant, la déduction de l'épargne retraite. Premier salaire : taux de prélèvement IR-CSG d'une personne seule qui prend un emploi au Smic, à mi-temps ou à temps plein (aujourd'hui : 9,5~\%). Taux d'entrée pairs seulement ; 10332 barèmes essayés, 1047 admissibles. La proposition retient la combinaison qui maximise la participation (franchise de 8\,000~€, taux d'entrée de 14~\%).}
\end{table}


\figbloc{ircsg_n17_designs}{Plus la franchise est forte, plus la participation progresse, et le gain total avec
elle}{fig:ircsg_designs}

\subsection{Les variantes}

Le tableau~\ref{tab:ircsg_dec} donne la variation par décile de toutes les variantes, le tableau~\ref{tab:ircsg_arbcas}
et la figure~\ref{fig:ircsg_cas} leurs effets sur des ménages types, la figure~\ref{fig:ircsg_arbitrages} le gain de
chaque arbitrage ; le tableau~\ref{tab:ircsg_arch} rappelle les architectures de la méthode v2, où chaque dispositif est
ajouté seul au barème neutre et financé par des points de taux pour tous. Un chèque par enfant égal au coût actuel du
quotient familial (765~€) libérerait 0,7 point de taux pour tous.

\begin{table}[!htbp]\centering\scriptsize\setlength{\tabcolsep}{3pt}
\caption{Variation de l'impôt par décile de revenu fiscal, toutes variantes (€ par foyer et par an)}\label{tab:ircsg_dec}
\resizebox{\linewidth}{!}{%
\begin{tabular}{@{}lrrrrrrrrrrrr@{}}
\toprule
& D1 & D2 & D3 & D4 & D5 & D6 & D7 & D8 & D9 & D10 & 1~\% (Md€) & Retraités (Md€) \\
\midrule
Réforme naïve (quotient, 3,99~\%) & $-$11 & $-$37 & +162 & +105 & $-$543 & $-$551 & $-$570 & +95 & $-$524 & +1\,841 & +2,6 & +8,2 \\
Réforme naïve, variante individuelle (1,64~\%) & $-$57 & $-$217 & $-$171 & $-$256 & $-$906 & $-$938 & $-$794 & +36 & +277 & +2\,953 & +2,8 & +2,7 \\
\textbf{Proposition} & $-$84 & $-$279 & +80 & +251 & $-$215 & $-$155 & $-$320 & +354 & $-$71 & +436 & 0,0 & +21,5 \\
Variante : franchise de 4\,000~€ & $-$85 & $-$149 & +147 & +225 & $-$286 & $-$129 & $-$237 & +304 & $-$158 & +367 & 0,0 & +14,3 \\
Variante : sans franchise & +2 & +19 & +243 & +318 & $-$225 & $-$202 & $-$414 & +113 & $-$259 & +399 & 0,0 & +6,9 \\
Quotient conjugal et familial & $-$84 & $-$285 & $-$18 & +119 & $-$276 & $-$137 & $-$149 & +468 & $-$99 & +460 & 0,0 & +18,9 \\
Abattement de 10~\% sur les pensions & $-$83 & $-$277 & +42 & +273 & $-$200 & $-$138 & $-$287 & +314 & $-$94 & +448 & 0,0 & +14,8 \\
Chèque par enfant (même coût) & $-$168 & $-$354 & +16 & +197 & $-$226 & $-$99 & $-$227 & +392 & +11 & +455 & 0,0 & +21,6 \\
Le 1~\% peut payer moins & $-$84 & $-$278 & +105 & +335 & $-$86 & $-$11 & $-$197 & +493 & +162 & $-$423 & $-$3,8 & +22,5 \\
Barème neutre (méthode v2) & 0 & +83 & +265 & +165 & $-$365 & $-$221 & $-$324 & +130 & $-$27 & +290 & $-$0,5 & +5,0 \\
Ancienne proposition (1~\% en baisse) & $-$85 & $-$145 & +190 & +342 & $-$150 & $-$12 & $-$121 & +450 & 0 & $-$453 & $-$2,5 & +13,9 \\
\bottomrule
\end{tabular}}
\par\vspace{3pt}\parbox{\linewidth}{\raggedright\scriptsize\color{encre}
Déciles de foyers fiscaux classés par revenu fiscal de référence (IRCOM, revenus 2024 ; montants de 2025). Moyennes de foyers représentatifs : elles cachent les écarts à l'intérieur de chaque décile (sous-partie~\ref{sec:ircsg}.4).}
\end{table}


\begin{table}[!htbp]\centering\scriptsize\setlength{\tabcolsep}{3pt}
\caption{Les arbitrages pour des ménages types (variation de l'IR + CSG + CRDS, € par an)}\label{tab:ircsg_arbcas}
\resizebox{\linewidth}{!}{%
\begin{tabular}{@{}lrrrrr@{}}
\toprule
& Proposition & Quotient & Abatt. 10~\% & Chèque & Franchise 4\,000~€ \\
\midrule
Jeune salarié au Smic & $-$447 & $-$635 & $-$259 & $-$359 & $-$243 \\
Cadre seul, 100 000 € brut & $-$4\,015 & $-$2\,733 & $-$3\,682 & $-$4\,962 & $-$3\,894 \\
Parent isolé, deux enfants, au Smic & $-$2\,027 & $-$688 & $-$2\,027 & $-$2\,740 & $-$1\,683 \\
Couple sans enfant à un seul revenu, 80 000 € & +3\,498 & +65 & +3\,814 & +2\,873 & +3\,677 \\
Couple avec deux enfants à un seul revenu, 50 000 € & $-$862 & $-$77 & $-$539 & +92 & $-$518 \\
Couple avec trois enfants, 60 000 € brut (60/40) & $-$2\,922 & $-$993 & $-$2\,591 & $-$3\,078 & $-$2\,452 \\
Couple avec trois enfants à un seul revenu, 60 000 € & $-$1\,003 & +102 & $-$682 & +315 & $-$670 \\
Retraité seul, pension de 16 000 € & +1\,005 & +877 & +871 & +1\,041 & +840 \\
Retraité seul, pension de 20 000 € & +1\,353 & +1\,109 & +1\,262 & +1\,521 & +940 \\
Retraité seul, pension de 25 000 € & +1\,815 & +1\,577 & +1\,519 & +2\,118 & +1\,092 \\
Couple de retraités, 45 000 € de pensions (60/40) & +1\,330 & +196 & +809 & +1\,667 & +309 \\
\bottomrule
\end{tabular}}
\par\vspace{3pt}\parbox{\linewidth}{\raggedright\scriptsize\color{encre}
Mêmes cas-types que le tableau~\ref{tab:ircsg_cas}. Quotient : imposition commune des couples, quotient conjugal et familial aux plafonds actuels à la place de la franchise par enfant ; abattement de 10~\% sur les pensions en plus de l'exonération des petites pensions ; chèque de 1\,191~€ par enfant, de même coût que la franchise de 6\,000~€ par enfant ; franchise d'activité de 4\,000~€ par adulte au lieu de 8\,000~€. Chaque variante a son propre barème, recherché sous les mêmes contraintes.}
\end{table}


\figbloc{ircsg_n15_cas_types}{Familles et salariés gagnent, retraités et couples à un seul revenu perdent}{fig:ircsg_cas}

\figbloc{ircsg_n14_arbitrages}{Chaque arbitrage déplace le gain de moins d'un milliard, mais la charge des retraités de
plusieurs}{fig:ircsg_arbitrages}

\figbloc{ircsg_n18_variantes_marges}{Étaler la sortie de l'exonération abaisse le taux des petites pensions et le
relève un peu plus haut}{fig:ircsg_variantes_marges}

\begin{table}[!htbp]\centering\scriptsize\setlength{\tabcolsep}{3pt}
\caption{Architectures comparées, à recette égale (220~Md€)}\label{tab:ircsg_arch}
\resizebox{\linewidth}{!}{%
\begin{tabular}{@{}lrrrrrrr@{}}
\toprule
& Points & Gain d'efficacité & Part du & Perdants & Retraités & Couple, & 10~\% \\
& pour tous & Md€ par an & gain & $>$ 500~€ & Md€ & 2 enfants & du haut \\
& & (fourchette) & atteignable & & & € par foyer & Md€ \\
\midrule
Barème neutre & 0,0 & 4,0 (1,9 à 9,2) & 52~\% & 25~\% & +7,0 & +570 & +2,5 \\
\multicolumn{8}{@{}l}{\emph{Un dispositif ajouté au barème neutre}} \\
\quad Franchise d'activité 4~000~€ & +1,8 & 4,6 (2,1 à 10,2) & 59~\% & 30~\% & +14,0 & +217 & +3,6 \\
\quad Franchise d'activité 8~000~€ & +3,6 & 5,0 (2,2 à 11,0) & 64~\% & 33~\% & +21,4 & $-$172 & +4,7 \\
\quad Franchise de 8~000~€ sur tous les revenus & +5,8 & 3,3 (1,4 à 7,4) & 43~\% & 21~\% & +5,4 & +775 & +9,6 \\
\quad Franchise de 6~000~€ par enfant & +1,5 & 3,6 (1,7 à 8,4) & 47~\% & 27~\% & +13,0 & $-$983 & +0,6 \\
\quad Chèque de 742~€ par enfant & +0,9 & 3,1 (1,5 à 7,3) & 40~\% & 26~\% & +10,7 & $-$361 & +3,1 \\
\quad Quotient conjugal et familial & +2,0 & 3,9 (1,9 à 9,0) & 50~\% & 28~\% & +12,1 & $-$965 & $-$3,0 \\
\quad Abattement de 10~\% sur les pensions & +0,6 & 3,4 (1,6 à 8,0) & 44~\% & 24~\% & +0,9 & +923 & +2,7 \\
\quad Petites pensions exonérées & +0,4 & 3,7 (1,7 à 8,5) & 47~\% & 16~\% & +3,2 & +786 & +4,0 \\
\midrule
\textbf{Proposition}, franchise d'activité 4~000~€ & --- & 4,9 (2,3 à 11,1) & 63~\% & 29~\% & +15,2 & $-$1\,195 & $-$1,7 \\
Proposition, franchise d'activité 8~000~€ & --- & 5,2 (2,4 à 11,8) & 67~\% & 26~\% & +22,8 & $-$1\,629 & $-$1,5 \\
\midrule
Optimum d'efficacité, sans franchise & --- & 7,4 (3,7 à 17,4) & 95~\% & 58~\% & +13,7 & $-$519 & $-$25,7 \\
Gain atteignable (optimum, franchise 4~000~€) & --- & 7,8 (3,8 à 18,2) & 100~\% & 49~\% & +20,3 & $-$869 & $-$25,1 \\
\bottomrule
\end{tabular}}
\par\vspace{3pt}\parbox{\linewidth}{\raggedright\scriptsize\color{encre}
Gain d'efficacité : baisse de la perte sèche sur l'activité réelle par rapport à aujourd'hui, hypothèses centrales (basses et hautes entre parenthèses) ; la réaction du revenu déclaré n'y est pas comptée. Points pour tous : hausse des taux du barème neutre qui finance le dispositif. Proposition : franchise d'activité, 6~000~€ par enfant, pensions exonérées jusqu'à 12~500~€ puis de moins en moins jusqu'à 25~000~€ ; barème le plus efficace parmi ceux qui ne changent la charge moyenne d'aucun décile de plus de 500~€ et n'alourdissent pas celle du 1~\% le plus aisé. Retraités, couples, 10~\% du haut : variation de leur impôt (positif : ils paient plus).}
\end{table}


\subsection{Sensibilités, haut de la distribution et transition}

\paragraph{Ce qui fait varier le gain.} La perception de la CSG domine toutes les autres sensibilités
(tableau~\ref{tab:ircsg_sens}) : la fusion n'est efficace que si la CSG d'aujourd'hui fait déjà réagir les salariés.

\begin{table}[!htbp]\centering\scriptsize\setlength{\tabcolsep}{4pt}
\caption{Ce qui fait varier le gain de la proposition et la neutralité du 1~\%}\label{tab:ircsg_sens}
\resizebox{\linewidth}{!}{%
\begin{tabular}{@{}lrl@{}}
\toprule
Hypothèse & Gain (Md€) & Autre effet \\
\midrule
Modèle, hypothèses centrales & 4,91 & participation 2,27~Md€ \\
Élasticités basses / hautes & 2,25 / 10,94 & \\
Élasticités de long terme (intensif 0,10 à 0,25 ; participation 0,10) & 7,30 & participation 0,91~Md€ \\
CSG perçue à moitié aujourd'hui & 1,29 & intensif $-$0,98~Md€ \\
CSG non perçue aujourd'hui & $-$2,09 & intensif $-$4,37~Md€ \\
Régime permanent : impôt futur sur les pensions vu à moitié par les actifs & 4,10 & \\
Régime permanent : vu en entier & 3,25 & coin $+$1,40~point \\
Répartition des revenus au PFU logés en haut (convention « haut ») & 4,5 & 1~\% $-$2,8~Md€ ; neutre si 50~\% à 110~k€ \\
Répartition des revenus au PFU au prorata & 5,3 & 1~\% +3,7~Md€ ; neutre si 50~\% à 858~k€ \\
\bottomrule
\end{tabular}}
\par\vspace{3pt}\parbox{\linewidth}{\raggedright\scriptsize\color{encre}
Gain : baisse de la perte sèche sur l'activité (jeu central sauf mention). CSG non perçue : le revenu d'activité ne réagit aujourd'hui qu'à l'IR (comme aux cotisations chez Lehmann, Marical et Rioux, 2013) et réagit après la fusion à tout l'impôt fusionné ; la marge de participation est gardée. Régime permanent : la hausse de l'impôt sur les pensions réduit la valeur des droits acquis par les cotisations contributives (27,4~\% du salaire brut) ; borne haute, le lien entre cotisation et pension étant faible au bas de l'échelle et nul au-delà de 8 plafonds. Convention centrale de répartition des revenus au PFU : « mixte » (controle\_hauts\_purs.py).}
\end{table}


\paragraph{Le haut de la distribution.} La DGFiP ne publie pas, par tranche, la part des revenus imposés au PFU. Dans la
convention centrale (« mixte »), la charge du 1~\% le plus aisé ne bouge pas : la baisse des hauts salaires y est
compensée par une hausse de 3 à 16 points pour les revenus du barème au-delà de 1~M€, auxquels cette convention prête un
taux actuel de 33 à 39~\%. Dans la convention « haut », plus proche du barème légal (44 à 51~\%) et que soutient un
calage sur l'IR observé (\texttt{revue/hypotheses\_ircsg.md}, H9), le 1~\% paierait 3,0~Md€ de moins, les foyers de 1 à
9~M€ de 2 à 8 points de moins ; garder sa charge demanderait le taux de 50~\% dès 91\,000~€ environ, mais le dixième le
plus aisé paierait alors 1\,300~€ de plus en moyenne. Au prorata, le 1~\% paierait 3,8~Md€ de plus
(\texttt{sorties/controle\_hauts\_purs.txt}). La disparition de la CEHR (3 à 4 points au très haut) sur les revenus du
barème, qui rend 0,3 à 0,7~Md€ au 1~\%, n'est pas comptée ici mais dans le bilan du haut (section~\ref{sec:ens_haut}).

\paragraph{La transition.} Le tableau~\ref{tab:ircsg_trans} donne la hausse annuelle des principaux perdants selon la
durée de la convergence.

\begin{table}[!htbp]\centering\scriptsize\setlength{\tabcolsep}{4pt}
\caption{Hausse annuelle des principaux perdants selon la durée de la transition (€ par an)}\label{tab:ircsg_trans}
\begin{tabular}{@{}lrrrr@{}}
\toprule
& Hausse finale & \multicolumn{3}{c}{Hausse chaque année, transition sur} \\ \cmidrule(l){3-5}
& & 3 ans & 5 ans & 10 ans \\
\midrule
Retraités, médiane & +1\,392 & 464 & 278 & 139 \\
Retraités, 90\textsuperscript{e} centile & +3\,423 & 1\,141 & 685 & 342 \\
Couples de retraités, 90\textsuperscript{e} centile & +5\,659 & 1\,886 & 1\,132 & 566 \\
Retraités, 99\textsuperscript{e} centile & +7\,423 & 2\,474 & 1\,485 & 742 \\
Tous les foyers, 90\textsuperscript{e} centile & +2\,245 & 748 & 449 & 224 \\
Retraité seul, pension de 16 000 € & +1\,050 & 350 & 210 & 105 \\
Retraité seul, pension de 20 000 € & +1\,390 & 463 & 278 & 139 \\
Retraité seul, pension de 25 000 € & +1\,679 & 560 & 336 & 168 \\
Couple de retraités, 45 000 € de pensions (60/40) & +1\,834 & 611 & 367 & 183 \\
Couple sans enfant à un seul revenu, 80 000 € & +3\,715 & 1\,238 & 743 & 372 \\
\bottomrule
\end{tabular}
\par\vspace{3pt}\parbox{\linewidth}{\raggedright\scriptsize\color{encre}
Convergence linéaire : chaque année, chaque foyer paie une part croissante du nouvel impôt et décroissante de l'ancien, ce qui tient la recette chaque année. Centiles : foyers représentatifs du modèle (la dispersion réelle des hausses est plus forte ; le 99\textsuperscript{e} centile en est un minorant). Cas-types : tableau~\ref{tab:ircsg_cas}.}
\end{table}


\paragraph{La valeur actualisée.} La hausse annuelle des retraités est multipliée par un facteur d'annuité : table de
Gompertz ($\mu(x) = A\,e^{0,10x}$) calée sur une espérance de vie à 65 ans de 21,8 ans (Insee 2025 provisoire) ;
cohortes de même taille dès 62 ans (âge moyen 76 ans, hypothèse) ; taux réel de 1 à 3~\%. Le facteur moyen va de 10,8 à
12,9 (8,6 à 10,7 avec une montée en charge sur cinq ans). Au taux de 2~\%, un euro de hausse annuelle vaut 17~€ à 65
ans, 12~€ à 75 ans et 7~€ à 85 ans (\texttt{sorties/transition\_ircsg.txt}).

\paragraph{Le gain permanent.} En régime permanent, la hausse de l'impôt sur les pensions (5,1 points de taux moyen)
réduit la valeur des droits à pension qu'ouvre chaque euro de salaire, soit environ 27 centimes par euro brut sous le
plafond de la Sécurité sociale (hypothèse) : elle alourdit d'environ 1,4 point le prélèvement implicite sur le salaire
des actifs. Selon que ceux-ci le perçoivent à moitié ou en entier, le gain permanent est de 3,8 ou de 3,0~Md€ au lieu de
4,7. C'est une borne : le lien entre un euro de salaire et la pension est faible au bas de l'échelle (minimum
contributif) et nul au-delà de huit fois le plafond.

\subsection{Hypothèses fortes et chiffres non vérifiés}

\begin{table}[!htbp]\centering\scriptsize\setlength{\tabcolsep}{3pt}
\caption{Hypothèses fortes et chiffres non vérifiés}\label{tab:ircsg_hyp}
\begin{tabular}{@{}p{0.43\linewidth}p{0.2\linewidth}p{0.33\linewidth}@{}}
\toprule
Hypothèse & Où elle joue & Ce qu'elle change si elle est fausse \\
\midrule
Perception de la CSG : la CSG et la CRDS d'aujourd'hui font réagir le revenu d'activité autant que l'IR. Aucune étude ne le mesure. & Tout le gain sur la marge intensive & Décisif. CSG non perçue : gain $-$2,5~Md€ au lieu de 4,7 ; perçue à moitié : 1,0. \\[2pt]
Élasticités : revenu réel 0,05 (moitié du bas) et 0,10 (au-dessus) ; participation 0,25 sous 2 Smic ; ce sont des choix, pas des estimations françaises. & Tous les gains & Fourchette 2,1 à 10,3~Md€ ; long terme 6,7~Md€. \\[2pt]
Les pensions, le chômage et l'épargne ne réagissent pas (âge de départ, surcote, cumul emploi-retraite, sortie en capital). & Mécanisme du financement par les pensions & Non chiffré ; sens ambigu (une pension nette plus faible pousse à partir plus tard ; un taux marginal élevé dans la zone de sortie de l'exonération pousse en sens inverse). \\[2pt]
Régime permanent : les actifs ne voient pas que leurs futures pensions seront plus taxées. & Part « permanente » du gain & Vu en entier : gain permanent 3,0~Md€ au lieu de 4,7. \\[2pt]
Répartition des revenus au PFU entre tranches (non publiée) ; convention « mixte ». & Charge du 1~\%, forme du haut du barème & Convention « haut » (que soutient un calage sur l'IR observé) : 1~\% $-$3,0~Md€ ; charge inchangée si 50~\% dès 91~k€. \\[2pt]
Foyers purs : pas de foyer mixte (couple dont un seul est retraité, retraité bailleur) ; loyers tous chez les actifs. & Coût de la franchise, charge des retraités & Franchise sous-estimée et hausse des retraités surestimée, jusqu'à 3~Md€ environ (borne grossière). \\[2pt]
Partage des revenus dans les couples (Insee Première 1492, quatre classes ; le même pour les retraités). & Coût de l'impôt individuel, charge des couples de retraités & Couples de retraités à 62/38 : charge des retraités plus faible de 2~Md€ environ (registre, H11, ancienne proposition). \\[2pt]
Enfants : un par parent isolé, deux par couple « avec enfants », aucun chez les retraités et les concubins. & Coût de la franchise par enfant, familles nombreuses & Familles de trois enfants absentes des agrégats (cas-types seulement). \\[2pt]
Charge actuelle par type de foyer répartie selon des cas-types sans niche. & Qui paie entre types de foyers & Charge actuelle des retraités (50,2~Md€) non observable ; niches réparties au prorata. \\[2pt]
Base 2025 : l'IR de 2025 est l'impôt encaissé dans l'année (comptes nationaux), dont la hausse (+7,9~\%) tient en partie au jeu des acomptes et des soldes du prélèvement à la source ; chaque agrégat de 2024 est porté en 2025 avec la même croissance dans toutes les tranches. & Recette à remplacer, charge actuelle par tranche & Si cette hausse est passagère, la recette est surestimée d'environ 4,5~Md€ (IR porté en 2025 comme les revenus), soit environ 0,4 point de taux. \\[2pt]
Épargne retraite : déduction de 10,6~Md€ (cases 2042) supposée proportionnelle aux revenus d'activité de la moitié haute. & Barème, haut de la distribution & Coût de la déduction 3,0~Md€ (0,29~point) ; répartition non vérifiée. \\[2pt]
Cas-types : législation des revenus 2025 d'OpenFisca, Smic de 21\,622~€ (sources secondaires) ; retraités en régime stable (taux de CSG le plus élevé dont le seuil est dépassé par le revenu fiscal calculé à ce taux). & Tous les cas-types, courbes de taux marginaux & Dans une bande de quelques centaines d'euros autour de chaque seuil de CSG, aucun taux n'est cohérent avec son propre revenu fiscal ; le taux inférieur est retenu, ce qui déplace le saut de moins de 1\,000~€. \\[2pt]
Valeur actualisée : table de Gompertz calée sur l'espérance de vie à 65 ans (21,8 ans) ; cohortes de retraités de même taille. & Valeur actualisée de la hausse des retraités & Ordre de grandeur ; 1 à 3~\% réel donne la fourchette. \\[2pt]
Chiffres non vérifiés sur une source primaire : CSG des pensions de 2024 (CCSS) ; croissance de la CSG par assiette en 2025 (régimes de base et FSV, appliquée à tous les affectataires) ; paramètres 2025 d'OpenFisca (seuils de CSG des pensions, abattement des plus de 65 ans) ; Smic 2025 ; part des familles au plafond du quotient familial ; partage des niches par tranche ; l'IRCOM inclut l'IR au PFU et exclut la CEHR. & Calage, cas-types, récit & Registre revue/hypotheses\_ircsg.md, section 4. \\[2pt]
\bottomrule
\end{tabular}
\par\vspace{3pt}\parbox{\linewidth}{\raggedright\scriptsize\color{encre}
Source : registre revue/hypotheses\_ircsg.md (hypothèses H1 à H20), chiffres recalculés pour la proposition par transition\_ircsg.py, controle\_hauts\_purs.py et proposition\_ircsg.py.}
\end{table}


\subsection{Programmes}

{\sloppy Dans \texttt{synthese/}, depuis \texttt{fiscalite/} : \texttt{python3.13 openfisca\_cas\_ircsg.py}, puis
\texttt{situation\_ircsg.py}, \texttt{architectures\_ircsg.py}, \texttt{naive\_ircsg.py},
\texttt{proposition\_ircsg.py}, \texttt{arbitrages\_prop\_ircsg.py}, \texttt{transition\_ircsg.py},
\texttt{controle\_hauts\_purs.py}, \texttt{tables\_ircsg.py} et \texttt{figures\_ircsg\_n.py} (deux fois, la seconde
avec \texttt{FIG\_NU=1}). Modèle : \texttt{modele\_purs.py} (\texttt{BASE\_IRCSG=2024} rend l'ancienne base) ;
outils : \texttt{outils\_ircsg.py} ; vérifications : \texttt{revue/calculs\_verif/ircsg/}.\par}
}{}
\section{Immobilier : calibrage, hypothèses et variantes}\label{ann:foncier}
\IfFileExists{parties/annexe_foncier.tex}{% Annexe de la partie Immobilier (refonte du 16/11/2026, coupée le 17/11/2026 selon synthese/CONSIGNES_COUPE.md).
% Chiffres : sorties/calculs_foncier.txt et .json (synthese/calculs_foncier.py) ; ../foncier/sorties/simulation.txt
% (foncier/simulation.py) ; registre : ../revue/hypotheses_foncier.md ; figures : ../foncier/figures_note.py. Base 2025
% (décision 5 du 17/11/2026) : données 2025 par ../foncier/donnees_2025.py. Figures n04, n05 et n10 venues du corps.
% Le tableau « base 2024 / base 2025 » de la version du 17/11 (avant coupe) est retiré : chiffres de la base 2024 par
% BASE_FONCIER=2024 (voir calculs_foncier.py) et dans l'historique git (commit e7232a5).

\subsection{Champ, données et calibrage}

\textbf{Champ.} Les 30,5~millions de logements des particuliers de France métropolitaine hors Alsace-Moselle (dont DVF
ne publie pas les ventes), environ 94~\% de la métropole : parc du recensement de 2022 moins les logements HLM loués,
résidences secondaires, logements vacants et environ un million de logements d'autres personnes morales compris. L'unité
de calcul est la cellule « commune × maisons » ou « commune × appartements », dont tous les logements sont supposés
identiques ; la taxe de chaque commune est répartie entre statuts d'occupation au prorata du nombre de logements
(figure~\ref{fig:fon_qui}).

\textbf{Base 2025.} Taxe foncière bâtie par commune (REI 2025), ventes de logements anciens (DVF 2025), comptes de
patrimoine et PIB de 2025 (Insee, comptes de la Nation publiés en 2026), droits et taxe foncière sur les logements (Compte
du logement 2025), droits de toutes les collectivités (OFGL, comptes de gestion 2025). Restent antérieurs le parc
(recensement de 2022), le rapport taxe/valeur de chaque commune (2017) et les prix des logements (ventes de 2023-2024).

\textbf{Valeur.} Prix médian au m² des ventes 2023-2024 de la commune (à défaut, de l'intercommunalité, puis du
département, sous dix ventes) × surface médiane des logements vendus : 6\,762~Md€, 20~\% de moins que les comptes de
patrimoine de l'Insee de fin 2025 ramenés au champ (8\,467~Md€). L'écart n'est pas uniforme : le prix moyen des ventes
dépasse la valeur modélisée de 24~\% dans le dixième de communes le moins cher, de 6 à 10~\% au milieu et de 21~\% dans le
plus cher. Les recettes sont justes, mais les taux en part de la valeur du modèle sont surestimés d'environ un cinquième ;
le corps du texte donne donc les taux en valeurs de l'Insee. Les prix de 2025 sont inférieurs de 1,5~\% à ceux de
2023-2024 (indices Notaires-Insee : Paris $-$2,0~\%, Île-de-France hors Paris $-$3,6~\%, province $-$1,0~\%).

\textbf{Part du sol.} Maisons : prix au m² des terrains à bâtir de la commune (à défaut, celui du département ajusté par
l'écart de prix des maisons, élasticité de 1,54 estimée sur 3\,307~communes) × parcelle médiane plafonnée à 1\,500~m².
Appartements : un moins le rapport entre un coût de reconstruction uniforme (1\,464~€/m², calé pour que le sol fasse
46,2~\% de la valeur, comme dans les comptes de patrimoine des ménages de fin 2025) et le prix au m². Bornes de 10 et
90~\%. Le sol vaut 3\,126~Md€ ; sa part va de 16~\% dans le dixième de logements le moins cher à 68~\% dans le plus cher,
et atteint 85~\% pour l'appartement parisien médian.

\textbf{Taxe foncière actuelle.} Rapport taxe/valeur de la commune en 2017 (André et Meslin, 2023) × valeur × croissance
de la taxe de la commune de 2017 à 2025 (REI), calé sur 19,2~Md€ (2017) × 1,333 (croissance de la taxe foncière bâtie de
la métropole de 2017 à 2025) × part du champ : 24,2~Md€ en 2025, 795~€ par logement, 0,36~\% de la valeur. Les communes
dont la taxe varie de moins de $-$20~\% ou de plus de +30~\% entre 2024 et 2025 (surtout des communes nouvelles,
0,2~\% des logements) gardent leur croissance de 2017 à 2024, prolongée de la croissance nationale de 2025 (+2,8~\%).

\textbf{Droits de mutation.} 5,81~\% (4,50~\% départemental, 1,20~\% communal, frais d'assiette de 2,37~\% de la part
départementale) des 212~Md€ de ventes de logements anciens de 2025 (840\,000 ventes, 2,8~\% du parc) : 12,3~Md€, dont 9,5
départementaux, 2,5 communaux et 0,2 de frais (196~Md€ de ventes et 11,4~Md€ de droits par an en 2023-2024). Rapportés au
parc, 404~€ par logement et par an, 0,18~\% de la valeur. Le relèvement temporaire de la part départementale à 5~\% n'est
pas compté : les droits perçus sur les logements en 2025 (13,6~Md€, France entière ; figure~\ref{fig:fon_cycle})
dépassent de 0,4~Md€ ce qu'aurait donné le taux permanent si le champ en représente la même part qu'en 2023-2024
(93,9~\%), soit un taux moyen effectif d'environ 6,0~\% (calcul de l'auteur, non vérifié).

\textbf{Scénarios.} Tous remplacent la moitié des droits sur les logements (6,2~Md€ sans réaction des ventes) par de la
taxe foncière, à recette constante. Réforme naïve : taxe actuelle relevée de 25,4~\% partout, soit 30,4~Md€ de taxe sur
les logements des particuliers. « N4 » (proposition de la version précédente) : chaque commune garde sa recette, répartie
selon deux taux (sol = 4 × bâti), et un taux national de 0,197~\% sur le sol des logements du modèle remplace les droits.
Proposition : même part communale, droits divisés par deux aussi sur les autres immeubles, taux national sur tous les
terrains privés ; la taxe des logements des particuliers atteint 28,9~Md€.

\figbloc{foncier_n04_qui_paie}{Les propriétaires occupants paient un peu plus de la moitié de la taxe foncière des logements des particuliers}{fig:fon_qui}

\figbloc{foncier_n05_cycle}{Les droits sur les ventes de logements ont perdu un tiers de leur produit entre 2022 et 2024, avant de se redresser en partie en 2025}{fig:fon_cycle}

\subsection{L'assiette du taux national : tous les terrains privés}

Le taux national porte sur le sol sous les bâtiments et les terrains non bâtis des ménages, des sociétés et des
associations, soit 6\,656~Md€ fin 2025, hors terres cultivées (902~Md€) et terrains publics (1\,094~Md€)
(tableau~\ref{tab:fon_insee}) ; le sol sous les bâtiments des ménages en représente 62,1~\%.

\begin{table}[!htbp]\centering\scriptsize
\caption{Valeur des terrains par secteur, fin 2025 (Md€)}\label{tab:fon_insee}
\begin{tabularx}{\linewidth}{@{}>{\raggedright\arraybackslash}X rrrr@{}}
\toprule
Secteur & Sol sous les bâtiments & Terres cultivées & Autres terrains (non bâtis) & Dans l'assiette \\
\midrule
Ménages & 4\,136 & 193 & 452 & 4\,588 \\
Sociétés non financières (dont HLM) & 1\,701 & 710 & 113 & 1\,814 \\
Sociétés financières & 200 & --- & --- & 200 \\
Institutions sans but lucratif & 54 & --- & --- & 54 \\
Administrations publiques & 941 & --- & 153 & exclues \\
Ensemble & 7\,032 & 902 & 717 & 6\,656 \\
\bottomrule
\end{tabularx}
\par\vspace{2pt}{\raggedright\scriptsize Source : Insee, comptes nationaux annuels base 2020, tableau 8.200, millésime
de juin 2026 (\texttt{foncier/donnees/terrains\_secteurs\_2025.csv}), cohérent avec l'Insee Première n° 2135 (novembre
2026 : 4\,190~Md€ de terrains bâtis pour les ménages au sens large). Ce millésime relève de 5~\% la valeur des terrains
bâtis de 2024 (7\,044 au lieu de 6\,729~Md€). Le sol est une valeur résiduelle (valeur de marché moins bâti au coût de
remplacement), sans marge d'incertitude publiée ; les « autres terrains » comprennent des terrains naturels et boisés
valorisés au prix des terres agricoles et des forêts.\par}
\end{table}

\textbf{Méthode.} Le taux national lève $R$ = 7,58~Md€ à l'échelle du champ (6,16~Md€ de droits sur les logements, 1,42
sur les autres immeubles). La part des logements des particuliers, $s = 4\,136/6\,656$, est répartie entre logements au
prorata de la valeur de leur sol dans le modèle : $sR/3\,126 = 0,151~\%$ des valeurs du modèle. En valeurs de l'Insee, le
taux vaut $R/(6\,656 \times 0,946) = 0,120~\%$, si le champ représente la même fraction (94,6~\%) des terrains de chaque
secteur. Les droits sur les autres immeubles (3,13~Md€ en 2025, 3,08~Md€ par an en 2023-2024) sont la différence entre le
total des collectivités (OFGL) et les logements (Compte du logement), deux sources de périmètres différents ; ils sont
ramenés au champ et au taux permanent par le rapport entre droits du modèle et droits perçus sur les logements en 2025
(0,907). Le sol des 4,2~millions de logements HLM loués est estimé à 603~Md€ en valeurs de l'Insee, avec la valeur
médiane des logements de leur commune (un majorant).

\begin{table}[!htbp]\centering\scriptsize
\caption{Le taux national selon l'assiette}\label{tab:fon_sens}
\begin{tabularx}{\linewidth}{@{}>{\raggedright\arraybackslash}X rrrrr@{}}
\toprule
Assiette & Taux (Insee) & Taux (modèle) & Paris & Dernier dixième & Hausses $>$ 300~€ \\
\midrule
Sol des logements seul, droits des logements seuls (N4) & 0,157~\% & 0,197~\% & +129~\% & +602~€ & 25,1~\% \\
Terrains bâtis privés & 0,131~\% & 0,165~\% & +108~\% & +500~€ & 20,7~\% \\
Tous les terrains privés (proposition) & 0,120~\% & 0,151~\% & +99~\% & +456~€ & 19,0~\% \\
Tous les terrains privés, sol des HLM exonéré & 0,132~\% & 0,166~\% & +109~\% & +503~€ & 20,9~\% \\
Tous les terrains privés, terres agricoles comprises & 0,106~\% & 0,133~\% & +87~\% & +399~€ & 15,8~\% \\
Tous les terrains privés, droits des seuls logements & 0,098~\% & 0,122~\% & +80~\% & +366~€ & 14,3~\% \\
\bottomrule
\end{tabularx}
\par\vspace{2pt}{\raggedright\scriptsize Sans réaction des ventes ; avec une reprise de 6~\% des ventes par point, taux
et hausses baissent d'environ 18~\%. Paris : variation moyenne de la taxe des logements ; dernier dixième : variation
moyenne par logement dans le dixième le plus cher. Les terres agricoles toucheraient des exploitants déjà soumis à la
taxe sur le foncier non bâti.\par}
\end{table}

\subsection{Résultats complets par dixième de logements}

\begin{table}[!htbp]\centering\scriptsize
\caption{Taxe foncière et droits par dixième de logements, selon le prix au m² de leur commune}\label{tab:fon_dec}
\setlength{\tabcolsep}{3.5pt}
\begin{tabularx}{\linewidth}{@{}c r *{4}{r} *{3}{r} *{3}{r}@{}}
\toprule
& Prix & \multicolumn{4}{c}{Taxe / valeur (\%)} & \multicolumn{3}{c}{Variation de taxe (€/an)} & \multicolumn{3}{c}{Taxe moins droits (€/an)} \\
\cmidrule(lr){3-6}\cmidrule(lr){7-9}\cmidrule(l){10-12}
Dixième & (€/m²) & auj. & naïve & N4 & prop. & naïve & N4 & proposition & naïve & N4 & prop. \\
\midrule
1 & 1\,035 & 0,54 & 0,67 & 0,55 & 0,55 & +114 & +16 & +10 (+2~\%) & +33 & $-$65 & $-$72 \\
2 & 1\,404 & 0,49 & 0,61 & 0,52 & 0,51 & +147 & +36 & +25 (+4~\%) & +42 & $-$69 & $-$80 \\
3 & 1\,693 & 0,49 & 0,61 & 0,53 & 0,52 & +175 & +53 & +38 (+5~\%) & +45 & $-$77 & $-$92 \\
4 & 2\,000 & 0,46 & 0,58 & 0,53 & 0,51 & +193 & +104 & +81 (+11~\%) & +47 & $-$42 & $-$65 \\
5 & 2\,356 & 0,43 & 0,54 & 0,51 & 0,49 & +212 & +144 & +113 (+14~\%) & +47 & $-$21 & $-$52 \\
6 & 2\,737 & 0,41 & 0,52 & 0,50 & 0,48 & +224 & +177 & +137 (+16~\%) & +36 & $-$10 & $-$50 \\
7 & 3\,176 & 0,39 & 0,49 & 0,49 & 0,47 & +235 & +236 & +185 (+20~\%) & +26 & +26 & $-$25 \\
8 & 3\,749 & 0,36 & 0,45 & 0,46 & 0,44 & +252 & +295 & +228 (+23~\%) & +20 & +63 & $-$4 \\
9 & 4\,746 & 0,31 & 0,38 & 0,41 & 0,39 & +256 & +358 & +274 (+27~\%) & $-$38 & +63 & $-$21 \\
10 & 6\,261 & 0,18 & 0,23 & 0,31 & 0,28 & +212 & +602 & +456 (+55~\%) & $-$258 & +132 & $-$14 \\
\bottomrule
\end{tabularx}
\par\vspace{2pt}{\raggedright\scriptsize Dixièmes de logements (et non de communes), classés selon le prix médian au m²
de leur commune. Taxe moins droits : baisse des droits rapportée à l'ensemble des logements de la commune. Sans réaction
des ventes.\par}
\end{table}

\textbf{Distribution individuelle (proposition).} Variation de la taxe foncière par logement : $-$9~€ au
10\textsuperscript{e} centile, +36~€ au 25\textsuperscript{e}, +110~€ à la médiane, +238~€ au 75\textsuperscript{e},
+422~€ au 90\textsuperscript{e}, +545~€ au 95\textsuperscript{e} et +647~€ au 99\textsuperscript{e} ; 7,2~\% des
logements dépassent +500~€, 0,9~\% +660~€ et 0,2~\% +1\,000~€ ; 11~\% voient leur taxe baisser (un sur neuf). Avec le taux
national sur le seul sol des logements, les mêmes centiles valent +4, +51, +141, +301, +529, +709 et +800~€
(figure~\ref{fig:fon_distrib}). Plus fortes hausses (cellules d'au moins 1\,000~logements) : maisons de Saint-Tropez
+2\,576~€ (+71~\%), de Megève +1\,927~€, appartements de Saint-Jean-Cap-Ferrat +1\,887~€, maisons de Ramatuelle +1\,799~€,
de La Clusaz +1\,742~€ et de Cassis +1\,673~€ (+109~\%). Hors Paris, le dixième le plus cher prend 40~\%. Par type, les
maisons prennent +112~€ et les appartements +212~€ en moyenne, mais +197 et +97~€ si l'on mesure le sol des maisons par
résidu, comme celui des appartements (coût de reconstruction recalé à 1\,973~€/m²).

\textbf{Vieillissement de la taxe actuelle.} Appliquer le rapport taxe/valeur de 2017 aux valeurs de 2023-2024, puis la
croissance de la taxe de 2017 à 2025, compte deux fois la hausse des prix de 2017 à 2023-2024. Une correction à trois
zones (indices Notaires-Insee : Paris ×1,113, Île-de-France hors Paris ×1,134, province ×1,281) ramène la hausse
parisienne de +99 à +88~\% et celle du dernier dixième de +55 à +52~\% ; la part des hausses de plus de 300~€ ne change
pas (19,0~\%). Une correction propre demande les valeurs de 2017 par commune.

\figbloc{foncier_n12_variantes}{Toutes les variantes de répartition : les deux taux nationaux inversent la régressivité, la taxe relevée l'aggrave}{fig:fon_var_toutes}

\begin{table}[!htbp]\centering\scriptsize
\caption{Toutes les variantes de répartition (droits sur les logements divisés par deux)}\label{tab:fon_variantes}
\begin{tabularx}{\linewidth}{@{}>{\raggedright\arraybackslash}X rrrr@{}}
\toprule
Variante & Premier dixième & Dernier dixième & Hausses $>$ 300~€ & Taxes qui doublent \\
\midrule
Taxe actuelle relevée (réforme naïve) & +114~€ (+25~\%) & +212~€ (+25~\%) & 14~\% & 0~\% \\
Valeur de marché, taux national unique & $-$72~€ ($-$16~\%) & +1\,247~€ (+149~\%) & 28~\% & 11,8~\% \\
Deux taux nationaux, sol = 2 × bâti & $-$149~€ ($-$33~\%) & +1\,559~€ (+187~\%) & 28~\% & 12,5~\% \\
Deux taux nationaux, sol = 4 × bâti & $-$213~€ ($-$47~\%) & +1\,821~€ (+218~\%) & 29~\% & 13,7~\% \\
Sol seul, taux national & $-$315~€ ($-$70~\%) & +2\,236~€ (+268~\%) & 29~\% & 16,3~\% \\
Deux taux par commune (2 × bâti), taux départemental sur le sol & +51~€ (+11~\%) & +463~€ (+55~\%) & 23~\% & 1,2~\% \\
Deux taux par commune (4 × bâti), taux départemental sur le sol & +45~€ (+10~\%) & +454~€ (+54~\%) & 25~\% & 1,9~\% \\
Deux taux par commune, taux national sur le sol des logements (N4) & +16~€ (+4~\%) & +602~€ (+72~\%) & 25~\% & 5,9~\% \\
N4 avec reprise des ventes (+6~\% par point) & +11~€ (+2~\%) & +487~€ (+58~\%) & 20~\% & 4,8~\% \\
Proposition (tous les terrains, tous les immeubles) & +10~€ (+2~\%) & +456~€ (+55~\%) & 19~\% & 4,5~\% \\
\bottomrule
\end{tabularx}
\par\vspace{2pt}{\raggedright\scriptsize Variation moyenne de la taxe foncière par logement. Taux nationaux, en part de
la valeur du modèle : 0,614 et 0,307~\% (sol = 2 × bâti), 0,753 et 0,188~\% (4 ×), 0,972~\% (sol seul). Taux
départemental : de 0,106~\% (Paris : 0,130~\%) à 0,789~\% (Creuse). Sources : \texttt{foncier/sorties/simulation.txt},
\texttt{sorties/calculs\_foncier.txt}.\par}
\end{table}

\figbloc{foncier_n13_distribution}{Le taux national sur tous les terrains réduit les plus fortes hausses d'un cinquième par rapport au taux sur le seul sol des logements}{fig:fon_distrib}

\subsection{Perte sèche}

Deux triangles (méthode : \texttt{revue/methode\_perte\_seche.md}). Sur les ventes, $\sigma V(\tau_1^2-\tau_0^2)/2$, où
$V$ = 212~Md€ de ventes (2025) et $\sigma$ la baisse des ventes par point de droits. Sur le capital bâti,
$\varepsilon_K (K/u)(t_1^2-t_0^2)/2$, où $t$ est l'impôt sur le bâti rapporté à sa valeur ($K$ = 3\,636~Md€), $u$ le
coût d'usage (4,5~\%) et $\varepsilon_K$ l'élasticité du capital logement (1). Le sol n'entre dans aucun triangle. Ces
pertes sèches sont statiques ; à long terme, que les ventes perdues en 2014 ne semblent pas l'avoir été durablement
(Bérard et Trannoy, 2018) plaide pour le bas de la fourchette.

\begin{table}[!htbp]\centering\scriptsize
\caption{Perte sèche évitée (Md€ par an ; signe moins = gain)}\label{tab:fon_ps}
\begin{tabularx}{\linewidth}{@{}>{\raggedright\arraybackslash}X r@{}}
\toprule
Ventes de logements, droits divisés par deux : $\sigma$ = 0 / 3 / 6 / 8 / 15 & 0 / $-$0,80 / $-$1,61 / $-$2,14 / $-$4,02 \\
Ventes de logements, baisse de 30~\% / de 100~\% ($\sigma$ = 6) & $-$1,09 / $-$2,14 \\
Bâti, proposition (impôt sur le bâti de 14,3 à 7,3~Md€) : $\varepsilon_K$ = 1 / 0,5 ; $u$ = 3,5 / 6~\% & $-$0,46 / $-$0,23 ; $-$0,59 / $-$0,34 \\
Bâti, rapport sol/bâti de 2 / 4 / sol seul (impôt sur le bâti de 10,7 / 7,3 / 0~Md€) & $-$0,27 / $-$0,46 / $-$0,62 \\
Bâti, réforme naïve (impôt sur le bâti de 14,3 à 17,9~Md€) & +0,36 \\
Bâti, en comptant les droits comme un impôt annuel sur la détention (0,18~\%) : proposition / réforme naïve & $-$0,98 / +0,04 \\
Ventes des autres immeubles ($\sigma$ = 6 transposé, 49~Md€ de ventes ; ordre de grandeur) & $-$0,37 \\
\midrule
Total, réforme naïve & $-$1,25 \\
Total, proposition, logements / avec les autres immeubles & $-$2,06 / $-$2,44 \\
Variante : coût des logements mal adaptés (Määttänen et Terviö, 2022) à la place du triangle sur les ventes & environ $-$3,8 \\
\bottomrule
\end{tabularx}
\par\vspace{2pt}{\raggedright\scriptsize La variante de Määttänen et Terviö applique aux 30,5~millions de logements le
gain par ménage de leur modèle d'Helsinki (euros de 2004), interpolé entre 30~€ par an pour un droit de 2~\% et 180~€ pour
6~\% (109~€ par logement, 3,3~Md€) ; elle remplace le triangle sur les ventes, qui mesure les mêmes échanges manqués.
Le loyer que se verse un propriétaire occupant n'étant pas imposé, alléger l'impôt sur le bâti peut aussi accroître une
distorsion existante : le signe du terme bâti n'est pas certain.\par}
\end{table}

\subsection{Capitalisation, report de paiement et transition}

\textbf{Capitalisation.} Un impôt annuel permanent $\Delta$ réduit la valeur d'un bien de $\Delta/r$, avec $r$ le taux
d'actualisation réel (3~\%, mesuré à 3-4~\% en Angleterre par Koster et Pinchbeck, 2022). Cas symétrique : la valeur
varie de $-(\Delta T + \Delta D)/r$, où $\Delta D$ est la baisse annuelle moyenne des droits par logement de la commune ;
cas asymétrique : de $-\Delta T/r$. Une montée linéaire sur $n$ années laisse supporter à l'annonce une part
$\sum_t \min(t/n, 1)(1+r)^{-t} \big/ \sum_t (1+r)^{-t}$ de la perte : 97~\% sur trois ans, 94~\% sur cinq, 88~\% sur dix,
77~\% sur vingt.

\begin{table}[!htbp]\centering\scriptsize
\caption{Variation de la valeur des logements des particuliers selon l'incidence (r = 3~\%)}\label{tab:fon_cap}
\begin{tabularx}{\linewidth}{@{}>{\raggedright\arraybackslash}p{3.4cm}*{4}{>{\raggedright\arraybackslash}X}@{}}
\toprule
& Ensemble & Somme des pertes & Appartement parisien médian & Logements perdant plus de 5~\% \\
\midrule
Proposition, cas symétrique & +48~Md€ (+0,7~\%) & 35~Md€ & $-$2\,970~€ ($-$0,7~\%) & 0,9~\% \\
Proposition, cas asymétrique & $-$157~Md€ ($-$2,3~\%) & 168~Md€ & $-$18\,152~€ ($-$4,4~\%) & 5,9~\% \\
Réforme naïve, cas asymétrique & $-$205~Md€ ($-$3,0~\%) & 205~Md€ & $-$4\,466~€ ($-$1,1~\%) & --- \\
\bottomrule
\end{tabularx}
\par\vspace{2pt}{\raggedright\scriptsize Cas symétrique : la hausse de taxe et la baisse des droits passent toutes deux
dans les prix ; cas asymétrique : seule la hausse de taxe y passe.\par}
\end{table}

\textbf{Report de paiement.} La hausse $\Delta$ reportée chaque année pendant $h$ ans au taux $i$ forme une dette de
$\Delta[(1+i)^h - 1]/i$. Pour l'appartement parisien médian ($\Delta$ = 545~€) : 6\,200~€ après dix ans, 10\,100~€ après
quinze, 14\,600~€ après vingt à 3~\% (1,5, 2,4 et 3,5~\% de la valeur) ; 9\,400~€ à 2~\% et 10\,900~€ à 4~\% après quinze
ans. Le taux de 3~\% est une convention : la règle est de facturer le taux d'emprunt de l'État.

\begin{table}[!htbp]\centering\scriptsize
\caption{Calendrier de la proposition sur cinq ans}\label{tab:fon_cal}
\begin{tabularx}{\linewidth}{@{}>{\raggedright\arraybackslash}X rrrrr@{}}
\toprule
Année & 1 & 2 & 3 & 4 & 5 \\
\midrule
Taux national (\% de la valeur du sol, Insee) & 0,024 & 0,048 & 0,072 & 0,096 & 0,120 \\
Recette du taux national (Md€) & 1,5 & 3,0 & 4,5 & 6,1 & 7,6 \\
Droits non remplacés (Md€) & 6,1 & 4,5 & 3,0 & 1,5 & 0 \\
Taxe de l'appartement parisien médian (€, 527 en 2025) & 636 & 745 & 854 & 963 & 1\,072 \\
\bottomrule
\end{tabularx}
\par\vspace{2pt}{\raggedright\scriptsize Droits divisés par deux dès l'annonce (année 0, dont les quelques mois avant le
1\textsuperscript{er}~janvier ne sont pas comptés) ; ventes et taxe foncière de 2025, en euros de 2025. Déficit cumulé :
15,2~Md€ (12,4 avec la reprise des ventes) ; sur trois ans 7,6 (6,2), sur dix ans 34,1 (27,8). Si la montée ne commence
qu'après trois années de droits baissés : 30,3 / 37,9 / 56,8~Md€ (montée sur 3, 5 ou 10 ans) ; avec une baisse des droits
en deux temps (un quart à l'annonce, un quart au démarrage), qui limite ce risque mais recrée l'attente : 18,9 / 26,5 /
45,5~Md€.\par}
\end{table}

\textbf{Autres calendriers.} Avec le taux national plein dès la première année, l'appartement parisien médian prendrait
d'un coup +534~€ (+698~€ avec le taux sur le seul sol des logements). Financer l'attente par une hausse uniforme de la
taxe actuelle demanderait 31~\% de hausse et porterait la taxe à 0,70~\% de la valeur dans le dixième le moins cher,
contre 0,24~\% dans le plus cher.

\textbf{Collectivités.} Droits perdus sur les logements (ventes de 2025, taux permanent) : 6,16~Md€, dont 4,77 pour les
départements, 1,27 pour les communes et 0,11 de frais de l'État ; s'y ajoutent 1,42~Md€ sur les autres immeubles. Reversé
à chaque département à hauteur de ses pertes, le taux national crée une péréquation : sur la seule part des logements,
Paris paie 621~M€ et perd 537~M€ (+83~M€), les Hauts-de-Seine +45~M€, la Creuse $-$4~M€ ; sept départements sont
contributeurs nets, pour 204~M€. Avec le taux sur le seul sol des logements, Paris paierait 811~M€ (+274~M€), les
Hauts-de-Seine +137~M€, et les transferts atteindraient 869~M€.

\subsection{Cas-types complets}

\begin{table}[!htbp]\centering\scriptsize
\caption{Le jeune ménage qui achète le logement médian de sa commune}\label{tab:fon_primo}
\setlength{\tabcolsep}{4pt}
\begin{tabularx}{\linewidth}{@{}>{\raggedright\arraybackslash}X rrrrrr@{}}
\toprule
Logement & Sol & Droits en moins & Taxe (€/an) & dont national & Revend après 7 / 15 ans & Économie rattrapée \\
\midrule
Paris, appartement & 85~\% & 12\,083 & +545 & 534 & +8\,690 / +5\,582 & 37 ans (naïve : jamais) \\
Lyon, appartement & 69~\% & 8\,105 & +303 & 288 & +6\,217 / +4\,488 & 55 ans (jamais) \\
Bordeaux, appartement & 66~\% & 6\,469 & +344 & 223 & +4\,326 / +2\,362 & 28 ans (jamais) \\
Nantes, maison & 46~\% & 13\,051 & +143 & 314 & +12\,161 / +11\,346 & jamais (65 ans) \\
Limoges, maison & 23~\% & 4\,988 & +253 & 61 & +3\,410 / +1\,965 & 30 ans (22 ans) \\
Saint-Étienne, appartement & 10~\% & 2\,266 & $-$42 & 12 & +2\,527 / +2\,766 & taxe en baisse (18 ans) \\
Annecy, maison & 30~\% & 20\,125 & $-$204 & 314 & +21\,394 / +22\,558 & taxe en baisse (jamais) \\
Arcachon, maison & 42~\% & 22\,700 & +123 & 496 & +21\,934 / +21\,233 & jamais (jamais) \\
Guéret, maison & 11~\% & 3\,193 & +27 & 18 & +3\,028 / +2\,876 & jamais (16 ans) \\
Ussel, maison & 10~\% & 2\,982 & +15 & 15 & +2\,886 / +2\,798 & jamais (13 ans) \\
\bottomrule
\end{tabularx}
\par\vspace{2pt}{\raggedright\scriptsize En euros ; taxe de 2025, droits au taux permanent, logement médian aux prix de
2023-2024. Revend : économie de droits moins la valeur actuelle (3~\%) de la hausse de taxe payée pendant la détention,
sans effet sur les prix. Économie rattrapée : durée au bout de laquelle la hausse actualisée égale l'économie (« jamais »
si la hausse perpétuelle actualisée reste inférieure) ; entre parenthèses, la réforme naïve. Part communale : variation
totale moins part nationale. Une famille nantaise qui passe d'un appartement à une maison de 450\,000~€ paie 13\,050~€ de
droits au lieu de 26\,100~€.\par}
\end{table}

\figbloc{foncier_n10_primo}{Un ménage qui achète le logement médian gagne 2\,000 à 13\,000~€ de droits ; la hausse de taxe foncière en reprend moins de la moitié en dix ans}{fig:fon_primo}

\subsection{Changements par rapport à la version précédente}

La version précédente reposait sur la taxe foncière de 2024, les ventes de 2023-2024 et les comptes de patrimoine de fin
2024. En base 2025, les ventes de logements anciens dépassent de 8~\% la moyenne de 2023-2024, la taxe foncière bâtie de
2,8~\% celle de 2024, et les terrains privés valent 4,9~\% de plus, surtout parce que l'Insee a relevé d'environ 5~\% ses
valeurs de 2024. Le taux national reste à 0,12~\%, la répartition des hausses ne change presque pas, et la bascule « taxe
moins droits » devient favorable dans tous les dixièmes ; les conclusions ne changent pas. Quatre chiffres de la version
précédente étaient faux (registre \texttt{revue/hypotheses\_foncier.md}, §~5.2) : la taxe foncière sur les logements
de 2024 valait 30,5~Md€, et non 31,6 (chiffre de 2025) ; le sol seul évite 0,6~Md€ de perte sèche sur le bâti, et non
0,5 ; la péréquation parisienne (« 213~M€ ») supposait un taux national limité à la part départementale des droits ;
et les « 660~€ » attribués à l'appartement parisien médian étaient le seuil des 5~\% de plus fortes hausses. Les dixièmes
portent désormais sur le prix au m² des communes et non sur leur richesse ; le coût des logements mal adaptés remplace le
triangle sur les ventes au lieu de s'y ajouter ; le taux national, qui porte sur tous les terrains privés, monte
progressivement, et la transition n'est plus financée par la taxe assise sur les valeurs de 1970.

\subsection{Hypothèses fortes et chiffres non vérifiés}

\begin{table}[!htbp]\centering\scriptsize
\caption{Hypothèses fortes}\label{tab:fon_hyp}
\begin{tabularx}{\linewidth}{@{}>{\raggedright\arraybackslash}p{4.3cm}>{\raggedright\arraybackslash}p{2.6cm}>{\raggedright\arraybackslash}X@{}}
\toprule
Hypothèse & Où elle joue & Ce qu'elle change si elle est fausse \\
\midrule
Dans une commune, taxe actuelle proportionnelle à la valeur de marché & variations individuelles & Gagnants et perdants plus nombreux et plus marqués (écarts entre quartiers : voir le corps) ; non chiffrable sans les fichiers fonciers \\
Logements d'une cellule commune × type identiques & parts de logements, cas-types & Parts de cellules : ordres de grandeur, pas des nombres d'avis \\
Valeur = prix médian × surface médiane (20~\% sous l'Insee, biais en U) & taux en part de la valeur & Taux du modèle surestimés d'un cinquième ; recettes justes ; régressivité confirmée après correction (environ 0,43 à 0,15~\%) \\
Sol des appartements : coût de reconstruction uniforme (1\,464~€/m²) & part nationale, bâti & 9~\% des logements à la borne de 10~\% ; avec 2\,500~€/m² à Paris, hausse parisienne inférieure d'environ 11~\% \\
Sol des maisons = terrains à bâtir × parcelle médiane & maisons / appartements & Mesuré par résidu, le partage s'inverse : ne rien en conclure \\
Taxe actuelle vieillie en comptant deux fois la hausse des prix & taxe actuelle par commune & Paris +88~\% au lieu de +99~\% ; dixièmes peu touchés \\
Calage sur 19,2~Md€ (2017), champ plus étroit que le modèle & hausses en pourcentage & Taxe actuelle peut-être sous-estimée de 7~\% au plus : hausses en pourcentage un peu surestimées \\
Ventes de 2025 au taux permanent de 5,81~\% & coût, taux national & Au niveau de 2021-2022, la baisse sur les logements coûterait près de 8~Md€ ; au taux temporaire de 6,32~\%, 9~\% de plus \\
Ventes en hausse de 6~\% par point de baisse (convention) & gain, coût & Sans réaction durable, 0,5~Md€ de gain (bâti) seulement ; gain linéaire en $\sigma$ \\
Incidence : trois cas polaires & capitalisation & Études contradictoires (acheteur en France en 2014, vendeur en Allemagne ; capitalisation complète ou nulle au Danemark) \\
Bâti : $\varepsilon_K$ = 1, $u$ = 4,5~\%, construction non bloquée par le droit des sols & gain sur le bâti & De 0,2 à 0,6~Md€ ; plus faible là où le zonage bloque la construction \\
Valeurs de l'Insee des terrains par secteur, même part de champ par secteur & assiette, taux national & Valeurs résiduelles, révisées de 9~\% (2023) puis de 5~\% (2024) ; terrains d'entreprise surestimés : taux et charge des logements plus élevés (sol bâti seul : 0,13~\%) \\
Charge sur les entreprises et les terrains nus non localisée & péréquation & Géographie des 2,9~Md€ non résidentiels inconnue : péréquation calculée pour les seuls logements \\
Droits sur les autres immeubles = différence de deux sources ; même réaction des ventes & coût (1,42~Md€), gain (0,4~Md€) & Montant et gain incertains ; variante limitée aux logements : 0,098~\% \\
Chaque commune garde sa recette ; locaux professionnels, HLM et taxe d'ordures ménagères hors modèle & part communale & Étendre la base aux locaux professionnels créerait des transferts non simulés \\
Valeurs provisoires du sol prêtes dès la première année & calendrier, déficit & Trois ans de retard : +23~Md€ de déficit ; validité juridique non expertisée \\
Actualisation et intérêt du report à 3~\% réels & valeurs actuelles, dette reportée & À 4~\%, la perte actualisée baisse d'un quart \\
Exonérations actuelles (plus de 75~ans modestes, neuf) non modélisées & retraités, taux national & À décider pour le taux national \\
Millésimes différents (taxe/valeur 2017, recensement 2022, prix 2023-2024, taxe et ventes 2025) & tout & Prix de 2025 inférieurs de 1,5~\% ; à actualiser pour une réforme en 2028-2030 \\
\bottomrule
\end{tabularx}
\end{table}

\begin{table}[!htbp]\centering\scriptsize
\caption{Chiffres non vérifiés sur une source primaire}\label{tab:fon_nv}
\begin{tabularx}{\linewidth}{@{}>{\raggedright\arraybackslash}p{5.2cm}>{\raggedright\arraybackslash}p{3.2cm}>{\raggedright\arraybackslash}X@{}}
\toprule
Chiffre & Où il sert & État \\
\midrule
Part immobilière des successions : 60~\% (12,7~Md€ en 2025) & inventaire & clé du CPO, aucune ventilation officielle \\
Taxe foncière hors logements : 11,5~Md€ en 2025 & inventaire & REI moins Compte du logement \\
Droits sur les autres biens : 3,1~Md€ en 2025 & coût de la baisse & OFGL moins Compte du logement \\
Frais d'assiette de l'État : 2,37~\% de la part départementale & taux de 5,81~\% & CPO et BOFiP ; article du CGI non relu \\
Taxe d'ordures ménagères assise sur la base de la taxe foncière & transition & article 1522 du CGI non relu \\
Plafond de revenu de l'exonération des plus de 75~ans (12\,793~€ pour une part en 2026) & retraités & source secondaire \\
Quotité moyenne financée de 84,7~\% (apport de 15~\%) & poids des droits dans l'apport & CCSF 2026 d'après la note de lecture \\
Taxe foncière non récupérable sur le locataire, déductible au régime réel & cas de la bailleuse & décret n°~87-713 et article 31 du CGI non relus \\
Coût de reconstruction de 1\,464~€/m² & part du sol & calé, jamais comparé à une source publiée \\
Produit du relèvement temporaire des droits en 2025 : 0,4~Md€ & droits de 2025 & suppose que le champ pèse la même part des droits qu'en 2023-2024 \\
Plus-values immobilières (2,8~Md€), contribution de sécurité immobilière (0,7~Md€), taxe sur les bureaux d'Île-de-France (1,1~Md€) & inventaire & chiffres de 2024 ; plus-values d'après l'audition de la Cour des comptes (rapport non lu) \\
Sol des HLM : 603~Md€ & variante d'exonération & estimation du modèle, majorant \\
Réaction des ventes : 6~\% par point & gain, coût & convention construite sur la littérature \\
Danemark : base de 80~\% de la valeur depuis 2024, report automatique des hausses & report par défaut & textes officiels lus via des outils de recherche (\texttt{lvt/notes/transitions\_et\_saillance.md}) \\
Territoire de la capitale australienne : transition ralentie & calendrier & presse et organismes favorables à la réforme \\
Recours constitutionnels pendants en Allemagne & expériences étrangères & source associative \\
Karlsruhe : maisons +76~\%, copropriétés $-$20~\% & expériences étrangères & calcul de la ville (présentation du 7/10/2024) \\
Banzhaf et Lavery : +3~points de croissance décennale du nombre de pièces & problème du bâti & version de travail de 2008 \\
Recours au report facultatif inférieur à 1~\% aux États-Unis & report par défaut & Baer (1998), cité par Shan (2010) \\
\bottomrule
\end{tabularx}
\end{table}
}{}
\section{Impôt sur les sociétés : calibrage, hypothèses et variantes}\label{ann:is}
\IfFileExists{parties/annexe_is.tex}{% Annexe de la partie IS (refonte du 16/11/2026 ; coupée le 17/11/2026 selon synthese/CONSIGNES_COUPE.md : figures
% is_n02 et is_n06 reçues du corps ; tableaux d'hypothèses fusionnés). Chiffres : ../is/sorties/calculs_note_is.txt ;
% registre des hypothèses : ../revue/hypotheses_is.md ; tableaux générés par ../is/calculs_note_is.py (parties/tab_is_*.tex).

\subsection{Calibrage}\label{ann:is_calibrage}

Chiffres : \texttt{is/calculs\_note\_is.py}, sur la mécanique de \texttt{is/simulation\_is.py} corrigée le 16 novembre
2026, recalée sur 2025 ; euros courants, l'année 1 étant 2027. Faute de données de 2025, la comparaison de l'OCDE
(2023), les taux implicites de l'Insee (2022) et la répartition des bénéfices par taille (2024) restent antérieurs.

\begin{table}[!htbp]\centering\scriptsize
\caption{Données et paramètres de la partie IS}\label{tab:is_calibrage}
\begin{tabularx}{\linewidth}{@{}>{\raggedright\arraybackslash}p{4.6cm}>{\raggedright\arraybackslash}p{3.3cm}>{\raggedright\arraybackslash}X@{}}
\toprule
Grandeur & Valeur & Source ou statut \\
\midrule
Recette de 2025 (2,8 millions de redevables) & IS net 59,9~Md€ au budget ; en prélèvements obligatoires, 59,3 (69,5
avant 10,1 de crédits d'impôt), contribution exceptionnelle 7,5, contribution sociale 1,15 : 68,0~Md€ (2,27~\% du PIB) ;
minimum mondial : 0,5~Md€ en 2026 & PLF 2027, rapport AN n°~2847 ; Insee 3.217 \\
Investissement éligible des sociétés non financières (autres machines et équipements, TIC) & 84,1~Md€ (2025) ; 89,2~Md€
(2027) & Insee 8.212 ; +3~\% par an (3,2~\% de 2022 à 2025) \\
Stock correspondant ; usure & 541~Md€ (fin 2025) ; 14~\% par an & Insee 8.212 \\
Bénéfices imposés en 2025 & 30,6~Md€ à 15~\% ; 256,8~Md€ à 25~\% ; déficits 111,3~Md€ & DGFiP n°~51 (provisoire) \\
Croissance du bénéfice fiscal & +2~\% par an en 2026 et 2027 & hypothèse ; de 2022 à 2025, la base au taux normal a crû
de 2,4~\%, celle à 15~\% de 15~\% \\
Un point de taux normal en 2027 & 2,99~Md€ (2,67 si le taux réduit est gardé) & bénéfices de 2025 $\times$ 1,040 \\
Fin du taux réduit & 2,15 à 3,18~Md€ & haut : 10 points sur 30,6~Md€ ; bas : coût officiel de 2019 ramené à 10 points,
porté à 3~\% par an \\
CIR apprécié au niveau du groupe ; fin des zones hors outre-mer & 0,59 ; 0,54~Md€ (IR et IS) & CPO (2025) ; estimation
de 2021 vieillie ; part de l'IS non publiée \\
Boîte à brevets à 17,5~\% & 0,21 à 0,62~Md€ & moitié de la dépense fiscale (1,26~Md€ en 2025, estimée par le PLF 2026) \\
Ressources de l'assiette & 3,50 / 4,22 / 4,94~Md€ & bas, central, haut, en 2027, puis +3~\% par an \\
Durée d'amortissement des équipements & 7 ans, linéaire ou dégressif (coefficient 2,25) & BOFiP, BOI-BIC-AMT-20-20-30 \\
\bottomrule
\end{tabularx}
\end{table}

\figbloc{is_n02_qui_paie}{370 grandes entreprises paient 29~\% de l'IS, 2,6 millions de très petites
sociétés 21~\% (2025)}{fig:is_qui}

\textbf{Le plafond de 1~M€.} Parmi les 1,2 million d'entreprises qui ont investi en 2021, les 2~\% qui ont investi au
moins 500\,000~€ chacune réalisent 87~\% des 248,7~Md€ d'investissements corporels, et les 59~\% qui ont investi moins
de 10\,000~€ n'en font que 1~\% (Insee, 2023). Une interpolation de Pareto entre ces seuils, prolongée jusqu'à 1~M€,
donne une part couverte de 21~\% (18~\% pour un plafond de 0,5~M€). Le champ de l'Insee est plus large que celui de la
mesure, et l'on ignore la composition des achats par taille : nous retenons 20~\% (15 à 30~\%). Au
Royaume-Uni, 23~Md£ d'investissement ont été déduits en 2019-2020 sous le plafond équivalent (IFS, 2023).

\textbf{Expériences étrangères.} La déduction totale britannique exclut voitures et biens loués ; elle coûte de
l'ordre de 10 à 11~Md£ par an au début (IFS, 2023). La déduction américaine de 100~\% de 2017, réduite de
20 points par an dès 2023, a été rétablie pour les biens acquis après le 19 janvier 2025, occasion comprise sauf entre
entreprises liées (EY, 2026).

\subsection{Méthode}\label{ann:is_methode}

\textbf{Coût annuel.} Les achats de l'année $v$ sont $I_v = I_{2027}(1+g)^{v}$. Avec une fraction $b_v$ déduite
immédiatement et le reste amorti selon le plan actuel $d_k$, le coût pour l'État l'année $t$ est
$\tau\, s \sum_{v \le t} b_v I_v\,(\mathbf{1}_{t=v} - d_{t-v})$, où $\tau = 25~\%$ et $s$ est la part imposable. Le
plafond est une fraction $\theta$ des achats déduite à 100~\% : $b_v = \theta + (1-\theta)\,p_v$, où $p_v$ est le
palier de l'année. Le cas central est la moyenne des plans linéaire et dégressif.

\textbf{Coût actualisé et coût permanent.} Une année d'achats coûte $\tau\, s\, I\,(1 - z)$ en valeur actualisée ($z$ :
valeur actuelle des déductions pour 1~€ : 0,847 en linéaire, 0,897 en dégressif à 5~\% ; 0,896 et 0,930 à
3,25~\%). En régime permanent, le coût annuel vaut $\tau\, s\, I_t\,(1 - z(g))$, où $z(g)$ est actualisé au taux de
croissance de l'investissement : il serait nul sans croissance.

\textbf{Dette et baisse finale.} La dette vaut celle de l'année précédente, intérêts à 3,25~\% compris, moins le solde
de l'année. L'année qui suit l'extinction, le taux normal baisse de $x = (F - C)/(P - C/25)$ points, où $F$, $C$ et $P$
sont les ressources, le coût de la déduction et la valeur d'un point en régime permanent (2040) ; $C/25$ traduit la
baisse du coût de la déduction avec le taux.

\textbf{Taux effectifs.} Méthode de Devereux et Griffith (2003) : financement par fonds propres, sans impôt personnel,
rendement réel exigé de 5~\%, inflation de 2~\%, usure de 14~\% (équipements) ou 3,1~\% (bâtiments, amortis sur 25 ans) ;
taux moyen effectif pour un projet rapportant 20~\% (tableau~\ref{tab:is_tef}). L'OCDE (2025), dont les conventions
diffèrent (note de méthode non lue), donne pour la France un taux moyen effectif de 23,6~\% contre 22,2~\% en moyenne
(13\up{e} sur 38) et un taux marginal effectif de 15,0~\% contre 16,5~\% (9,8~\% contre 10,4~\% pour les équipements),
sans semble-t-il la contribution exceptionnelle. Seuls les écarts entre variantes se comparent.

\begin{table}[!htbp]\centering\scriptsize
\caption{Taux effectifs d'imposition d'un investissement, aujourd'hui et selon les réformes}\label{tab:is_tef}
\begin{tabularx}{\linewidth}{@{}>{\raggedright\arraybackslash}X>{\raggedleft\arraybackslash}p{1.4cm}>{\raggedleft\arraybackslash}p{2.2cm}>{\raggedleft\arraybackslash}p{2.6cm}@{}}
\toprule
Cas & Taux légal & Taux marginal effectif & Taux moyen effectif (projet à 20~\%) \\
\midrule
Équipement, aujourd'hui, amortissement linéaire & 25~\% & 20,7~\% & 23,6~\% \\
Équipement, aujourd'hui, amortissement dégressif & 25~\% & 15,1~\% & 22,1~\% \\
Équipement, baisse du taux de même coût actualisé que la déduction & 24,0~\% & 19,8~\% & 22,7~\% \\
Équipement, taux à 22~\% & 22~\% & 18,1~\% & 20,8~\% \\
Équipement, taux à 20~\% & 20~\% & 16,3~\% & 18,9~\% \\
Équipement, déduction immédiate de 50~\% & 25~\% & 11,5~\% & 21,2~\% \\
Équipement, déduction immédiate de 100~\% & 25~\% & 0 & 18,8~\% \\
Équipement, petite société sous 42\,500~€, aujourd'hui & 15~\% & 12,1~\% & 14,2~\% \\
Bâtiment, aujourd'hui & 25~\% & 22,0~\% & 24,0~\% \\
Bâtiment, petite société sous 42\,500~€, aujourd'hui & 15~\% & 13,0~\% & 14,4~\% \\
Bâtiment, taux à 22~\% & 22~\% & 19,2~\% & 21,1~\% \\
\bottomrule
\end{tabularx}
\par\vspace{3pt}\parbox{\linewidth}{\raggedright\scriptsize Coin sur le coût d'usage des équipements, au taux nominal de
5~\% : 5,1~\% en linéaire, 3,5~\% en dégressif, 0 avec la déduction immédiate. Source : \texttt{is/calculs\_note\_is.py}.}
\end{table}

\textbf{Part d'une baisse du taux qui touche l'investissement en équipements.} L'impôt actualisé sur le rendement normal
d'une année d'achats d'équipements égale le coût actualisé de leur déduction immédiate, 1,73 à 2,55~Md€. Un point de taux
en retire un vingt-cinquième, 0,07 à 0,10~Md€, pour un coût de 2,67~Md€ par an : 2,6 à 3,8~\% du coût, un peu plus si
l'on comptait les autres actifs neufs. Par taille, la baisse irait pour 30~\% aux grandes entreprises, 25~\% aux
entreprises de taille intermédiaire, 28 à 31~\% aux PME et 14 à 17~\% aux très petites sociétés (bénéfices de 2024).

\textbf{Production et bien-être.} Perte sèche : \texttt{synthese/sorties/perte\_seche\_autres.txt} ; élasticité :
H18-H19. Avec 50~\% sans montée, les gains de production et de bien-être sont divisés par deux ;
avec le plafond seul, le gain de production n'est que de 0,02~\% environ. Le suramortissement français de 2015 n'a pas
d'évaluation causale connue ; la seule étude du suramortissement robotique de 2019-2020 trouve un effet de 12 à 14~\%
sur l'ensemble des PME industrielles, difficile à concilier avec ses 2\,400 utilisateurs (DGE, 2026).

\textbf{Cas types.} Bénéfice et achats constants ; valeur de la déduction immédiate égale à l'avance d'impôt actualisée à
5~\% par année d'achats, $\tau\, I\,(1 - z)$, soit 3,8~\% des achats en linéaire ; avance de trésorerie de la première
année, $\tau\, I\,(1 - d_0)$ ; contribution exceptionnelle du grand groupe calculée sur un impôt supposé stable
(tableau~\ref{tab:is_cas}).

\begin{table}[!htbp]\centering\scriptsize
\caption{Quatre sociétés types : ce que change la proposition (€ par an)}\label{tab:is_cas}
\begin{tabularx}{\linewidth}{@{}>{\raggedright\arraybackslash}X>{\raggedleft\arraybackslash}p{1.45cm}>{\raggedleft\arraybackslash}p{1.45cm}>{\raggedleft\arraybackslash}p{1.55cm}>{\raggedleft\arraybackslash}p{1.9cm}>{\raggedleft\arraybackslash}p{1.4cm}@{}}
\toprule
Société & Impôt actuel & Fin du taux réduit & Valeur de la déduction & Variation nette & Trésorerie 2027 \\
\midrule
N'investit pas, bénéfice 42\,500~€ & 6\,375 & +4\,250 & 0 & +4\,250 (+67~\%) & +4\,250 \\
PME de services, bénéfice 40\,000~€, achats 10\,000~€ & 6\,000 & +4\,000 & $-$382 & +3\,618 (+60~\%) & +1\,679 \\
PME industrielle, bénéfice 300\,000~€, achats 150\,000~€ & 70\,750 & +4\,250 & $-$5\,727 & $-$1\,477 ($-$2~\%) & $-$30\,571 \\
Grand groupe, bénéfice 500~M€, achats 300~M€ & 129,1~M & 0 & $-$11,5~M & $-$11,5~M ($-$9~\%) & $-$34,9~M \\
Grand groupe : fin de la contribution exceptionnelle & 180,6~M (2026) & & & $-$51,5~M ($-$29~\%) & $-$51,5~M \\
\bottomrule
\end{tabularx}
\par\vspace{3pt}\parbox{\linewidth}{\raggedright\scriptsize Signe positif : la société paie davantage. Régime permanent,
avant la baisse finale du taux (3 à 6,5~M€ par an pour le grand groupe). Les PME achètent sous le plafond et déduisent
100~\% dès 2027 ; le grand groupe déduit 50~\% de ses achats au-delà de 1~M€ en 2027. Une société qui amortit
aujourd'hui en dégressif doit acheter 164\,000~€ par an pour compenser la perte de 4\,250~€.}
\end{table}

\subsection{Variantes}\label{ann:is_variantes}

\textbf{Calendriers.} Le tableau~\ref{tab:is_options} compare sept calendriers, la figure~\ref{fig:is_calendriers} le
coût annuel de trois d'entre eux, la figure~\ref{fig:is_dette_cal} la dette de ceux qui en créent une. L'amortissement
neutre coûte à l'État, à son taux, autant que la déduction immédiate (1,2 à 1,7~Md€ par année d'achats), mais ne vaut
pour les entreprises que 1,1 à 1,6~Md€, contre 1,7 à 2,6. Une marche unique en 2033 ferait bondir le coût de 2 à 10~Md€ ; le
gain d'attente qu'elle crée vaut un mois de rendement d'une machine rapportant 15 à 20~\% par an, contre une semaine
avec des paliers (H28). Le tableau~\ref{tab:is_fin} donne la dette selon le financement.

\begin{table}[!htbp]\centering\scriptsize
\caption{Les calendriers de la déduction immédiate comparés (Md€ ; part imposable 75~\%)}\label{tab:is_options}
\begin{tabularx}{\linewidth}{@{}>{\raggedright\arraybackslash}X>{\raggedleft\arraybackslash}p{1.45cm}>{\raggedleft\arraybackslash}p{1.45cm}>{\raggedleft\arraybackslash}p{1.45cm}>{\raggedleft\arraybackslash}p{1.45cm}>{\raggedleft\arraybackslash}p{1.45cm}@{}}
\toprule
Calendrier & Coût 2027 & Plus forte hausse annuelle & Coût actualisé par année d'achats & Dette maximale (cas central) & Extinction (central ; bornes) \\
\midrule
50~\% de 2027 à 2032, 100~\% en 2033 (ancienne note) & 7,0--7,8 & 7,9 & 1,7--2,6 & 11,2 & 2041 (2030 à 2051) \\
50~\% en 2027, +10 points par an, 100~\% en 2032 & 7,0--7,8 & 0,1 & 1,7--2,6 & 21,4 & 2042 (2036 à 2053) \\
100~\% sous 1 M€ dès 2027 ; au-delà, paliers 50 à 100~\% (2032) & 8,4--9,3 & 0,1 & 1,7--2,6 & 22,9 & 2042 (2036 à 2053) \\
100~\% sous 1 M€ par groupe, sans plus (AIA britannique seule) & 2,8--3,1 & 0,0 & 0,3--0,5 & aucune & --- \\
100~\% pour tous dès 2027 & 14,0--15,5 & 0,1 & 1,7--2,6 & 30,5 & 2043 (2036 à 2054) \\
50~\% pour tous, sans montée & 7,0--7,8 & 0,0 & 0,9--1,3 & 5,7 & 2033 (2030 à 2036) \\
Amortissement neutre (annuités revalorisées à 3,25~\%) & 0 & 0,4 & 1,1--1,6 & aucune & --- \\
\bottomrule
\end{tabularx}
\par\vspace{3pt}\parbox{\linewidth}{\raggedright\scriptsize Coût 2027 : régime actuel dégressif -- linéaire. Plus forte hausse annuelle : saut du coût d'une année sur l'autre (cas central) ; la falaise de 2033 en est le cas extrême. Coût actualisé à 5~\% (taux des entreprises), en régime permanent. Bornes de l'extinction : cas favorable (dégressif, ressources hautes) et défavorable (linéaire, ressources basses). Plafond : 1~M€ par groupe, couvrant 20~\% de l'investissement éligible (hypothèse). Amortissement neutre : plans actuels, annuités revalorisées au taux d'emprunt de l'État. Source : \texttt{is/calculs\_note\_is.py}.}
\end{table}


\figbloc{is_n06_calendriers}{Le calendrier proposé évite la marche de 2033 : son coût décroît dès 2027 au lieu de
bondir de 8~Md€}{fig:is_calendriers}

\figbloc{is_n09_dette_calendriers}{La dette de transition selon le calendrier : 29~Md€ au plus pour 100~\% dès 2027,
22~Md€ pour la proposition, 9~Md€ pour l'ancien calendrier}{fig:is_dette_cal}

\begin{table}[!htbp]\centering\scriptsize
\caption{La dette de transition selon le financement (proposition, cas central pour le coût)}\label{tab:is_fin}
\begin{tabularx}{\linewidth}{@{}>{\raggedright\arraybackslash}X rrr@{}}
\toprule
Financement & Ressources 2027 (Md€) & Dette maximale (Md€) & Extinction \\
\midrule
Assiette (proposition) & 3,5 à 4,9 & 16 à 28 & 2038 à 2046 \\
Assiette, taux réduit supprimé en deux ans (20~\% en 2027) & 2,9 en 2027 (central) & 23 (central) & 2041 (central) \\
Assiette, taux réduit porté à 20~\% & 2,4 à 3,4 & 29 à 40 & 2047 à 2067 \\
Assiette, taux réduit retiré aux seules sociétés immobilières et financières (hypothèse) & 1,9 à 2,7 & 36 (haut) ; 75 en
2066 (bas) & 2057 (haut) ; jamais (bas) \\
Assiette, taux réduit maintenu & 1,3 à 1,8 & 88 à 143 en 2066 & jamais \\
Taux normal à 28~\%, taux réduit gardé, sans autre mesure & 4,6 à 8,0 & 1 à 18 & 2030 à 2039 \\
Assiette et une dernière année de contribution exceptionnelle (5,0 en 2027, 0,5 en 2028) & 3,5 à 4,9 + 5,5 & 10 à 21
(central 15) & 2036 à 2043 (central 2039) \\
\bottomrule
\end{tabularx}
\par\vspace{3pt}\parbox{\linewidth}{\raggedright\scriptsize Régime actuel moitié linéaire, moitié dégressif ; part
imposable 75~\%. « Jamais » : dette non éteinte en 2066. Sociétés immobilières et financières : base à 15~\% supposée
répartie comme le nombre de sociétés (31~\% en 2025) ou l'IS brut (24~\%), invérifiable (H25). Source :
\texttt{is/calculs\_note\_is.py}.}
\end{table}

\subsection{La transition, année par année}\label{ann:is_transition}

\begin{table}[!htbp]\centering\scriptsize
\setlength{\tabcolsep}{3.5pt}
\caption{Transition année par année, avec les bornes (Md€ courants)}\label{tab:is_transition_complet}
\begin{tabular}{@{}l>{\raggedleft\arraybackslash}p{1.3cm}>{\raggedleft\arraybackslash}p{1.3cm}>{\raggedleft\arraybackslash}p{1.3cm}>{\raggedleft\arraybackslash}p{1.3cm}>{\raggedleft\arraybackslash}p{1.3cm}>{\raggedleft\arraybackslash}p{1.3cm}>{\raggedleft\arraybackslash}p{1.3cm}>{\raggedleft\arraybackslash}p{1.3cm}@{}}
\toprule
& \multicolumn{4}{c}{Cas central} & \multicolumn{2}{c}{Cas favorable} & \multicolumn{2}{c}{Cas défavorable}\\
\cmidrule(lr){2-5}\cmidrule(lr){6-7}\cmidrule(lr){8-9}
Année & Coût & Ressources & Solde & Dette & Coût & Dette & Coût & Dette \\
\midrule
2027 & 8,9 & 4,1 & $-$4,8 & 4,8 & 8,4 & 3,8 & 9,3 & 5,8 \\
2028 & 8,3 & 4,2 & $-$4,1 & 9,1 & 7,1 & 6,3 & 9,4 & 11,8 \\
2029 & 7,9 & 4,3 & $-$3,6 & 12,9 & 6,3 & 7,9 & 9,4 & 17,9 \\
2030 & 7,5 & 4,4 & $-$3,1 & 16,4 & 5,9 & 8,9 & 9,2 & 23,9 \\
2031 & 7,2 & 4,6 & $-$2,6 & 19,6 & 5,6 & 9,6 & 8,8 & 29,6 \\
2032 & 6,8 & 4,7 & $-$2,0 & 22,3 & 5,2 & 9,8 & 8,3 & 34,8 \\
2033 & 4,7 & 4,8 & 0,1 & 22,9 & 3,4 & 7,9 & 6,1 & 37,8 \\
2034 & 3,6 & 5,0 & 1,4 & 22,2 & 2,6 & 5,1 & 4,5 & 39,2 \\
2035 & 2,9 & 5,1 & 2,2 & 20,7 & 2,1 & 1,6 & 3,7 & 39,8 \\
2036 & 2,4 & 5,3 & 2,9 & 18,5 & 1,8 & $-$2,6 & 3,0 & 39,5 \\
2037 & 2,1 & 5,5 & 3,4 & 15,7 & 1,6 & $-$7,3 & 2,6 & 38,7 \\
2038 & 1,9 & 5,6 & 3,7 & 12,5 & 1,5 & $-$12,5 & 2,3 & 37,5 \\
2039 & 1,9 & 5,8 & 3,9 & 9,1 & 1,5 & $-$17,9 & 2,3 & 36,0 \\
2040 & 2,0 & 6,0 & 4,0 & 5,4 & 1,6 & $-$23,7 & 2,4 & 34,4 \\
2041 & 2,0 & 6,1 & 4,1 & 1,4 & 1,6 & $-$29,8 & 2,4 & 32,7 \\
2042 & 2,1 & 6,3 & 4,2 & $-$2,7 & 1,7 & $-$36,3 & 2,5 & 30,8 \\
2043 & 2,2 & 6,5 & 4,3 & $-$7,2 & 1,7 & $-$43,1 & 2,6 & 28,8 \\
2044 & 2,2 & 6,7 & 4,5 & $-$11,9 & 1,8 & $-$50,4 & 2,7 & 26,6 \\
2045 & 2,3 & 6,9 & 4,6 & $-$16,9 & 1,8 & $-$58,1 & 2,7 & 24,3 \\
2046 & 2,4 & 7,1 & 4,8 & $-$22,2 & 1,9 & $-$66,2 & 2,8 & 21,8 \\
2047 & 2,4 & 7,3 & 4,9 & $-$27,8 & 2,0 & $-$74,7 & 2,9 & 19,1 \\
2048 & 2,5 & 7,6 & 5,0 & $-$33,8 & 2,0 & $-$83,7 & 3,0 & 16,2 \\
2049 & 2,6 & 7,8 & 5,2 & $-$40,0 & 2,1 & $-$93,2 & 3,1 & 13,1 \\
2050 & 2,7 & 8,0 & 5,3 & $-$46,7 & 2,1 & $-$103,2 & 3,2 & 9,8 \\
2051 & 2,7 & 8,3 & 5,5 & $-$53,7 & 2,2 & $-$113,8 & 3,3 & 6,3 \\
2052 & 2,8 & 8,5 & 5,7 & $-$61,1 & 2,3 & $-$124,9 & 3,4 & 2,6 \\
2053 & 2,9 & 8,8 & 5,8 & $-$69,0 & 2,3 & $-$136,6 & 3,5 & $-$1,4 \\
2054 & 3,0 & 9,0 & 6,0 & $-$77,2 & 2,4 & $-$148,9 & 3,6 & $-$5,6 \\
2055 & 3,1 & 9,3 & 6,2 & $-$86,0 & 2,5 & $-$161,8 & 3,7 & $-$10,1 \\
2056 & 3,2 & 9,6 & 6,4 & $-$95,1 & 2,6 & $-$175,4 & 3,8 & $-$14,8 \\
\bottomrule
\end{tabular}
\par\vspace{3pt}\parbox{\linewidth}{\raggedright\scriptsize Cas favorable : régime actuel dégressif, ressources hautes (4,62~Md€ en 2027). Cas défavorable : régime actuel linéaire, ressources basses (3,50~Md€). Cas central : moitié-moitié, ressources au milieu (4,06~Md€). Une dette négative est un excédent cumulé : dans la proposition, il n'apparaît pas, car l'excédent est rendu par la baisse du taux normal dès l'année qui suit l'extinction (le tableau montre ce qui se passerait sans cette baisse). Part imposable : 75~\%. Source : \texttt{is/calculs\_note\_is.py}.}
\end{table}


\subsection{Hypothèses fortes et chiffres non vérifiés}\label{ann:is_hypotheses}

Registre : \texttt{revue/hypotheses\_is.md} (H1 à H22 ; H23 à H28 ajoutées ici). Effets sur la proposition, cas
central, sauf mention.

\begin{table}[!htbp]\centering\scriptsize
\caption{Hypothèses fortes de la partie IS}\label{tab:is_hyp1}
\begin{tabularx}{\linewidth}{@{}>{\raggedright\arraybackslash}p{4.4cm}>{\raggedright\arraybackslash}p{2.4cm}>{\raggedright\arraybackslash}X@{}}
\toprule
Hypothèse & Où elle joue & Ce qu'elle change si elle est fausse \\
\midrule
H1. 75~\% de l'investissement éligible fait par des sociétés imposables (60 à 90~\%), sans statistique publiée & coûts,
dette, PIB & L'hypothèse qui pèse le plus (§~\ref{sec:is_transition}). \\
H2. Les sociétés imposables utilisent toute la déduction l'année de l'achat, les autres ne coûtent rien & profil du
coût & Si les autres imputaient un jour tous leurs déficits : coût actualisé majoré d'un tiers au plus (0,6 à
0,85~Md€ par année d'achats). Revalorisation des déficits non chiffrable. \\
H3. Champ : machines hors transport et informatique des sociétés non financières & coûts & Avec le transport,
déduction totale à 23,3~Md€ la première année au lieu de 15,5 ; entreprises à l'IR : au plus 14,5~Md€ d'assiette. \\
H4. Amortissement type sur 7 ans & coût, coin & À 5 ou 10 ans, coût actualisé de 1,9 ou 3,5~Md€ au lieu de 2,6 ; à 10
ans, coin de 7,0~\% au lieu de 5,1. \\
H5. Plan actuel moitié linéaire, moitié dégressif (mélange réel non publié) & fourchettes & Dette de 13 à 32~Md€ avec
les ressources centrales. \\
H6. Investissement +3~\% par an & coût permanent, dette & À 2 ou 4~\% : dette de 20,2 ou 22,9~Md€. \\
H7. Bénéfice fiscal +2~\% par an en 2026 et 2027 & point de taux & De 0 à 2,5~\% : point de 2,87 à 3,02~Md€. \\
H8-H9. Actualisation à 5~\% pour les entreprises, emprunt de l'État à 3,25~\% & valeurs, seuils, dette & À 8~\%, la
déduction vaut plus pour les entreprises ; de 2,5 à 5~\%, dette de 20,9 à 23,1~Md€. \\
H10-H11. Aucun investissement suscité ; coûts comptés l'année de l'exercice & premières années, calendrier & Surcoût
temporaire, récupéré ensuite ; en caisse, un an plus tard. \\
H12. Bénéfices déclarés : 0 à 1,5~\% de moins par point de taux, sur tous les bénéfices & IS à 28~\%, baisse du taux &
Mesuré sur les multinationales : borne basse du 28~\% pessimiste. \\
H13. Coût réel du transfert de bénéfices : 0 à 100~\% de l'impôt perdu & perte sèche du 28~\% & Avec un pays d'accueil
à 12,5-15~\%, borne haute de 1,5 au lieu de 3,2~Md€. \\
H14. Fin du taux réduit : borne basse au coût officiel de 2019, non relevée au recalage (68~\% du rendement statique,
contre 75~\% auparavant) & ressources & Réactions (passage à l'IR, rémunération en
salaire) non chiffrées. \\
H15. Boîte à brevets : 0,62~Md€, 0,2 si la base fuit (borne arbitraire) & ressources & Quelques Md€ de dette. \\
H16-H17. CIR au niveau du groupe et zones : chiffrages du CPO & ressources & CIR estimé en 2021 (liasses de 2016) ;
part de l'IS des zones non publiée. Au plus 0,5~Md€ d'écart. \\
H18-H19. $r+\delta = 19~\%$ ; élasticité du stock d'équipements de 0,5 à 1,6 (0,5 sans source) & PIB, perte sèche &
Effet proportionnel ; dépend de la taille si la déduction agit par la trésorerie. \\
H20. Gain de PIB hors effet du financement & PIB & Voir §~\ref{sec:is_prop} (petites sociétés). \\
H21-H22. Taux de 25~\% pour tous dans le coût ; référence : droit en vigueur & coûts, recette & Négligeable ; fin de la
contribution exceptionnelle affichée (§~\ref{sec:is_arbitrages}). \\
H23. Plafond couvrant 20~\% de l'investissement éligible (15 à 30~\%) & premières années & Coût 2027 de 8,5 à
9,6~Md€, dette de 21,2 à 22,3~Md€. \\
H24. Ressources au milieu de la fourchette (4,22~Md€) & transition & Bornes de 3,5 et 4,9~Md€ : dette de 8 à 40~Md€
avec les bornes de H5. \\
H25. Taux réservé aux sociétés opérationnelles : base à 15~\% répartie comme le nombre de sociétés ou l'IS brut &
variante de financement & Rendement de 0,5 à 1,0~Md€ : ordre de grandeur indicatif, invérifiable. \\
H26. Taux effectifs : conventions de Devereux et Griffith, fonds propres & figures~\ref{fig:is_cout_egal} et
\ref{fig:is_avant_apres} & Avec emprunt, l'équipement est subventionné ; nous ne durcissons pas le plafonnement des
charges financières (3~M€ ou 30~\% du résultat avant intérêts et amortissements), qui épargne la plupart des PME. \\
H27. Cas types : déduction valorisée à 5~\% ; grand groupe sans CIR ni brevets & cas types & À 8-10~\%, la déduction
vaut 50 à 80~\% de plus pour les PME. \\
H28. Gain d'attente comparé à un rendement brut de 15 à 20~\% par an & calendrier & Ordre de grandeur, non mesuré. \\
\bottomrule
\end{tabularx}
\end{table}

\begin{table}[!htbp]\centering\scriptsize
\caption{Chiffres non vérifiés sur une source primaire (partie IS)}\label{tab:is_nonverif}
\begin{tabularx}{\linewidth}{@{}>{\raggedright\arraybackslash}p{5.6cm}>{\raggedright\arraybackslash}X@{}}
\toprule
Chiffre ou fait & Statut \\
\midrule
Environ 450 redevables de la contribution exceptionnelle & Source : le cours ; ni la DGFiP ni le PLF 2027 ne le donnent \\
Contribution exceptionnelle de 2026 : environ 7,0~Md€ & Déduite du PLF 2027 (5,0~Md€ en 2027, « en diminution de
2,0~Md€ ») \\
Suppression de l'article 15 du PLF 2027 en première lecture & Source secondaire (Légifiscal) \\
Bénéfices de 2025 ; croissance en 2026 et 2027 & DGFiP n°~51 provisoire ; tome I des Voies et moyens non publié \\
Contribution sociale de 2025 : 1,15~Md€ & Insee 3.217 ; le budget de l'État inscrit 1,5~Md€ en 2026, concept non rapproché \\
Contribution exceptionnelle dans les taux effectifs de l'OCDE & Déduite des données \\
Maffini, Xing et Devereux (2019) : +11~\% & Version de travail ; le résumé publié donne +2,1 à 2,5 points de taux
d'investissement \\
Bernini et Treibich (2016) ; Devereux, Liu et Loretz (2014) & Résumés seulement lus \\
Règles américaines (biens d'occasion, entreprises liées, 19 janvier 2025) & Source secondaire (EY, 2026) \\
CIR calculé société par société & D'après le CPO (2025) ; CGI non relu \\
Répartition de la base à 15~\% par type de société & Aucune statistique \\
\bottomrule
\end{tabularx}
\end{table}
}{}
\section{Revenus du capital : calibrage, hypothèses et variantes}\label{ann:pfu}
\IfFileExists{parties/annexe_pfu.tex}{% Annexe de la partie « Revenus du capital » (refonte du 16/11/2026). Chiffres : ../pfu/sorties/calculs_pfu.txt
% (programme ../pfu/calculs_pfu.py) ; données : ../pfu/donnees/ ; vérifications : ../revue/hypotheses_pfu.md.
% Coupe du 17/11/2026 : reçoit du corps les figures pfu_n02, n04 et n09, le détail de la charge du 1 % et des
% variantes du report d'apport-cession.

Chiffres produits par \texttt{pfu/calculs\_pfu.py} (DGFiP, Insee, Banque de France, Tax Foundation). L'assiette de
2025 applique aux cases de la déclaration de 2024 l'évolution des revenus déclarés mesurée par la DGFiP ($-$12,9~\%
pour les dividendes et les intérêts, $-$26,4~\% pour les plus-values).

\subsection{Assiette, recettes et taux cumulés}

\begin{table}[!htbp]\centering\scriptsize
\caption{Assiette et recettes des prélèvements sur les revenus du capital (Md€)}\label{tab:pfu_a_assiette}
\begin{tabularx}{\linewidth}{@{}>{\raggedright\arraybackslash}X>{\raggedleft\arraybackslash}p{1.5cm}>{\raggedright\arraybackslash}p{6.2cm}@{}}
\toprule
Poste & Montant & Précisions \\
\midrule
Dividendes (case 2DC), revenus de 2024 & 50,1 & \\
Intérêts et autres distributions (2TR, 2TS) & 11,8 & \\
Plus-values et actifs numériques, moins les pertes (3VG, 3AN, 3VH) & 17,7 & \\
\textbf{Assiette du PFU, revenus de 2024} & \textbf{79,6} & y compris les 2,54~M de foyers au barème (majorant
d'écart faible) \\
\textbf{Assiette du PFU, revenus de 2025 (estimée)} & \textbf{67,0} & dividendes 43,7, intérêts 10,3, plus-values
13,0 ; contrôles : 65,5 et 72,5 \\
Plus-values placées en report d'apport-cession (3WH) & 22,2 & 33,6 en 2023 \\
Stock de reports (8UT) & 173,3 & majorant ; 127 en 2023 (évaluation préalable du PLF 2027) \\
Plus-values en sursis d'\emph{exit tax} (3WA) & 4,7 & 358 foyers \\
\midrule
IR au PFU (impôt « hors barème »), revenus de 2025 & 8,4 & 10,2 sur les revenus de 2024 (acompte de 2024 : 7,7) \\
Prélèvements sociaux sur l'assiette du PFU de 2025 & 12,5 & au taux de 2026 ; 11,5 au taux de 2025 \\
Prélèvements sociaux sur tous les revenus du capital, encaissés en 2025 & 36,2 & au taux de 17,2~\% (assiette de
210~Md€) \\
CEHR (tous revenus), revenus de 2025 & 1,4 & 1,9 sur ceux de 2024 ; part payée sur le capital inconnue \\
CDHR, revenus de 2025 & 0,62 & environ 14\,000 foyers ; acompte de 2025 : 0,37 (7\,203 foyers) \\
\bottomrule
\end{tabularx}
\par\vspace{3pt}\parbox{\linewidth}{\raggedright\scriptsize Sources : DGFiP (totaux 2042 ; Statistiques n°~41, 43
et 53) ; Assemblée nationale, rapports n°~2847 et 3056.}
\end{table}

Taux cumulés sur 100~€ de bénéfice (droit de 2026) : 48,6~\% au PFU avec l'IS de 25~\% ; avec l'IS de 25,8~\%
(contribution sociale de 3,3~\% sur l'impôt qui dépasse 763\,000~€), 49,1~\%, 52,1~\% avec la CEHR et 54,5~\% au
plancher de la CDHR ; avec l'IS de 36,1~\%, 56,2, 58,7 et 60,8~\%. Holding : 26,8~\%, 26,0~\% en intégration
fiscale. Salaire : 55,4~\% hors cotisations, comparaison favorable au capital (le bénéfice répond à un coût du travail
cotisations patronales comprises). La Tax Foundation retient pour la France l'IS
de 36,1~\% et un taux personnel de 34~\%, soit 57,9~\% ; avec l'IS de 25,8~\%, 51,0~\% en 2025 (5\textsuperscript{e}
sur 31) et 52,1~\% en 2026 (4\textsuperscript{e}, face aux taux 2025 des autres pays).

\figbloc{pfu_n04_europe}{Avec l'IS de droit commun, la France reste parmi les cinq pays européens qui taxent le plus le bénéfice distribué}{fig:pfu_europe}

\figbloc{pfu_n02_concentration}{Quatre mille foyers reçoivent un tiers des dividendes et 70~\% des plus-values imposées}{fig:pfu_concentration}

\subsection{Changer le taux : méthode et résultats}

\emph{Méthode.} L'assiette $B$ réagit de $\Delta Q = e\,B\,\Delta\ln(1-t)$, $e$ étant l'élasticité des revenus
déclarés au taux net (0,77 chez Lefebvre, Lehmann et Sicsic ; 0,65 au CAE ; 1 en hypothèse haute). Premières années :
toute la réaction compte comme recette, $t_1(B+\Delta Q) - t_0 B$ (ce qui redonne le chiffrage du CAE : 1,16~Md€
recalculé, 1,2 publié). Régime permanent : seule rapporte la part $\pi$ de la réaction qui aurait sinon échappé pour
toujours à l'impôt (purge au décès, holding jamais vidée), d'où $(t_1-t_0)B + \pi\,t_1\,\Delta Q$ ; $\pi$ n'est
mesuré nulle part.

\begin{table}[!htbp]\centering\scriptsize
\caption{Recette (+) ou coût (--) d'un changement du taux du PFU (Md€ par an)}\label{tab:pfu_a_taux}
\begin{tabularx}{\linewidth}{@{}>{\raggedright\arraybackslash}X>{\raggedleft\arraybackslash}p{1.4cm}>{\raggedleft\arraybackslash}p{3.2cm}>{\raggedleft\arraybackslash}p{3.3cm}@{}}
\toprule
Réforme, base & Mécanique & Premières années ($e$ = 0,65 / 0,77 / 1) & Permanent ($e$ = 0,77 ; $\pi$ = 0 / 0,5 / 1) \\
\midrule
PFU à 25~\%, revenus de 2025 & $-$4,29 & $-$3,32 / $-$3,14 / $-$2,79 & $-$4,29 / $-$3,71 / $-$3,14 \\
PFU à 25~\%, revenus de 2024 & $-$5,10 & $-$3,94 / $-$3,73 / $-$3,32 & $-$5,10 / $-$4,41 / $-$3,73 \\
PFU à 33~\%, revenus de 2025 & 1,07 & 0,73 / 0,67 / 0,55 & 1,07 / 0,87 / 0,67 \\
PFU à 33~\%, revenus de 2024 & 1,27 & 0,87 / 0,80 / 0,65 & 1,27 / 1,04 / 0,80 \\
\bottomrule
\end{tabularx}
\end{table}

\emph{Taux qui maximise la recette}, sans réaction croisée : $1/(1+e)$, soit 60,6 / 56,5 / 50,0~\% pour $e$ = 0,65 /
0,77 / 1 ; lever un euro de plus à 31,4~\% demande d'alourdir de 1,42 / 1,54 / 1,84~€ l'impôt calculé à comportements
inchangés (ce ne sont pas des pertes sèches). Avec la réaction des revenus du travail : 43~\% (Lefebvre, Lehmann et
Sicsic ; 57~\% sans elle).

\emph{Perte sèche d'un PFU à 25~\%.} Formule de Harberger appliquée au seul volet réel : la part $\pi$ de la réaction
compte comme perte sèche évitée, au coin moyen ; la recette perdue en régime permanent est levée ailleurs au « prix
d'un milliard » $\mu-1$ = 0,04 / 0,08 / 0,17 (barème fusionné, \texttt{synthese/sorties/architectures\_ircsg.txt}).
Avec $e$ = 0,77, revenus de 2025 : +0,17 / +0,34 / +0,73~Md€ si $\pi = 0$ ; $-$0,50 / $-$0,35 / $-$0,02 si
$\pi = 0{,}5$ ; $-$1,17 / $-$1,05 / $-$0,76 si $\pi = 1$. La baisse allège la perte sèche si la moitié au moins de
la réaction est une évasion définitive.

\emph{Bloc épargne.} Avec une élasticité réelle de l'épargne au rendement net $\varepsilon_s$, la perte sèche sur
l'assiette du PFU vaut $\tfrac12\,\varepsilon_s\,t^2/(1-t)\,B$ ; passer de 31,4 à 25~\% la réduirait de 0,20 / 0,40 /
1,01~Md€ pour $\varepsilon_s$ = 0,1 / 0,2 / 0,5 (revenus de 2025). Aucune de ces valeurs n'est estimée sur données
françaises.

\emph{Retour au barème.} Hypothèses fortes : 96~\% des dividendes au taux marginal de 45~\% et 4~\% à 30~\%, intérêts
à 30~\%, plus-values à 45~\%, abattement de 40~\% sur les dividendes. Revenus de 2025 : 10,1~Md€ en statique, avant
les foyers déjà au barème, la CDHR perdue (0,4 à 0,6) et les foyers à 41~\% ; 5,1 / 4,2 / 2,4~Md€ après réaction
($e$ = 0,65 / 0,77 / 1) ; 2,8 si les dividendes déclarés baissaient de 40~\%. L'Institut Montaigne, cité par le CAE
(2025), donne 3,2 à 3,6~Md€ en statique. Un chiffrage propre demande les fichiers POTE.

\emph{CDHR.} Rentier qui ne vit que de dividendes, à 400\,000~€ comme à 3~M€ : 38,6~\% aujourd'hui, avec un PFU à
25~\% obtenu par la part d'IR et sans la majoration de la CEHR (la CDHR reprend l'écart) ; 32,2~\% si le PFU baisse
par les prélèvements sociaux ; au barème, 38,6 et 45,3~\%. Baisser la part d'IR de 12,8 à 6,4~\% multiplie la CDHR payée par 2,05 et 2,77 : la reprise serait de l'ordre de 0,7 à 1,1~Md€ si ses
redevables étaient des rentiers purs (1,05 à 1,77 fois son rendement de 0,62~Md€), base non publiée.

\emph{Exil.} Avec 0,02 à 0,23~\% du centile aux plus hauts revenus du capital qui part à long terme par point
d'imposition effective (Bach et al., 2025, CAE), un PFU à 25~\% en retiendrait 0,13 à 1,5~\%, soit 0,02 à 0,17~Md€
de recette gardée (assiette du PFU du centile : 47~Md€ en 2025).

\subsection{La capitalisation dans une holding}

100~€ de bénéfice d'une filiale avant IS. A : distribués au PFU. B : remontés en holding, placés $N$ ans à $R$
(revenus imposés à 1,25~\%), puis distribués. C : comme B, puis transmis, les héritiers sortant la valeur sans impôt
(avec la proposition : au PFU, dix ans après). D : distribués, puis placés personnellement, revenus au PFU chaque année.
Actualisation au taux $R$.

\begin{table}[!htbp]\centering\scriptsize
\caption{Impôt total en valeur actualisée selon le chemin du bénéfice (\%)}\label{tab:pfu_a_holding}
\begin{tabularx}{\linewidth}{@{}>{\raggedright\arraybackslash}X*{6}{>{\raggedleft\arraybackslash}p{1.45cm}}@{}}
\toprule
$R$, $N$ & A & B & C & C proposition & D & Richesse B/D \\
\midrule
3~\%, 25 ans & 49,1 & 49,6 & 26,5 & 49,8 & 59,6 & +25~\% \\
3~\%, 40 ans & 49,1 & 49,9 & 26,9 & 50,0 & 64,8 & +42~\% \\
5~\%, 25 ans & 49,1 & 49,9 & 26,9 & 50,2 & 65,1 & +44~\% \\
5~\%, 40 ans & 49,1 & 50,3 & 27,6 & 50,6 & 72,1 & +78~\% \\
7~\%, 25 ans & 49,1 & 50,1 & 27,3 & 50,6 & 69,7 & +65~\% \\
\bottomrule
\end{tabularx}
\par\vspace{3pt}\parbox{\linewidth}{\raggedright\scriptsize Note : si la holding place en titres de portefeuille,
dont les revenus paient l'IS de 25~\%, B vaut 62,3~\% et C 45,0~\% ($R$ = 5~\%, $N$ = 25) : l'avantage de
capitalisation disparaît.}
\end{table}

\subsection{L'effacement à la transmission : chiffrage}

\emph{Méthode.} Plus-values effacées par an : $G = E\,b\,\lambda$ (titres concernés, part transmise chaque année,
part de plus-value latente), plus les reports transmis $S/T$ (stock sur délai moyen avant transmission). Impôt
effacé : $L = 31{,}4~\%\times G + 30~\%\times S/T$ (taux d'expiration du report non tranché). Une cohorte paie une
part $a$ de son impôt latent dans l'année, puis $(1-a)\,h\,(1-h)^{k-2}$ la $k$-ième année ($a$ = 0,1 pour les
reports) ; valeur actualisée : $L\,[a/(1+r) + (1-a)\,h/((1+r)(r+h))]$. Valeur des titres supposée constante après la
transmission (prudent).

\begin{table}[!htbp]\centering\scriptsize
\caption{Effacement à la transmission : hypothèses et résultats}\label{tab:pfu_a_effacement}
\begin{tabularx}{\linewidth}{@{}>{\raggedright\arraybackslash}X*{3}{>{\raggedleft\arraybackslash}p{2.55cm}}@{}}
\toprule
 & Bas & Central & Haut \\
\midrule
Titres concernés $E$ (Md€) & 1\,200 & 1\,450 & 1\,700 \\
Part transmise par an $b$ & 1,7~\% & 2,3~\% & 3,0~\% \\
Part de plus-value latente $\lambda$ & 40~\% & 55~\% & 70~\% \\
Ventes dans l'année $a$ ; ensuite $h$ par an & 15~\% ; 2~\% & 30~\% ; 4~\% & 45~\% ; 6~\% \\
Stock de reports $S$ (Md€) ; délai $T$ (ans) & 127 ; 35 & 144 ; 25 & 173 ; 15 \\
\midrule
Plus-values effacées par an, hors reports (Md€) & 8,2 & 18,3 & 35,7 \\
Reports transmis par an (Md€) & 3,6 & 5,8 & 11,6 \\
Impôt effacé par an (Md€) & 3,65 & 7,49 & 14,68 \\
\quad en \% du PIB & 0,12 & 0,25 & 0,49 \\
Recette, an 1 / 5 / 10 (Md€) & 0,49 / 0,74 / 1,02 & 1,90 / 2,74 / 3,62 & 5,39 / 7,43 / 9,36 \\
Recette, an 20 / 30 (Md€) & 1,50 / 1,89 & 4,91 / 5,78 & 11,81 / 13,13 \\
Valeur actualisée d'une année de transmissions, $r$ = 2,5 / 3,5 / 5~\% & 1,85 / 1,59 / 1,33 & 5,21 / 4,71 / 4,17
& 11,65 / 10,87 / 9,96 \\
\quad avec intérêt sur l'impôt différé & 3,65 & 7,49 & 14,68 \\
Clause de droits acquis : recette an 10 ; cumul ans 1-10 & 0,47 ; 2,4 & 1,43 ; 7,5 & 3,04 ; 16,4 \\
Sans clause : cumul ans 1-10 & 7,6 & 28,1 & 75,4 \\
Charge du 1~\% le plus aisé : moyenne ans 1-10 ; an 10 & 0,69 ; 0,92 & 2,53 ; 3,26 & 6,78 ; 8,42 \\
\bottomrule
\end{tabularx}
\par\vspace{3pt}\parbox{\linewidth}{\raggedright\scriptsize Note : jeux bas et haut = toutes les hypothèses basses
ou toutes les hautes (des bornes). $E$ part des 1\,969~Md€ d'actions et de fonds en actions des ménages fin 2025
(Banque de France), moins les parts de sociétés immobilières, le PEA et l'épargne salariale. Repère de $b$ : 300~Md€
de transmissions par an (CAE, 2021) pour 15\,300~Md€ de patrimoine net (Insee). Droits acquis : valeurs en hausse de
4~\% par an, reports existants exclus.}
\end{table}

\emph{Repères.} Le Trésor américain évalue l'effacement au décès à 37,8~Md\$ en 2025 (\emph{Tax Expenditures}
FY2026), soit 5,0~Md€ transposés au PIB et au taux français, sans les donations ; la Tax Foundation (Li, 2020) chiffre
sa suppression à 11 à 12~Md\$ par an, réactions comprises, soit 2,0~Md€. Le premier repère encadre l'impôt effacé, le
second la recette des premières années.

\figbloc{pfu_n10_sensibilite}{La recette de la dixième année dépend d'abord de la part transmise, de la plus-value latente et du rythme des ventes}{fig:pfu_sensibilite}

\emph{Variantes.} Une franchise de 100\,000~€ par transmission coûterait 0,4 à 1,1~Md€ par an. Déduire l'impôt latent
de l'assiette des droits de succession, exactement neutre entre vendre avant de mourir et transmettre, coûterait 0,7
à 3,0~Md€ par an selon le taux de droits effectivement payé, pacte Dutreil compris.

\subsection{La charge du 1~\% le plus aisé}

La fin de l'effacement coûterait au 1~\% 3,3~Md€ la dixième année, et 0,7~Md€ par an en moyenne avec la clause de
droits acquis. Le bilan d'ensemble (section~\ref{sec:ens_haut})
compte aussi la suppression de la majoration ex-CEHR sur le capital (0,7 à 1,1~Md€ rendus au haut) ; selon la façon
dont on répartit les revenus au PFU, la fusion peut alléger de 2,8~Md€ la charge du 1~\% ; et la fin de la
contribution exceptionnelle sur l'IS rend 5 à 7~Md€ aux quelque 450 entreprises qui la paient, dont une part
inconnue revient aux mêmes ménages.

\subsection{Le stock de reports d'apport-cession}

Valeur actualisée au taux $r$ de l'impôt sur le stock, avec trois taux de sortie annuels : expiration du vivant $e$
(environ 0,9~\% du stock en 2024), transmission $m = 1/T$, vente par les héritiers $h$. Droit actuel :
$I\,e/(r+e+m)$ ; expiration à la transmission (PLF 2027, art.~5), ou report transmis avec intérêt au taux $r$ :
$I\,(e+m)/(r+e+m)$ ; report transmis : $I\,[e+m\,h/(r+h)]/(r+e+m)$.

\begin{table}[!htbp]\centering\scriptsize
\caption{Impôt sur le stock de reports d'apport-cession, en valeur actualisée ($r$ = 3,5~\%, Md€)}\label{tab:pfu_a_reports}
\begin{tabularx}{\linewidth}{@{}>{\raggedright\arraybackslash}X*{3}{>{\raggedleft\arraybackslash}p{2.1cm}}@{}}
\toprule
 & Bas & Central & Haut \\
\midrule
Impôt nominal (30~\% du stock) & 38,1 & 43,2 & 52,0 \\
Droit actuel (effacement à la transmission) & 4,6 & 4,5 & 4,1 \\
Expiration à toute transmission (PLF 2027, art.~5) & 19,7 & 25,1 & 35,5 \\
Report transmis sans effacement (proposition) & 10,1 & 15,5 & 23,9 \\
Report transmis, avec intérêt sur l'impôt différé & 19,7 & 25,1 & 35,5 \\
Recette annuelle de l'art.~5 en régime de croisière (Md€) & 1,09 & 1,73 & 3,47 \\
\bottomrule
\end{tabularx}
\par\vspace{3pt}\parbox{\linewidth}{\raggedright\scriptsize Note : l'évaluation préalable de l'article~5 (0,5~Md€ en
2027, « en tenant compte de l'étalement ») est compatible avec une année pleine de 2 à 2,5~Md€. Borne de très long
terme : 24,7~Md€ d'entrées par an (2022-2024) $\times$ 30~\% = 7,4~Md€.}
\end{table}

\emph{L'entrepreneur qui cède} (7,5~M€ de plus-value, $h$ = 4~\% ; figure~\ref{fig:pfu_entrepreneur}) : impôt
actualisé à 3,5~\%, article~5 / report transmis, pour un décès 15, 20 ou 30 ans après : 1,41 / 0,75, 1,18 / 0,63,
0,84 / 0,45~M€ ; de 1,63 / 1,00 à 0,54 / 0,24~M€ avec un taux de 2,5 ou 5~\%. Faire expirer le report quand la
holding vend ses actifs ou devient une simple société de placement serait plus ciblé, mais difficile à contrôler et
non chiffrable. Les reports déjà transmis par donation, dont le délai de
conservation court encore, garderaient l'ancienne règle.

\figbloc{pfu_n09_entrepreneur}{Pour un entrepreneur qui cède, le report transmis vaut la moitié de l'impôt dû au décès, l'intérêt sur l'impôt différé le rétablit}{fig:pfu_entrepreneur}

\subsection{Expériences étrangères}

D'autres pays reportent la plus-value sur l'héritier (la moitié des pays de l'OCDE selon une étude citée par le
rapport n°~3056, non lue). L'évaluation des titres à la date de la réforme a compté parmi les difficultés américaines
de 1976 (non vérifié).

\begin{table}[!htbp]\centering\scriptsize
\caption{Le sort de la plus-value latente à la donation et au décès}\label{tab:pfu_a_etranger}
\begin{tabularx}{\linewidth}{@{}>{\raggedright\arraybackslash}p{1.8cm}>{\raggedright\arraybackslash}X>{\raggedright\arraybackslash}p{4.2cm}@{}}
\toprule
Pays & Règle & Source, statut \\
\midrule
France & Prix d'acquisition = valeur retenue pour les droits de mutation (CGI, art.~150-0~D) ; report effacé au
décès, et après donation si le donataire conserve six ou onze ans & Rapport n°~3056 ; évaluation préalable du PLF
2027 (lus) \\
Canada & Depuis 1972 (et la fin de l'impôt fédéral sur les successions), vente réputée faite juste avant le décès ;
conjoint au prix d'acquisition ; plus-values antérieures exonérées (« jour V ») & Agence du revenu du Canada ;
Vaillancourt et Kerkhoff (2019) (lus) ; jour V lu par résumé \\
États-Unis & Donataire : prix du donateur (IRC §~1015) ; héritier : valeur au décès (§~1014). Reprise du prix au décès
votée en 1976, abrogée en 1980. Proposition de 2015 : imposition à la donation et au décès, franchise de
100\,000~\$, report pour les entreprises familiales, paiement sur quinze ans &
\emph{U.S. Code} §~1015 ; Trésor, \emph{General Explanations} FY2016, p.~156-157 (lus) \\
Allemagne & Participations d'au moins 1~\% reçues à titre gratuit : prix du dernier acquéreur à titre onéreux (EStG,
§~17) & Loi (lue) ; autres titres non vérifiés \\
Australie & Biens acquis depuis le 20 septembre 1985 : l'héritier reprend le prix du défunt & Administration fiscale,
lue par résumé \\
\bottomrule
\end{tabularx}
\end{table}

\subsection{Hypothèses fortes et chiffres non vérifiés}

\begin{table}[!htbp]\centering\scriptsize
\caption{Hypothèses fortes (registre \texttt{revue/hypotheses\_pfu.md})}\label{tab:pfu_a_hyp}
\begin{tabularx}{\linewidth}{@{}>{\raggedright\arraybackslash}p{4.4cm}>{\raggedright\arraybackslash}p{3.3cm}>{\raggedright\arraybackslash}X@{}}
\toprule
Hypothèse & Où elle joue & Si elle est fausse \\
\midrule
Élasticité des revenus déclarés de 0,77 (0,65 à 1) & Changement de taux & À 0,5, un PFU à 25~\% coûte 3,5~Md€ les
premières années ; à 1,2, 2,5. Mesurée sur des foyers qui ont aussi des revenus d'activité \\
Réaction des dividendes surtout temporelle ($\pi$ faible) & Coût permanent d'une baisse ; perte sèche & Si
$\pi = 1$, coût permanent de 3,1~Md€ et perte sèche réduite par la baisse ; si $\pi = 0$, une hausse rapporte
durablement presque son rendement mécanique \\
Épargne insensible au rendement & « Une baisse n'achète rien » & Élasticité de 0,2 à 0,5 : la baisse gagne 0,4 à
1,0~Md€ de perte sèche, autant ou plus que son financement \\
Assiette de 2025 estimée (67~Md€) et représentative & Tous les chiffrages de taux & Celle de 2024, gonflée par la
CDHR, donne des chiffres plus élevés d'un cinquième ; celle de 2025 est peut-être déprimée par la CDHR elle-même \\
Réaction au taux attendu & Stabilité du taux & Documentée en 2024-2025 ; « suggérée » pour 2013 \\
Le bénéfice gardé en holding finit par sortir & « Report, pas exonération » & Sinon, 27~\% pour cette part ; aucune
donnée \\
$E$ = 1\,450~Md€ (1\,200 à 1\,700) ; $b$ = 2,3~\% (1,7 à 3) ; $\lambda$ = 55~\% (40 à 70) & Recette de la fin de
l'effacement & Proportionnelle à chacun ; $b$ et $\lambda$ ensemble : recette de l'an 10 de 2,3 à 5,6~Md€. Pas de
statistique sur la composition des successions depuis 2006 \\
Ventes des héritiers : $a$ = 30~\%, puis $h$ = 4~\% & Profil de la recette & Les héritiers garderont plus leurs
titres : le jeu bas ramène l'an 10 de 3,6 à 2,1~Md€ ; les donateurs n'attendront plus le décès pour vendre (non
chiffré) \\
90~\% de la recette payée par le 1~\% & Charge du 1~\% & À 75~\%, 2,1~Md€ au lieu de 2,5. Le dernier centile paie
72~\% de l'impôt hors barème de 2025 (78~\% en 2024) \\
Impôt personnel du 1~\% sur le capital : 16~Md€ (21 en 2024, année gonflée par la CDHR) & « 20~\% » & Assiette du
centile déduite de son impôt hors barème (47~Md€ en 2025, 63 en 2024), plus 65~\% de la CEHR et la CDHR \\
Stock de reports de 144~Md€ (127 à 173), transmis en 25 ans (15 à 35), imposé à 30~\% & Report & Taux non
tranché ; abattements sur les entrées de 2013-2017 \\
CEHR payée à 50 à 80~\% sur le capital & Majoration (0,7 à 1,1~Md€ ; 1,0 à 1,6 en 2024) & Inconnue \\
Constitutionnalité (plus-values acquises avant la réforme ; effet dès l'annonce) & Transition & Non expertisée \\
\bottomrule
\end{tabularx}
\end{table}

\begin{table}[!htbp]\centering\scriptsize
\caption{Chiffres non vérifiés sur une source primaire}\label{tab:pfu_a_nonverif}
\begin{tabularx}{\linewidth}{@{}>{\raggedright\arraybackslash}p{4.4cm}>{\raggedright\arraybackslash}p{3.3cm}>{\raggedright\arraybackslash}X@{}}
\toprule
Chiffre ou fait & Où il joue & Statut \\
\midrule
Revenus de 2025 (impôt hors barème, CEHR, CDHR, revenus déclarés) & Situation & DGFiP n°~53, provisoire ; assiette
par nature de revenu estimée par l'auteur \\
Taux total de 46~\% à l'entrée du top 0,1~\%, 26~\% pour le top 0,0002~\% (2016) & Diagnostic & IPP, note n°~92,
non lue (via le rapport n°~3056) \\
84 foyers, 23,5~Md€ mis en report en 2023 & Diagnostic & Via le rapport n°~3056 ; incohérent avec les 18,8~Md€
du top 0,01~\% \\
\emph{Exit tax} dégrevée après deux ou cinq ans, quinze avant 2019 & Transition & Sources secondaires \\
Droits de succession sur les entreprises différés avec intérêt & Intérêt sur l'impôt différé & Sources
secondaires ; taux non vérifié \\
Rôle de l'évaluation au 31 décembre 1976 dans l'abandon américain & Clause de droits acquis & Non vérifié \\
Régime allemand hors participations de 1~\% ; règles australiennes ; PIB américain & Exemples ; repères &
Non vérifiés sur les textes \\
CSG déductible de 6,8~\% ; barème des revenus de 2024 ; IS de 15~\% jusqu'à 42\,500~€ & Cas types & Non
revérifiés pour 2026 ; effets de second ordre \\
\bottomrule
\end{tabularx}
\end{table}
}{}
\section{Méthode commune}\label{ann:methode}
\IfFileExists{parties/annexe_methode.tex}{% Annexe de méthode commune, écrite le 16/11/2026 ; base 2025 le 17/11/2026 ; coupée le 17/11/2026
% (synthese/CONSIGNES_COUPE.md) : reçoit le bilan détaillé du 1 % (ann:meth_haut, tab:ens_haut), sorti de la synthèse.
% Sources : revue/methode_perte_seche.md, revue/methode_v2.md, synthese/perte_seche_autres.py,
% synthese/synthese_ensemble.py. Les calibrages propres à chaque réforme sont dans les annexes de partie (ann:ircsg,
% ann:foncier, ann:is, ann:pfu) et ne sont pas repris ici.

\subsection{Une seule mesure de l'efficacité : la perte sèche généralisée}\label{ann:meth_ps}

La note mesure de combien chaque réforme réduit la perte sèche par rapport à aujourd'hui, à recette égale, avec la
même formule pour les quatre (Harberger, 1964 ; Hines, 1999) :
\[
\Delta PS \;=\; -\sum_j \big(\bar T_j - s_j\,\tau_{k(j)}\big)\,\Delta Q_j ,\qquad
\Delta Q_j = \varepsilon_j\,Q_j\,\Delta\ln(1-t_j).
\]
L'indice $j$ parcourt les marges sur lesquelles la réforme agit ; $Q_j$ est la quantité taxée (revenu d'activité,
ventes de logements, stock d'équipements, revenus du capital déclarés) ; $t_j$, le taux réformé auquel elle réagit ;
$\bar T_j$, le coin total sur la marge, moyenne de ses valeurs avant et après ; $s_j$, la part de la réaction qui se
déplace vers une autre assiette taxée au taux $\tau_k$ (Chetty, 2009). Sur une marge déjà taxée, un point de plus coûte
au premier ordre le coin existant multiplié par la quantité perdue : beaucoup là où le coin est lourd, comme au
voisinage du Smic.

\paragraph{Le coin total} compte tout ce qu'un euro d'activité perdu fait perdre à la collectivité : pour l'IR-CSG,
l'impôt fusionné, les cotisations qui n'ouvrent pas de droits et le retrait de la prime d'activité (\ref{ann:ircsg}) ;
pour l'immobilier, le droit de mutation sur les ventes et l'impôt sur le bâti rapporté au coût d'usage sur la
construction, le sol n'ayant pas de marge ; pour l'IS, le coin sur le coût d'usage des équipements (Hall-Jorgenson).

\paragraph{Le réel et le déclaratif.} Seule la baisse réelle de l'activité ou de l'investissement entre dans le gain.
Un revenu déclaré plus tard, autrement ou dans une société est traité comme un effet sur la recette, net de ce qui est
récupéré ailleurs (la moitié de la réaction des hauts revenus passe à l'IS et au prélèvement forfaitaire, à 48,6~\%,
convention), donné à part et sans rien financer, ce qui minore le gain si l'optimisation consomme des ressources
réelles. Une réaction de calendrier, comme celle des dividendes en 2013, 2018 et 2024, ne compte presque pas.

\paragraph{Les marges.} Les heures réagissent au taux marginal ; la décision de travailler, sous deux Smic, au taux de
prélèvement sur tout le premier salaire. Ni les pensions, ni les allocations chômage, ni l'épargne ne réagissent, faute
de mesure ; l'âge de départ et le cumul emploi-retraite ne sont pas modélisés.

\subsection{Quatre jeux d'hypothèses sur les comportements}\label{ann:meth_jeux}

Les élasticités ont été mesurées face à de petites réformes, d'une année sur l'autre ; nul ne sait comment les
comportements réagiraient à un système entièrement différent. D'où trois jeux, bas, central (prudent, calé sur les
études françaises autant que possible) et haut, plus pour la fusion un jeu de long terme (tableaux~\ref{tab:ircsg_jeux}
et~\ref{tab:meth_jeux}).

\begin{table}[!htbp]\centering\scriptsize\setlength{\tabcolsep}{4pt}
\caption{Les paramètres de comportement de l'immobilier, de l'IS et des revenus du capital}\label{tab:meth_jeux}
\begin{tabularx}{\linewidth}{@{}>{\raggedright\arraybackslash}X>{\raggedright\arraybackslash}p{1.5cm}>{\raggedright\arraybackslash}p{1.5cm}>{\raggedright\arraybackslash}p{1.9cm}>{\raggedright\arraybackslash}p{3.4cm}@{}}
\toprule
Paramètre & Bas & Central & Haut & Appui \\
\midrule
Immobilier : ventes en moins par point de droits & 0 & 6~\% & 15~\% & convention, d'après Bérard et Trannoy (2018), Eerola et al. (2021) \\
Immobilier : capital logement, élasticité au coût d'usage & 0,5 & 1 & 1 (coût d'usage de 3,5~\%) & convention \\
IS : stock d'équipements, élasticité au coût d'usage & 0,5 & 1 & 1,6 & partie IS (H18-H19) ; borne basse sans source \\
Revenus du capital : revenus déclarés, élasticité au taux net & 0,65 & 0,77 & 1 & Lefebvre, Lehmann et Sicsic (2025) ; CAE (2025) \\
Revenus du capital : part de la réaction qui échappe définitivement à l'impôt & 0 & 0 & 1 & non mesurée \\
Épargne : élasticité au rendement net & 0 & 0 & 0 (0,1 à 0,5 en variante) & aucune estimation française \\
\bottomrule
\end{tabularx}
\par\vspace{2pt}{\raggedright\scriptsize Des choix appuyés sur les études, non des estimations. Les jeux bas et haut
combinent toutes les hypothèses basses ou hautes : ce sont des bornes, pas des scénarios.\par}
\end{table}

\paragraph{Le prix d'un milliard.} L'écart de recette entre deux variantes est supposé levé par un point de taux de
l'impôt fusionné pour tous, au prix de 0,04~Md€ de perte sèche par milliard dans le jeu bas, 0,08 dans le central et
0,17 dans le haut. C'est peu parce que les élasticités sont basses : avec 0,3 et un coin de 40 à 50~\%, la formule
usuelle donnerait 0,25 à 0,43 (avis Cochrane).

\subsection{Conventions communes}\label{ann:meth_conv}

\paragraph{Référence.} Le droit en vigueur en novembre 2026, sans le projet de loi de finances pour 2027, donc avec la
fin de la contribution exceptionnelle sur l'IS et l'extinction de la CRDS, que la synthèse affiche à part. Chaque
réforme est calibrée à recette constante avant toute réaction.

\paragraph{Base 2025.} Les chiffrages sont calés sur les prélèvements obligatoires de 2025, dans le sens que précise
l'annexe de chaque partie (pour l'IS, bénéfices de 2025 portés en 2027 ; pour la CEHR, revenus de 2025). Restent en
2024 la répartition entre foyers, le détail des déclarations, l'impôt sur le revenu imputé aux revenus du capital et
les comparaisons de l'OCDE. La synthèse additionne ces montants sans les indexer : l'erreur, de quelques pour cent, est
petite devant les fourchettes.

\paragraph{Actualisation.} 3,25~\% pour l'État (dette de transition, déduction immédiate), 5~\% pour les entreprises,
3~\% réels pour les ménages (taxe foncière capitalisée, report de paiement, droits annualisés), 3,5~\% nominal (2,5 à
5~\%) pour les héritiers, 1 à 3~\% réels pour la valeur actualisée de la hausse des retraités.

\paragraph{Incidence.} L'impôt sur le revenu et la CSG pèsent sur ceux qui les paient ; la hausse de taxe foncière
passe dans les prix et la baisse des droits profite à l'acheteur (cas central) ; l'IS est compté au niveau de la
société.

\subsection{Les calculs de la synthèse}\label{ann:meth_synthese}

\paragraph{Trajectoire.} Soldes des parties : immobilier sur cinq ans (droits divisés par deux dès l'annonce, taux
national montant d'un cinquième par an), IS au cas central (comptes de l'exercice, en caisse un an plus tard),
effacement au jeu central ; un an de retard sur 2027 décalerait la dette sans la changer. Le cas défavorable combine
90~\% d'investissement fait par des sociétés imposables et le jeu bas de l'effacement, le favorable 60~\%, le jeu haut
et la reprise des ventes.

\paragraph{Gain total.} Celui de l'IS est la moyenne des deux régimes d'amortissement actuels à une élasticité de 1,
borné par le dégressif à 0,5 et le linéaire à 1,6. Les bornes de l'immobilier sont les jeux bas et haut du
tableau~\ref{tab:meth_jeux}. Le gain permanent retire de celui de la fusion la part des retraités actuels, selon que
les actifs perçoivent à moitié ou entièrement l'impôt futur sur leur pension (\ref{ann:ircsg}).

\paragraph{Cas-types.} Impôt fusionné des cas-types OpenFisca de la partie IR-CSG ; logement médian de la commune
(Limoges, 171\,800~€ ; Paris, 416\,200~€ ; Lyon, 279\,200~€), sans effet de la taxe sur les prix ni sur les loyers ;
droits annualisés sur dix ans à 3~\% ; PME industrielle de la partie IS ; impôt transmis à 31,4~\% de la plus-value
latente supposée (80\,000~€ pour le couple de retraités, 120\,000~€ pour le cadre, 5~M€ pour le dirigeant, 40~M€ pour
le rentier).

\begin{table}[!htbp]\centering\scriptsize
\caption{Hypothèses propres à la synthèse}\label{tab:meth_hyp}
\begin{tabularx}{\linewidth}{@{}>{\raggedright\arraybackslash}p{5.2cm}>{\raggedright\arraybackslash}p{2.2cm}>{\raggedright\arraybackslash}X@{}}
\toprule
Hypothèse & Où elle joue & Si elle est fausse \\
\midrule
Le 1~\% paie 90~\% de la fin de l'effacement & bilan du haut & À 75~\%, 2,1~Md€ au lieu de 2,5 \\
Le 1~\% paie 0 à 0,5~Md€ du taux national sur le sol des entreprises et les terrains nus & bilan du haut & Non vérifiable :
détention non publiée par niveau de revenu \\
50 à 80~\% de la CEHR portent sur les revenus du capital & bilan du haut & Ne fait que répartir un coût de 1,4~Md€ \\
IR imputé au capital (revenus de 2024) : 10,2~Md€ au lieu de 11,5 en variante & bilan du haut & Majorant ; l'impôt de
l'assurance-vie, non chiffré, réduirait un peu la correction \\
Report de taxe foncière réservé aux 65 ans et plus : moitié des hausses & dette brute & Ordre de grandeur non vérifié ;
majorant de 20~Md€ fin 2036 \\
Taxe sur le sol sans effet sur les loyers & cas-types & Effet théorique, non mesuré \\
\bottomrule
\end{tabularx}
\end{table}

\subsection{Le bilan du 1~\% le plus aisé}\label{ann:meth_haut}

Le tableau~\ref{tab:ens_haut} fait le compte pour le centile des foyers au revenu fiscal de référence le plus élevé.
La fusion supprime la CEHR sur les revenus du barème, ce que la partie IR-CSG ne compte pas. Restent non répartis le
taux national sur le sol des entreprises et des terrains nus (1,4~Md€ net) et la fin du taux réduit de l'IS, qui touche
des holdings patrimoniales et des sociétés de professions libérales.

\begin{table}[!htbp]\centering\scriptsize
\caption{La charge du 1~\% le plus aisé (Md€ par an ; + : paie davantage)}\label{tab:ens_haut}
\begin{tabularx}{\linewidth}{@{}>{\raggedright\arraybackslash}X>{\raggedright\arraybackslash}p{3.2cm}>{\raggedright\arraybackslash}p{2.6cm}>{\raggedright\arraybackslash}p{2.6cm}@{}}
\toprule
 & Convention centrale (fourchette) & IR imputé au capital corrigé & Revenus du capital logés tout en haut \\
\midrule
Fusion IR-CSG & 0,0 & $-$1,1 & $-$3,0 \\
CEHR sur les revenus du barème & $-$0,5 ($-$0,3 à $-$0,7) & $-$0,5 & $-$0,5 \\
Fin de l'effacement (moyenne sur dix ans) & +2,5 (0,7 à 6,8) & +2,5 & +2,5 \\
Immobilier (sol des entreprises et terrains nus) & 0 à +0,5, non compté & id. & id. \\
Fin du taux réduit de l'IS & non chiffré & id. & id. \\
\textbf{Bilan} & \textbf{+2,0 ($-$0,1 à 7,0)} & \textbf{+0,9} & \textbf{$-$1,0} \\
Suppression de la majoration ex-CEHR sur le capital & $-$0,9 & $-$0,9 & $-$0,9 \\
\textbf{Après suppression} & \textbf{+1,1 ($-$1,0 à 6,1)} & \textbf{0,0} & \textbf{$-$1,9} \\
Barème adouci au-delà de 150\,000~€ & $-$1,1, soit 0,0 & --- & --- \\
\bottomrule
\end{tabularx}
\par\vspace{2pt}{\raggedright\scriptsize Hypothèses : tableau~\ref{tab:meth_hyp}. Source :
\texttt{sorties/synthese\_ensemble.txt}, \texttt{controle\_hauts\_purs.txt}.\par}
\end{table}

Le bilan dépend surtout de la répartition par tranche de revenu des revenus au prélèvement forfaitaire, que la DGFiP ne
publie pas. Le modèle IR-CSG les répartit au prorata des autres revenus, les loge d'abord tout en haut (« haut », que
suggère un calage sur l'impôt observé), ou prend la moyenne des deux (« mixte », centrale). Second écart, de même sens :
sur les revenus de 2024, le modèle retire de l'impôt sur le revenu 11,5~Md€ imputés aux revenus du capital, la partie
PFU au plus 10,2, sur l'assiette du prélèvement forfaitaire (79,6~Md€), qui exclut l'assurance-vie et les plus-values
immobilières. Corrigé, l'impôt des revenus du barème est sous-estimé d'environ 1,4~Md€ en 2025 : la fusion rendrait
cette somme, dont 1,1~Md€ au 1~\%, et la recette à remplacer passerait à environ 231,3~Md€ ; tenir inchangée la charge du 1~\% à recette constante
obligerait alors à atteindre 50~\% dès 180\,000~€ et ferait bouger de 540~€ la charge moyenne d'un dixième. Le modèle porte en outre cet
impôt au niveau de 2025 (+7,9~\%), quand la partie PFU le mesure en baisse de 18~\% sur les revenus de 2025 : ces
conventions ne sont pas réconciliées.

La suppression de la majoration ex-CEHR coûte 0,9~Md€ (0,7 à 1,1), un peu moins tant que la contribution
différentielle sur les hauts revenus reprend la baisse des rentiers à son plancher : le rentier des cas-types n'y gagne
rien avant qu'elle disparaisse, 107\,500~€ par an ensuite. Calculée sur la CEHR des revenus de 2024, encaissée en
2025, elle coûterait 1,9~Md€, le 1~\% paierait 0,5~Md€ de plus et la dette de transition culminerait à 35~Md€.

Si le calage confirme la convention centrale, rendre encore environ 1,1~Md€ au 1~\% passerait au mieux par le haut du
barème : le garder jusqu'à 150\,000~€ de revenu imposé (41,5~\%), puis n'atteindre 50~\% qu'à 1,25~M€ au lieu de
231\,000~€. Seul le 1~\% est touché, aucun dixième ne bouge de plus de 360~€, le gain de la fusion monte de 4,7 à
4,9~Md€ et la dette de transition à 39~Md€. Le faire à recette constante ferait payer les autres foyers, dont les
retraités ($+$0,4~Md€) : nous ne le recommandons pas. Si le calage loge les revenus du capital tout en haut, il faut
garder la majoration.

\paragraph{Calcul.} Le modèle IR-CSG est relancé sans modification dans chaque convention, et une fois avec 10,2~Md€
d'impôt imputé, réduit dans la même proportion dans chaque tranche. Le barème adouci ajoute au barème proposé un nœud à
150\,000~€, d'où le taux monte en ligne droite jusqu'à 50~\% à un seuil cherché par dichotomie.

\subsection{Programmes}\label{ann:meth_programmes}

{\small\raggedright Depuis \texttt{fiscalite/} : \texttt{synthese/synthese\_ensemble.py} (sorties
\texttt{synthese/sorties/synthese\_ensemble.txt} et \texttt{.json}), puis \texttt{synthese/figures\_ensemble.py}, une
seconde fois avec \texttt{FIG\_NU=1} ; perte sèche de l'IS : \texttt{synthese/perte\_seche\_autres.py} ; méthode :
\texttt{revue/methode\_v2.md}. Les programmes des parties sont dans leurs annexes.\par}
}{}

\vspace{8pt}
{\scriptsize\setlength{\parskip}{2pt}\textbf{Références.}
André et Meslin (2023), \emph{Insee Analyses} n°~91.
André et Sireyjol (2021), « Effets redistributifs de l'imposition des couples et des familles : une étude par
microsimulation de l'impôt sur le revenu », \emph{Économie et Statistique} n°~526-527.
Assemblée nationale (2026), commission d'enquête sur l'imposition des plus hauts patrimoines et des revenus les
plus élevés, rapport n°~3056.
Avenancio-León et Howard (2022), « The Assessment Gap », \emph{QJE} 137(3).
Bach, Bozio, Fabre, Guillouzouic, Leroy et Malgouyres (2019), note IPP n°~46.
Bach, Bozio, Grimprel, Guillouzouic, Landais et Malgouyres (2025), Focus du CAE n°~118.
Bachelet et Poulhès (2017), CGDD.
Banzhaf et Lavery (2010), \emph{Journal of Urban Economics} 67(2).
Beer, de Mooij et Liu (2020), \emph{Journal of Economic Surveys} 34(3).
Bérard et Trannoy (2018), \emph{Économie et Statistique} n°~500-501-502.
Bernini et Treibich (2016), \emph{Industrial and Corporate Change} 25(5).
Berry (2021), « Reassessing the Property Tax », Harris School.
Best et Kleven (2018), \emph{Review of Economic Studies} 85(1).
Bontemps, Cherbonnier et Magnac (2023), TSE WP 23-1486.
Cabannes, Houdré et Landais (2014), \emph{Économie et Statistique} 467-468.
CAE (2025, 2026), Focus n°~124 et n°~147.
Carbonnier (2021), « Imposition jointe des revenus et emploi des femmes mariées : estimation à partir du cas
français », \emph{Revue économique} 72(2).
Chetty (2009), « Is the Taxable Income Elasticity Sufficient to Calculate Deadweight Loss? », \emph{AEJ:
Economic Policy} 1(2).
Chetty, Friedman, Olsen et Pistaferri (2011), \emph{QJE} 126(2).
Chetty, Guren, Manoli et Weber (2013), « Does Indivisible Labor Explain the Difference between Micro and Macro
Elasticities? », \emph{NBER Macroeconomics Annual} 27.
Chodorow-Reich, Smith, Zidar et Zwick (2026), NBER WP 32180 (2024), version de janvier 2026.
Conseil constitutionnel, décisions 2000-437 DC, 2012-662 DC et 2019-796 DC.
Conseil d'État, 13 mars 2025, n°~481538.
Conseil des prélèvements obligatoires (2017, 2023, 2025), « Adapter l'impôt sur les sociétés à une économie
ouverte » ; « Les différences d'imposition sur les bénéfices entre PME et grandes entreprises » ; « Tracer un
cadre fiscal et social pluriannuel pour l'industrie française ».
Dachis, Duranton et Turner (2012), \emph{Journal of Economic Geography} 12(2).
Dolls, Fuest, Krolage et Neumeier (2025), \emph{Journal of Urban Economics} 145.
Edmark, Liang, Mörk et Selin (2016), « The Swedish Earned Income Tax Credit: Did It Increase Employment? »,
\emph{FinanzArchiv} 72(4).
Eerola, Harjunen, Lyytikäinen et Saarimaa (2021), \emph{Journal of Urban Economics} 124.
Feldstein (1999), \emph{REStat} 81(4).
Ferreira (2010), \emph{Journal of Public Economics} 94(9-10).
France Stratégie (2023), comité d'évaluation des réformes de la fiscalité du capital, quatrième rapport.
Fritzsche et Vandrei (2019), \emph{Regional Science and Urban Economics} 74.
Harberger (1964), « Taxation, Resource Allocation, and Welfare », dans NBER et Brookings Institution, \emph{The
Role of Direct and Indirect Taxes in the Federal Revenue System}, Princeton University Press.
Hines (1999), « Three Sides of Harberger Triangles », \emph{Journal of Economic Perspectives} 13(2).
HMRC (2023), évaluation de l'allègement pour primo-accédants.
Høj, Jørgensen et Schou (2018), \emph{Fiscal Studies} 39(2).
Hourani, Millar-Powell, Perret et Ramm (2023), OCDE, \emph{Taxation Working Paper} n°~65.
IFS (2023), \emph{Green Budget}, ch.~10.
IPP (2023), rapport n°~47.
Kleven et Kreiner (2006), « The Marginal Cost of Public Funds: Hours of Work versus Labor Force
Participation », \emph{Journal of Public Economics} 90(10-11).
Koster et Pinchbeck (2022), \emph{AEJ: Economic Policy} 14(1).
Lefebvre, Lehmann et Sicsic (2025), \emph{Scandinavian Journal of Economics} 127(2) (IZA DP 16112, 2023).
Lehmann, Marical et Rioux (2013), « Labor Income Responds Differently to Income-Tax and Payroll-Tax Reforms »,
\emph{Journal of Public Economics} 99.
Lutz (2015), \emph{AEJ: Economic Policy} 7(1).
Määttänen et Terviö (2022), \emph{Economic Journal} 132(644).
Maffini, Xing et Devereux (2019), \emph{AEJ: Economic Policy} 11(3).
Martínez, Saez et Siegenthaler (2021), \emph{AER} 111(2).
Morin (2014), \emph{Insee Première} 1492.
Nielsson, Wroblewski et Yding (2024), document de travail.
OCDE, \emph{Revenue Statistics}, \emph{Corporate Tax Statistics} (2025), \emph{Taxing Wages} (2025).
Paquier et Sicsic (2022), \emph{Économie et Statistique} n°~530-31.
Sénat (1977), M.~Blin, rapport général sur le PLF pour 1978, n°~88 ; Sénat (2026), mission d'information sur le
poids des prélèvements obligatoires en France et ses conséquences sur la compétitivité des entreprises, rapport
n°~942 (2025-2026).
Sicsic (2022), « Does Labour Income React More to Income Tax or Means-Tested Benefits Reforms? », \emph{Fiscal
Studies} 43(3).
Tax Foundation, \emph{International Tax Competitiveness Index} 2025 et 2026 ; \emph{Dividend Tax Rates in
Europe} 2026.
Yagan (2015), \emph{AER} 105(12).
Zwick et Mahon (2017), \emph{AER} 107(1).
Textes : LF 2024 (art.~19 et 33), LF 2025, LFSS et LF 2026, PLF 2027 (art.~5).
Détail des sources et statut de lecture : \texttt{ir\_csg/}, \texttt{foncier/notes/}, \texttt{is/notes/},
\texttt{pfu/notes/}, \texttt{revue/}.\par}


\end{document}
